\documentclass[11pt]{article}
\usepackage{ifthen}
\usepackage[utf8]{inputenc}
\usepackage[T1]{fontenc}
\usepackage[brazil]{babel}
\usepackage{geometry}
\geometry{a4paper,margin=2cm}
\usepackage{longtable,graphicx,multirow,enumitem,tabularx,setspace,ragged2e}
\setlist{nosep,leftmargin=*}
\renewcommand\arraystretch{1.15}
% Caixa de seleção
\newcommand{\chk}[1]{\ifx#1X\setlength\fboxsep{1pt}\fbox{\rule{1.3ex}{0pt}X}%
\else\setlength\fboxsep{1pt}\fbox{\rule{1.3ex}{0pt}\rule{1.3ex}{0pt}}\fi}
\pagestyle{empty}
\begin{document}
\begin{center}
\begin{tabular}{lccr}
 \multirow{3}{*}{\includegraphics[height=2.7cm]{brasao.png}} &
 \multicolumn{2}{c}{\bfseries UNIVERSIDADE FEDERAL DE OURO PRETO} &
 \ \ \ \ \multirow{3}{*}{\includegraphics[height=2.7cm]{ufop.png}} \\
 & \multicolumn{2}{c}{\bfseries PRÓ-REITORIA DE GRADUAÇÃO} & \\
 & \multicolumn{2}{c}{} & \\
 & \multicolumn{2}{c}{\Large\bfseries PROGRAMA DE DISCIPLINA} & \\
\end{tabular}
\end{center}

\begin{center}
\begin{longtable}{|p{4cm}|p{4cm}|p{4cm}|p{4cm}|}
\hline
\multicolumn{3}{|p{12cm}|}{Nome do Componente Curricular em Português:} &
\multicolumn{1}{p{4cm}|}{Código:} \\ 
\multicolumn{3}{|p{12cm}|}{\textbf{Introdução à Inteligência Artificial}} &
\textbf{BIA101}\\ 
\multicolumn{3}{|p{12cm}|}{Nome do Componente Curricular em Inglês:} & \\ 
\multicolumn{3}{|p{12cm}|}{\textbf{Introduction to Artificial Intelligence}} & \\ 
\hline
\multicolumn{3}{|p{12cm}|}{Nome e Sigla do Departamento} & Unidade Acadêmica: \\ 
\multicolumn{3}{|p{12cm}|}{Departamento de Computação (DECOM)} & {ICEB} \\ 
\hline
\multicolumn{4}{|p{16cm}|}{Modalidade de Oferta:
[X] presencial \hspace{0.6cm}
[ ] semipresencial \hspace{0.6cm}
[ ] à distância}\\
\hline
\multicolumn{2}{|p{8cm}|}{Carga horária semestral} &
\multicolumn{2}{p{8cm}|}{Carga horária semanal}\\
\hline
\multicolumn{1}{|p{4cm}|}{Total} &
\multicolumn{1}{p{4cm}|}{Extensionista} &
\multicolumn{1}{p{4cm}|}{Teórica} &
\multicolumn{1}{p{4cm}|}{Prática} \\ 
\multicolumn{1}{|p{4cm}|}{30\,horas} &
\multicolumn{1}{p{4cm}|}{0\;horas} &
\multicolumn{1}{p{4cm}|}{2\;horas/aula} &
\multicolumn{1}{p{4cm}|}{0\;horas/aula} \\ 
\hline
\multicolumn{4}{|p{16cm}|}{Ementa:}\\
\multicolumn{4}{|p{16cm}|}{}\\
\multicolumn{4}{|p{\dimexpr 16cm + 6\tabcolsep\relax}|}{O que é Inteligência Artificial; o comportamento do aluno de Inteligência Artificial; áreas de pesquisa do Departamento de Computação; áreas nas quais atuam os profissionais em Inteligência Artificial.}\\
\multicolumn{4}{|p{16cm}|}{}\\
\hline
\multicolumn{4}{|p{16cm}|}{Conteúdo programático:}\\*
\multicolumn{4}{|p{\dimexpr 16cm + 6\tabcolsep\relax}|}{\begin{enumerate}[start=1]\item A área de Inteligência Artificial e suas oportunidades atuais\end{enumerate}}\\
\multicolumn{4}{|p{\dimexpr 16cm + 6\tabcolsep\relax}|}{\begin{enumerate}[start=2]\item Áreas de atuação e mercado de trabalho em Inteligência Artificial\end{enumerate}}\\
\multicolumn{4}{|p{\dimexpr 16cm + 6\tabcolsep\relax}|}{\begin{enumerate}[start=3]\item O curso de Inteligência Artificial: grade curricular, áreas, relações entre as disciplinas\end{enumerate}}\\
\multicolumn{4}{|p{\dimexpr 16cm + 6\tabcolsep\relax}|}{\begin{enumerate}[start=4]\item Organização política da universidade e institutos/unidades\end{enumerate}}\\
\multicolumn{4}{|p{\dimexpr 16cm + 6\tabcolsep\relax}|}{\begin{enumerate}[start=5]\item Representação acadêmica: centro acadêmico e movimentos estudantis\end{enumerate}}\\
\multicolumn{4}{|p{\dimexpr 16cm + 6\tabcolsep\relax}|}{\begin{enumerate}[start=6]\item Áreas de pesquisa e extensão\end{enumerate}}\\
\multicolumn{4}{|p{\dimexpr 16cm + 6\tabcolsep\relax}|}{\begin{enumerate}[start=7]\item Laboratórios de pesquisa e extensão\end{enumerate}}\\
\multicolumn{4}{|p{\dimexpr 16cm + 6\tabcolsep\relax}|}{\begin{enumerate}[start=8]\item O comportamento do aluno de Inteligência Artificial: organização e conselhos para estudar melhor e ter sucesso no curso\end{enumerate}}\\
\multicolumn{4}{|p{\dimexpr 16cm + 6\tabcolsep\relax}|}{\begin{enumerate}[start=9]\item Seminários de profissionais sobre o mercado de trabalho e carreiras na área de Inteligência Artificial\end{enumerate}}\\
\multicolumn{4}{|p{\dimexpr 16cm + 6\tabcolsep\relax}|}{\begin{enumerate}[start=10]\item Implicações éticas, sociais e ambientais do uso de Inteligência Artificial\end{enumerate}}\\
\multicolumn{4}{|p{\dimexpr 16cm + 6\tabcolsep\relax}|}{}\\
\hline
\multicolumn{4}{|p{16cm}|}{Bibliografia Básica:}\\*
\multicolumn{4}{|p{\dimexpr 16cm + 6\tabcolsep\relax}|}{\begin{itemize}\item Morais, Regis de (org.). Filosofia da ciência e da tecnologia: introdução metodológica e crítica. 1. ed. Campinas: Papirus, 2013. E-book.\end{itemize}}\\
\multicolumn{4}{|p{\dimexpr 16cm + 6\tabcolsep\relax}|}{\begin{itemize}\item Medeiros, Luciano Frontino de. Inteligência artificial aplicada: uma abordagem introdutória. Curitiba, PR: Intersaberes, 2018. E-book.\end{itemize}}\\
\multicolumn{4}{|p{\dimexpr 16cm + 6\tabcolsep\relax}|}{\begin{itemize}\item Dias, Ana Francisca Pinto et al.; Guimarães, João Alexandre Silva Alves; Alves, Rodrigo Vitorino Souza (org.). Os direitos humanos e a ética na era da inteligência artificial. Indaiatuba, SP: Foco, 2023. E-book.\end{itemize}}\\
\multicolumn{4}{|p{\dimexpr 16cm + 6\tabcolsep\relax}|}{\begin{itemize}\item Poole, David L.; Mackworth, Alan K. Artificial Intelligence: Foundations of Computational Agents. 3. ed. Cambridge: Cambridge University Press, 2023. ISBN 978-1-009-25819-7.\end{itemize}}\\
\multicolumn{4}{|p{\dimexpr 16cm + 6\tabcolsep\relax}|}{\begin{itemize}\item Crawford, Kate. Atlas da I.A.: poder, política e os custos planetários da inteligência artificial. São Paulo: Edições Sesc São Paulo, 2025. 138 p. ISBN 978-85-9493-328-7.\end{itemize}}\\
\multicolumn{4}{|p{16cm}|}{}\\
\hline
\multicolumn{4}{|p{16cm}|}{Bibliografia Complementar:}\\*
\multicolumn{4}{|p{\dimexpr 16cm + 6\tabcolsep\relax}|}{\begin{itemize}\item Krelling Neto, Antonio Osmar. Responsabilidade civil: cibercrimes. 1. ed. São Paulo: Contentus, 2020. E-book.\end{itemize}}\\
\multicolumn{4}{|p{\dimexpr 16cm + 6\tabcolsep\relax}|}{\begin{itemize}\item Taurion, Cezar. Big data. 1. ed. Rio de Janeiro: Brasport, 2013. E-book.\end{itemize}}\\
\multicolumn{4}{|p{\dimexpr 16cm + 6\tabcolsep\relax}|}{\begin{itemize}\item Munhoz, Antonio Siemsen. Responsabilidade e autoridade social das empresas. 1. ed. Curitiba: Intersaberes, 2015. E-book.\end{itemize}}\\
\multicolumn{4}{|p{\dimexpr 16cm + 6\tabcolsep\relax}|}{\begin{itemize}\item Flores, Márcio José das; Bess, Alexandre Leal. Inteligência artificial aplicada a negócios. Curitiba, PR: Intersaberes, 2023. E-book.\end{itemize}}\\
\multicolumn{4}{|p{\dimexpr 16cm + 6\tabcolsep\relax}|}{\begin{itemize}\item Muniz, Antonio et al. Inteligência artificial: entenda como a IA pode impactar no mercado de trabalho e na sociedade. [S.l.]: Brasport, 2024. E-book.\end{itemize}}\\
\hline
\end{longtable}
\end{center}

\clearpage
\begin{center}
\begin{tabular}{lccr}
 \multirow{3}{*}{\includegraphics[height=2.7cm]{brasao.png}} &
 \multicolumn{2}{c}{\bfseries UNIVERSIDADE FEDERAL DE OURO PRETO} &
 \ \ \ \ \multirow{3}{*}{\includegraphics[height=2.7cm]{ufop.png}} \\
 & \multicolumn{2}{c}{\bfseries PRÓ-REITORIA DE GRADUAÇÃO} & \\
 & \multicolumn{2}{c}{} & \\
 & \multicolumn{2}{c}{\Large\bfseries PROGRAMA DE DISCIPLINA} & \\
\end{tabular}
\end{center}

\begin{center}
\begin{longtable}{|p{4cm}|p{4cm}|p{4cm}|p{4cm}|}
\hline
\multicolumn{3}{|p{12cm}|}{Nome do Componente Curricular em Português:} &
\multicolumn{1}{p{4cm}|}{Código:} \\ 
\multicolumn{3}{|p{12cm}|}{\textbf{Introdução à Programação}} &
\textbf{BCC201}\\ 
\multicolumn{3}{|p{12cm}|}{Nome do Componente Curricular em Inglês:} & \\ 
\multicolumn{3}{|p{12cm}|}{\textbf{Introduction to Programming}} & \\ 
\hline
\multicolumn{3}{|p{12cm}|}{Nome e Sigla do Departamento} & Unidade Acadêmica: \\ 
\multicolumn{3}{|p{12cm}|}{Departamento de Computação (DECOM)} & {ICEB} \\ 
\hline
\multicolumn{4}{|p{16cm}|}{Modalidade de Oferta:
[X] presencial \hspace{0.6cm}
[ ] semipresencial \hspace{0.6cm}
[ ] à distância}\\
\hline
\multicolumn{2}{|p{8cm}|}{Carga horária semestral} &
\multicolumn{2}{p{8cm}|}{Carga horária semanal}\\
\hline
\multicolumn{1}{|p{4cm}|}{Total} &
\multicolumn{1}{p{4cm}|}{Extensionista} &
\multicolumn{1}{p{4cm}|}{Teórica} &
\multicolumn{1}{p{4cm}|}{Prática} \\ 
\multicolumn{1}{|p{4cm}|}{90\,horas} &
\multicolumn{1}{p{4cm}|}{0\;horas} &
\multicolumn{1}{p{4cm}|}{4\;horas/aula} &
\multicolumn{1}{p{4cm}|}{2\;horas/aula} \\ 
\hline
\multicolumn{4}{|p{16cm}|}{Ementa:}\\
\multicolumn{4}{|p{16cm}|}{}\\
\multicolumn{4}{|p{\dimexpr 16cm + 6\tabcolsep\relax}|}{Introdução à lógica de programação; conceitos básicos sobre algoritmos, utilização e formas de representação (fluxograma e Portugol); tipos de dados; variáveis e constantes; expressões e operadores relacionais, aritméticos e lógicos; estruturas condicionais e de repetição; subprogramação: modularização de programas (funções e procedimentos); estruturas de dados homogêneas (vetores e matrizes) e heterogêneas (estruturas/registros); manipulação de cadeias de caracteres; ponteiros; alocação dinâmica de memória; processamento de arquivos.}\\
\multicolumn{4}{|p{16cm}|}{}\\
\hline
\multicolumn{4}{|p{16cm}|}{Conteúdo programático:}\\*
\multicolumn{4}{|p{\dimexpr 16cm + 6\tabcolsep\relax}|}{\begin{enumerate}[start=1]\item Representação de dados\end{enumerate}}\\
\multicolumn{4}{|p{\dimexpr 16cm + 6\tabcolsep\relax}|}{\begin{enumerate}[start=2]\item Conceitos e representação de algoritmos\end{enumerate}}\\
\multicolumn{4}{|p{\dimexpr 16cm + 6\tabcolsep\relax}|}{\begin{enumerate}[start=3]\item Fluxograma e Portugol\end{enumerate}}\\
\multicolumn{4}{|p{\dimexpr 16cm + 6\tabcolsep\relax}|}{\begin{enumerate}[start=4]\item Conceitos básicos de programação: valores, tipos e expressões\end{enumerate}}\\
\multicolumn{4}{|p{\dimexpr 16cm + 6\tabcolsep\relax}|}{\begin{enumerate}[start=5]\item Variáveis, comandos de atribuição e de entrada e saída\end{enumerate}}\\
\multicolumn{4}{|p{\dimexpr 16cm + 6\tabcolsep\relax}|}{\begin{enumerate}[start=6]\item Comandos de controle de fluxo\end{enumerate}}\\
\multicolumn{4}{|p{\dimexpr 16cm + 6\tabcolsep\relax}|}{\begin{enumerate}[start=7]\item Comando de decisão (if)\end{enumerate}}\\
\multicolumn{4}{|p{\dimexpr 16cm + 6\tabcolsep\relax}|}{\begin{enumerate}[start=8]\item Comandos de decisão múltipla e de salto (switch, break)\end{enumerate}}\\
\multicolumn{4}{|p{\dimexpr 16cm + 6\tabcolsep\relax}|}{\begin{enumerate}[start=9]\item Comandos de repetição (while, do-while, for)\end{enumerate}}\\
\multicolumn{4}{|p{\dimexpr 16cm + 6\tabcolsep\relax}|}{\begin{enumerate}[start=10]\item Subprogramação: funções, procedimentos e parâmetros\end{enumerate}}\\
\multicolumn{4}{|p{\dimexpr 16cm + 6\tabcolsep\relax}|}{\begin{enumerate}[start=11]\item Estruturas de dados homogêneas (vetores)\end{enumerate}}\\
\multicolumn{4}{|p{\dimexpr 16cm + 6\tabcolsep\relax}|}{\begin{enumerate}[start=12]\item Cadeia de caracteres (strings)\end{enumerate}}\\
\multicolumn{4}{|p{\dimexpr 16cm + 6\tabcolsep\relax}|}{\begin{enumerate}[start=13]\item Estruturas de dados homogêneas (matrizes)\end{enumerate}}\\
\multicolumn{4}{|p{\dimexpr 16cm + 6\tabcolsep\relax}|}{\begin{enumerate}[start=14]\item Estruturas de dados heterogêneas\end{enumerate}}\\
\multicolumn{4}{|p{\dimexpr 16cm + 6\tabcolsep\relax}|}{\begin{enumerate}[start=15]\item Apontadores e memória dinâmica (ponteiros)\end{enumerate}}\\
\multicolumn{4}{|p{\dimexpr 16cm + 6\tabcolsep\relax}|}{\begin{enumerate}[start=16]\item Arquivos\end{enumerate}}\\
\multicolumn{4}{|p{\dimexpr 16cm + 6\tabcolsep\relax}|}{}\\
\hline
\multicolumn{4}{|p{16cm}|}{Bibliografia Básica:}\\*
\multicolumn{4}{|p{\dimexpr 16cm + 6\tabcolsep\relax}|}{\begin{itemize}\item DEITEL, P.; DEITEL, H. M. C: como programar. 6. ed. São Paulo: Pearson Education, 2011.\end{itemize}}\\
\multicolumn{4}{|p{\dimexpr 16cm + 6\tabcolsep\relax}|}{\begin{itemize}\item DEITEL, H. M.; DEITEL, P. J. C++: como programar. 5. ed. São Paulo: Pearson Prentice Hall, 2006.\end{itemize}}\\
\multicolumn{4}{|p{\dimexpr 16cm + 6\tabcolsep\relax}|}{\begin{itemize}\item SOUZA, M. A. F. de. Algoritmos e lógica de programação. São Paulo: Cengage Learning, 2005.\end{itemize}}\\
\multicolumn{4}{|p{16cm}|}{}\\
\hline
\multicolumn{4}{|p{16cm}|}{Bibliografia Complementar:}\\*
\multicolumn{4}{|p{\dimexpr 16cm + 6\tabcolsep\relax}|}{\begin{itemize}\item ASCENCIO, A. F. G.; CAMPOS, E. A. V. Fundamentos da programação de computadores: algoritmos, Pascal e C/C++. 2. ed. São Paulo: Prentice Hall, 2007.\end{itemize}}\\
\multicolumn{4}{|p{\dimexpr 16cm + 6\tabcolsep\relax}|}{\begin{itemize}\item GUEDES, S. Lógica de programação algorítmica. São Paulo: Pearson Education do Brasil, 2014.\end{itemize}}\\
\multicolumn{4}{|p{\dimexpr 16cm + 6\tabcolsep\relax}|}{\begin{itemize}\item MIZRAHI, V. V. Treinamento em linguagem C. 2. ed. São Paulo: Pearson Education, 2010.\end{itemize}}\\
\multicolumn{4}{|p{\dimexpr 16cm + 6\tabcolsep\relax}|}{\begin{itemize}\item MIZRAHI, V. V. Treinamento em linguagem C++: módulo 2. São Paulo: Pearson Education, 2006.\end{itemize}}\\
\multicolumn{4}{|p{\dimexpr 16cm + 6\tabcolsep\relax}|}{\begin{itemize}\item SAVITCH, W. J. C++ absoluto. São Paulo: Pearson Education; Addison-Wesley, 2004.\end{itemize}}\\
\hline
\end{longtable}
\end{center}

\clearpage
\begin{center}
\begin{tabular}{lccr}
 \multirow{3}{*}{\includegraphics[height=2.7cm]{brasao.png}} &
 \multicolumn{2}{c}{\bfseries UNIVERSIDADE FEDERAL DE OURO PRETO} &
 \ \ \ \ \multirow{3}{*}{\includegraphics[height=2.7cm]{ufop.png}} \\
 & \multicolumn{2}{c}{\bfseries PRÓ-REITORIA DE GRADUAÇÃO} & \\
 & \multicolumn{2}{c}{} & \\
 & \multicolumn{2}{c}{\Large\bfseries PROGRAMA DE DISCIPLINA} & \\
\end{tabular}
\end{center}

\begin{center}
\begin{longtable}{|p{4cm}|p{4cm}|p{4cm}|p{4cm}|}
\hline
\multicolumn{3}{|p{12cm}|}{Nome do Componente Curricular em Português:} &
\multicolumn{1}{p{4cm}|}{Código:} \\ 
\multicolumn{3}{|p{12cm}|}{\textbf{Teoria dos Grafos}} &
\textbf{BCC204}\\ 
\multicolumn{3}{|p{12cm}|}{Nome do Componente Curricular em Inglês:} & \\ 
\multicolumn{3}{|p{12cm}|}{\textbf{Graph Theory}} & \\ 
\hline
\multicolumn{3}{|p{12cm}|}{Nome e Sigla do Departamento} & Unidade Acadêmica: \\ 
\multicolumn{3}{|p{12cm}|}{Departamento de Computação (DECOM)} & {ICEB} \\ 
\hline
\multicolumn{4}{|p{16cm}|}{Modalidade de Oferta:
[X] presencial \hspace{0.6cm}
[ ] semipresencial \hspace{0.6cm}
[ ] à distância}\\
\hline
\multicolumn{2}{|p{8cm}|}{Carga horária semestral} &
\multicolumn{2}{p{8cm}|}{Carga horária semanal}\\
\hline
\multicolumn{1}{|p{4cm}|}{Total} &
\multicolumn{1}{p{4cm}|}{Extensionista} &
\multicolumn{1}{p{4cm}|}{Teórica} &
\multicolumn{1}{p{4cm}|}{Prática} \\ 
\multicolumn{1}{|p{4cm}|}{60\,horas} &
\multicolumn{1}{p{4cm}|}{0\;horas} &
\multicolumn{1}{p{4cm}|}{4\;horas/aula} &
\multicolumn{1}{p{4cm}|}{0\;horas/aula} \\ 
\hline
\multicolumn{4}{|p{16cm}|}{Ementa:}\\
\multicolumn{4}{|p{16cm}|}{}\\
\multicolumn{4}{|p{\dimexpr 16cm + 6\tabcolsep\relax}|}{Grafos orientados e não-orientados; caminhos; planaridade; conectividade; coloração; grafos infinitos; problemas intratáveis; busca em largura e profundidade; algoritmos do menor caminho; árvore geradora; ordenação topológica.}\\
\multicolumn{4}{|p{16cm}|}{}\\
\hline
\multicolumn{4}{|p{16cm}|}{Conteúdo programático:}\\*
\multicolumn{4}{|p{\dimexpr 16cm + 6\tabcolsep\relax}|}{\begin{enumerate}[start=1]\item Introdução e estruturas de dados para grafos\end{enumerate}}\\
\multicolumn{4}{|p{\dimexpr 16cm + 6\tabcolsep\relax}|}{\begin{enumerate}[start=2]\item Formalização: definições\end{enumerate}}\\
\multicolumn{4}{|p{\dimexpr 16cm + 6\tabcolsep\relax}|}{\begin{enumerate}[start=3]\item Isomorfismo\end{enumerate}}\\
\multicolumn{4}{|p{\dimexpr 16cm + 6\tabcolsep\relax}|}{\begin{enumerate}[start=4]\item Complementaridade e subgrafos\end{enumerate}}\\
\multicolumn{4}{|p{\dimexpr 16cm + 6\tabcolsep\relax}|}{\begin{enumerate}[start=5]\item Teorema do aperto de mãos e bipartição\end{enumerate}}\\
\multicolumn{4}{|p{\dimexpr 16cm + 6\tabcolsep\relax}|}{\begin{enumerate}[start=6]\item Passeio, cadeia e caminho\end{enumerate}}\\
\multicolumn{4}{|p{\dimexpr 16cm + 6\tabcolsep\relax}|}{\begin{enumerate}[start=7]\item Transitividade e conectividade\end{enumerate}}\\
\multicolumn{4}{|p{\dimexpr 16cm + 6\tabcolsep\relax}|}{\begin{enumerate}[start=8]\item Busca em grafos: busca em profundidade e largura\end{enumerate}}\\
\multicolumn{4}{|p{\dimexpr 16cm + 6\tabcolsep\relax}|}{\begin{enumerate}[start=9]\item Algoritmos de caminhos mínimos: Dijkstra, Bellman-Ford, Floyd-Warshall\end{enumerate}}\\
\multicolumn{4}{|p{\dimexpr 16cm + 6\tabcolsep\relax}|}{\begin{enumerate}[start=10]\item Ordenação topológica\end{enumerate}}\\
\multicolumn{4}{|p{\dimexpr 16cm + 6\tabcolsep\relax}|}{\begin{enumerate}[start=11]\item Fluxo em redes: Ford-Fulkerson\end{enumerate}}\\
\multicolumn{4}{|p{\dimexpr 16cm + 6\tabcolsep\relax}|}{\begin{enumerate}[start=12]\item Problemas Intratáveis\end{enumerate}}\\
\multicolumn{4}{|p{\dimexpr 16cm + 6\tabcolsep\relax}|}{\begin{enumerate}[start=13]\item Casamento em grafos e Algoritmo Húngaro\end{enumerate}}\\
\multicolumn{4}{|p{\dimexpr 16cm + 6\tabcolsep\relax}|}{\begin{enumerate}[start=14]\item Conjuntos independentes, cliques e conjuntos dominantes\end{enumerate}}\\
\multicolumn{4}{|p{\dimexpr 16cm + 6\tabcolsep\relax}|}{\begin{enumerate}[start=15]\item O problema das 4 cores: coloração de mapas\end{enumerate}}\\
\multicolumn{4}{|p{\dimexpr 16cm + 6\tabcolsep\relax}|}{\begin{enumerate}[start=16]\item Coloração de grafos\end{enumerate}}\\
\multicolumn{4}{|p{\dimexpr 16cm + 6\tabcolsep\relax}|}{\begin{enumerate}[start=17]\item Planaridade em grafos\end{enumerate}}\\
\multicolumn{4}{|p{\dimexpr 16cm + 6\tabcolsep\relax}|}{\begin{enumerate}[start=18]\item Busca de soluções usando grafos\end{enumerate}}\\
\multicolumn{4}{|p{\dimexpr 16cm + 6\tabcolsep\relax}|}{}\\
\hline
\multicolumn{4}{|p{16cm}|}{Bibliografia Básica:}\\*
\multicolumn{4}{|p{\dimexpr 16cm + 6\tabcolsep\relax}|}{\begin{itemize}\item Goldbarg, Marco; Goldbarg, Elizabeth. Grafos: conceitos, algoritmos e aplicações. Rio de Janeiro: Elsevier, 2012.\end{itemize}}\\
\multicolumn{4}{|p{\dimexpr 16cm + 6\tabcolsep\relax}|}{\begin{itemize}\item Boaventura Netto, Paulo Oswaldo. Grafos: teoria, modelos, algoritmos. 2. ed. rev. e ampl. São Paulo: Edgard Blucher, 2001.\end{itemize}}\\
\multicolumn{4}{|p{16cm}|}{}\\
\hline
\multicolumn{4}{|p{16cm}|}{Bibliografia Complementar:}\\*
\multicolumn{4}{|p{\dimexpr 16cm + 6\tabcolsep\relax}|}{\begin{itemize}\item Gross, Jonathan L. Graph theory and its applications. 2. ed. Boca Raton: CRC Press, 2006.\end{itemize}}\\
\multicolumn{4}{|p{\dimexpr 16cm + 6\tabcolsep\relax}|}{\begin{itemize}\item Jungnickel, D. Graphs, networks, and algorithms. 3. ed. Berlin.\end{itemize}}\\
\multicolumn{4}{|p{\dimexpr 16cm + 6\tabcolsep\relax}|}{\begin{itemize}\item Wilson, Robin J. Introduction to graph theory. 3. ed. New York: John Wiley \& Sons, 1990.\end{itemize}}\\
\multicolumn{4}{|p{\dimexpr 16cm + 6\tabcolsep\relax}|}{\begin{itemize}\item New directions in the theory of graphs. New York.\end{itemize}}\\
\multicolumn{4}{|p{\dimexpr 16cm + 6\tabcolsep\relax}|}{\begin{itemize}\item Gibbons, Alan. Algorithmic graph theory. New York: Cambridge, 1994.\end{itemize}}\\
\multicolumn{4}{|p{\dimexpr 16cm + 6\tabcolsep\relax}|}{\begin{itemize}\item Cormen, Thomas H. et al.; tradução Vandenberg D. de Souza. Algoritmos: teoria e prática. Rio de Janeiro: Campus, 2002.\end{itemize}}\\
\hline
\end{longtable}
\end{center}

\clearpage
\begin{center}
\begin{tabular}{lccr}
 \multirow{3}{*}{\includegraphics[height=2.7cm]{brasao.png}} &
 \multicolumn{2}{c}{\bfseries UNIVERSIDADE FEDERAL DE OURO PRETO} &
 \ \ \ \ \multirow{3}{*}{\includegraphics[height=2.7cm]{ufop.png}} \\
 & \multicolumn{2}{c}{\bfseries PRÓ-REITORIA DE GRADUAÇÃO} & \\
 & \multicolumn{2}{c}{} & \\
 & \multicolumn{2}{c}{\Large\bfseries PROGRAMA DE DISCIPLINA} & \\
\end{tabular}
\end{center}

\begin{center}
\begin{longtable}{|p{4cm}|p{4cm}|p{4cm}|p{4cm}|}
\hline
\multicolumn{3}{|p{12cm}|}{Nome do Componente Curricular em Português:} &
\multicolumn{1}{p{4cm}|}{Código:} \\ 
\multicolumn{3}{|p{12cm}|}{\textbf{Matemática Discreta I}} &
\textbf{BCC101}\\ 
\multicolumn{3}{|p{12cm}|}{Nome do Componente Curricular em Inglês:} & \\ 
\multicolumn{3}{|p{12cm}|}{\textbf{Discrete Mathematics I}} & \\ 
\hline
\multicolumn{3}{|p{12cm}|}{Nome e Sigla do Departamento} & Unidade Acadêmica: \\ 
\multicolumn{3}{|p{12cm}|}{Departamento de Computação (DECOM)} & {ICEB} \\ 
\hline
\multicolumn{4}{|p{16cm}|}{Modalidade de Oferta:
[X] presencial \hspace{0.6cm}
[ ] semipresencial \hspace{0.6cm}
[ ] à distância}\\
\hline
\multicolumn{2}{|p{8cm}|}{Carga horária semestral} &
\multicolumn{2}{p{8cm}|}{Carga horária semanal}\\
\hline
\multicolumn{1}{|p{4cm}|}{Total} &
\multicolumn{1}{p{4cm}|}{Extensionista} &
\multicolumn{1}{p{4cm}|}{Teórica} &
\multicolumn{1}{p{4cm}|}{Prática} \\ 
\multicolumn{1}{|p{4cm}|}{60\,horas} &
\multicolumn{1}{p{4cm}|}{0\;horas} &
\multicolumn{1}{p{4cm}|}{4\;horas/aula} &
\multicolumn{1}{p{4cm}|}{0\;horas/aula} \\ 
\hline
\multicolumn{4}{|p{16cm}|}{Ementa:}\\
\multicolumn{4}{|p{16cm}|}{}\\
\multicolumn{4}{|p{\dimexpr 16cm + 6\tabcolsep\relax}|}{Introdução à teoria de conjuntos: definições de conjuntos, operações sobre conjuntos, cardinalidade de conjuntos. Funções: conceitos básicos, composição, funções recursivas. Lógica proposicional e lógica de predicados: sintaxe, semântica e sistema de dedução. Estratégias de prova. Indução e recursão.}\\
\multicolumn{4}{|p{16cm}|}{}\\
\hline
\multicolumn{4}{|p{16cm}|}{Conteúdo programático:}\\*
\multicolumn{4}{|p{\dimexpr 16cm + 6\tabcolsep\relax}|}{\begin{enumerate}[start=1]\item Introdução e Revisão de Teoria de Conjuntos\end{enumerate}}\\
\multicolumn{4}{|p{\dimexpr 16cm + 6\tabcolsep\relax}|}{\begin{enumerate}[start=2]\item Sintaxe e Semântica da Lógica Proposicional\end{enumerate}}\\
\multicolumn{4}{|p{\dimexpr 16cm + 6\tabcolsep\relax}|}{\begin{enumerate}[start=3]\item Sistema de Dedução da Lógica Proposicional\end{enumerate}}\\
\multicolumn{4}{|p{\dimexpr 16cm + 6\tabcolsep\relax}|}{\begin{enumerate}[start=4]\item Álgebra Booleana\end{enumerate}}\\
\multicolumn{4}{|p{\dimexpr 16cm + 6\tabcolsep\relax}|}{\begin{enumerate}[start=5]\item Sintaxe e Semântica da Lógica de Predicados\end{enumerate}}\\
\multicolumn{4}{|p{\dimexpr 16cm + 6\tabcolsep\relax}|}{\begin{enumerate}[start=6]\item Sistema de Dedução - Lógica de Predicados\end{enumerate}}\\
\multicolumn{4}{|p{\dimexpr 16cm + 6\tabcolsep\relax}|}{\begin{enumerate}[start=7]\item Álgebra de Predicados\end{enumerate}}\\
\multicolumn{4}{|p{\dimexpr 16cm + 6\tabcolsep\relax}|}{\begin{enumerate}[start=8]\item Estratégias de prova\end{enumerate}}\\
\multicolumn{4}{|p{\dimexpr 16cm + 6\tabcolsep\relax}|}{\begin{enumerate}[start=9]\item Indução e Recursão\end{enumerate}}\\
\multicolumn{4}{|p{\dimexpr 16cm + 6\tabcolsep\relax}|}{\begin{enumerate}[start=10]\item Provas e correção de provas\end{enumerate}}\\
\multicolumn{4}{|p{\dimexpr 16cm + 6\tabcolsep\relax}|}{}\\
\hline
\multicolumn{4}{|p{16cm}|}{Bibliografia Básica:}\\*
\multicolumn{4}{|p{\dimexpr 16cm + 6\tabcolsep\relax}|}{\begin{itemize}\item VELLEMAN, Daniel. J. How to Prove it: A Structured Approach. Cambridge: Cambridge University Press, 2006.\end{itemize}}\\
\multicolumn{4}{|p{\dimexpr 16cm + 6\tabcolsep\relax}|}{\begin{itemize}\item ROSEN, Kenneth. H. Matemática Discreta e suas Aplicações. 6.ed. São Paulo: McGraw-Hill, 2009.\end{itemize}}\\
\multicolumn{4}{|p{\dimexpr 16cm + 6\tabcolsep\relax}|}{\begin{itemize}\item O´DONNELL, Jonh.; HALL, Cordelia.; PAGE, Rex. Discrete Mathematics Using a Computer. Glasgow: Springer-Verlag, 2000.\end{itemize}}\\
\multicolumn{4}{|p{16cm}|}{}\\
\hline
\multicolumn{4}{|p{16cm}|}{Bibliografia Complementar:}\\*
\multicolumn{4}{|p{\dimexpr 16cm + 6\tabcolsep\relax}|}{\begin{itemize}\item HUTH, Michael; RYAN, Mark. Lógica em Ciência da Computação: Modelagem e Argumentação sobre Sistemas. 2. ed. Rio de Janeiro: LTC, 2008.\end{itemize}}\\
\multicolumn{4}{|p{\dimexpr 16cm + 6\tabcolsep\relax}|}{\begin{itemize}\item SCHEINERMAN, Edward. R. Matemática Discreta: Uma Introdução. São Paulo: Cengage Learning, 2011.\end{itemize}}\\
\multicolumn{4}{|p{\dimexpr 16cm + 6\tabcolsep\relax}|}{\begin{itemize}\item GERSTING, Judith L. Fundamentos Matemáticos para a Ciência da Computação. 5. ed. Rio de Janeiro: LTC, 2004.\end{itemize}}\\
\hline
\end{longtable}
\end{center}

\clearpage
\begin{center}
\begin{tabular}{lccr}
 \multirow{3}{*}{\includegraphics[height=2.7cm]{brasao.png}} &
 \multicolumn{2}{c}{\bfseries UNIVERSIDADE FEDERAL DE OURO PRETO} &
 \ \ \ \ \multirow{3}{*}{\includegraphics[height=2.7cm]{ufop.png}} \\
 & \multicolumn{2}{c}{\bfseries PRÓ-REITORIA DE GRADUAÇÃO} & \\
 & \multicolumn{2}{c}{} & \\
 & \multicolumn{2}{c}{\Large\bfseries PROGRAMA DE DISCIPLINA} & \\
\end{tabular}
\end{center}

\begin{center}
\begin{longtable}{|p{4cm}|p{4cm}|p{4cm}|p{4cm}|}
\hline
\multicolumn{3}{|p{12cm}|}{Nome do Componente Curricular em Português:} &
\multicolumn{1}{p{4cm}|}{Código:} \\ 
\multicolumn{3}{|p{12cm}|}{\textbf{Inteligência Artificial Clássica}} &
\textbf{BIA102}\\ 
\multicolumn{3}{|p{12cm}|}{Nome do Componente Curricular em Inglês:} & \\ 
\multicolumn{3}{|p{12cm}|}{\textbf{Classical Artificial Intelligence}} & \\ 
\hline
\multicolumn{3}{|p{12cm}|}{Nome e Sigla do Departamento} & Unidade Acadêmica: \\ 
\multicolumn{3}{|p{12cm}|}{Departamento de Computação (DECOM)} & {ICEB} \\ 
\hline
\multicolumn{4}{|p{16cm}|}{Modalidade de Oferta:
[X] presencial \hspace{0.6cm}
[ ] semipresencial \hspace{0.6cm}
[ ] à distância}\\
\hline
\multicolumn{2}{|p{8cm}|}{Carga horária semestral} &
\multicolumn{2}{p{8cm}|}{Carga horária semanal}\\
\hline
\multicolumn{1}{|p{4cm}|}{Total} &
\multicolumn{1}{p{4cm}|}{Extensionista} &
\multicolumn{1}{p{4cm}|}{Teórica} &
\multicolumn{1}{p{4cm}|}{Prática} \\ 
\multicolumn{1}{|p{4cm}|}{60\,horas} &
\multicolumn{1}{p{4cm}|}{0\;horas} &
\multicolumn{1}{p{4cm}|}{4\;horas/aula} &
\multicolumn{1}{p{4cm}|}{0\;horas/aula} \\ 
\hline
\multicolumn{4}{|p{16cm}|}{Ementa:}\\
\multicolumn{4}{|p{16cm}|}{}\\
\multicolumn{4}{|p{\dimexpr 16cm + 6\tabcolsep\relax}|}{Estudo dos fundamentos da Inteligência Artificial baseada em agentes. Resolução de problemas por busca, satisfação de restrições, raciocínio lógico e probabilístico, planejamento determinístico e sob incerteza. Representação de conhecimento, inferência e tomada de decisão. Introdução a sistemas multiagentes e discussão de aplicações, limitações e impactos sociais da Inteligência Artificial.}\\
\multicolumn{4}{|p{16cm}|}{}\\
\hline
\multicolumn{4}{|p{16cm}|}{Conteúdo programático:}\\*
\multicolumn{4}{|p{\dimexpr 16cm + 6\tabcolsep\relax}|}{\begin{enumerate}[start=1]\item Inteligência Artificial e Agentes.\end{enumerate}}\\
\multicolumn{4}{|p{\dimexpr 16cm + 6\tabcolsep\relax}|}{\begin{enumerate}[start=2]\item Resolução de Problemas por Busca\end{enumerate}}\\
\multicolumn{4}{|p{\dimexpr 16cm + 6\tabcolsep\relax}|}{\begin{enumerate}[start=3]\item Problemas de Satisfação de Restrições\end{enumerate}}\\
\multicolumn{4}{|p{\dimexpr 16cm + 6\tabcolsep\relax}|}{\begin{enumerate}[start=4]\item Proposições e Inferência\end{enumerate}}\\
\multicolumn{4}{|p{\dimexpr 16cm + 6\tabcolsep\relax}|}{\begin{enumerate}[start=5]\item Planejamento Determinístico\end{enumerate}}\\
\multicolumn{4}{|p{\dimexpr 16cm + 6\tabcolsep\relax}|}{\begin{enumerate}[start=6]\item Raciocínio Probabilístico\end{enumerate}}\\
\multicolumn{4}{|p{\dimexpr 16cm + 6\tabcolsep\relax}|}{\begin{enumerate}[start=7]\item Planejamento com Incerteza\end{enumerate}}\\
\multicolumn{4}{|p{\dimexpr 16cm + 6\tabcolsep\relax}|}{\begin{enumerate}[start=8]\item Sistemas Multiagentes\end{enumerate}}\\
\multicolumn{4}{|p{\dimexpr 16cm + 6\tabcolsep\relax}|}{\begin{enumerate}[start=9]\item Impacto Social da Inteligência Artificial\end{enumerate}}\\
\multicolumn{4}{|p{\dimexpr 16cm + 6\tabcolsep\relax}|}{}\\
\hline
\multicolumn{4}{|p{16cm}|}{Bibliografia Básica:}\\*
\multicolumn{4}{|p{\dimexpr 16cm + 6\tabcolsep\relax}|}{\begin{itemize}\item POOLE, David L.; MACKWORTH, Alan K. Artificial Intelligence: Foundations of Computational Agents. 3. ed. Cambridge: Cambridge University Press, 2023. ISBN 978-1-009-25819-7. Disponível em: <https://artint.info/>. Acesso em: 17 dez. 2025.\end{itemize}}\\
\multicolumn{4}{|p{\dimexpr 16cm + 6\tabcolsep\relax}|}{\begin{itemize}\item RUSSELL, Stuart; NORVIG, Peter. Inteligência artificial. 3. ed. Rio de Janeiro: Elsevier, 2013. E-book. Disponível em: https://plataforma.bvirtual.com.br. Acesso em: jun/2025.\end{itemize}}\\
\multicolumn{4}{|p{\dimexpr 16cm + 6\tabcolsep\relax}|}{\begin{itemize}\item BARROS, L. N. de. Introdução à inteligência artificial. 1. ed. Rio de Janeiro: LTC, 2010. E-book. Disponível em: https://plataforma.bvirtual.com.br. Acesso em: jun/2025.\end{itemize}}\\
\multicolumn{4}{|p{16cm}|}{}\\
\hline
\multicolumn{4}{|p{16cm}|}{Bibliografia Complementar:}\\*
\multicolumn{4}{|p{\dimexpr 16cm + 6\tabcolsep\relax}|}{\begin{itemize}\item MEDEIROS, Luciano Frontino de. Inteligência artificial aplicada: uma abordagem introdutória. Curitiba, PR: Intersaberes, 2018. E-book. Disponível em: https://plataforma.bvirtual.com.br. Acesso em: jun/2025.\end{itemize}}\\
\multicolumn{4}{|p{\dimexpr 16cm + 6\tabcolsep\relax}|}{\begin{itemize}\item BARR, Avron; FEIGENBAUM, Edward A. (ed.). The Handbook of Artificial Intelligence. Volume I. Stanford, CA: HeurisTech Press; Los Altos, CA: William Kaufmann, 1981.\end{itemize}}\\
\multicolumn{4}{|p{\dimexpr 16cm + 6\tabcolsep\relax}|}{\begin{itemize}\item LUGER, George F. Inteligência artificial: estruturas e estratégias para a solução de problemas complexos. 6. ed. Porto Alegre: Bookman, 2004. E-book. Disponível em: https://plataforma.bvirtual.com.br. Acesso em: jun/2025.\end{itemize}}\\
\multicolumn{4}{|p{\dimexpr 16cm + 6\tabcolsep\relax}|}{\begin{itemize}\item DIAS, Ana Francisca Pinto et al. (org.). Os direitos humanos e a ética na era da inteligência artificial. Indaiatuba, SP: Foco, 2023. E-book. Disponível em: https://plataforma.bvirtual.com.br. Acesso em: jun/2025.\end{itemize}}\\
\multicolumn{4}{|p{\dimexpr 16cm + 6\tabcolsep\relax}|}{\begin{itemize}\item LAUFER, C. Sistemas especialistas e representação do conhecimento. 1. ed. São Paulo: Érica, 2009. E-book. Disponível em: https://plataforma.bvirtual.com.br. Acesso em: jun/2025.\end{itemize}}\\
\hline
\end{longtable}
\end{center}

\clearpage
\begin{center}
\begin{tabular}{lccr}
 \multirow{3}{*}{\includegraphics[height=2.7cm]{brasao.png}} &
 \multicolumn{2}{c}{\bfseries UNIVERSIDADE FEDERAL DE OURO PRETO} &
 \ \ \ \ \multirow{3}{*}{\includegraphics[height=2.7cm]{ufop.png}} \\
 & \multicolumn{2}{c}{\bfseries PRÓ-REITORIA DE GRADUAÇÃO} & \\
 & \multicolumn{2}{c}{} & \\
 & \multicolumn{2}{c}{\Large\bfseries PROGRAMA DE DISCIPLINA} & \\
\end{tabular}
\end{center}

\begin{center}
\begin{longtable}{|p{4cm}|p{4cm}|p{4cm}|p{4cm}|}
\hline
\multicolumn{3}{|p{12cm}|}{Nome do Componente Curricular em Português:} &
\multicolumn{1}{p{4cm}|}{Código:} \\ 
\multicolumn{3}{|p{12cm}|}{\textbf{Fundamentos de Arquitetura e Organização de Computadores}} &
\textbf{BIA103}\\ 
\multicolumn{3}{|p{12cm}|}{Nome do Componente Curricular em Inglês:} & \\ 
\multicolumn{3}{|p{12cm}|}{\textbf{Computer Architecture and Organization}} & \\ 
\hline
\multicolumn{3}{|p{12cm}|}{Nome e Sigla do Departamento} & Unidade Acadêmica: \\ 
\multicolumn{3}{|p{12cm}|}{Departamento de Computação (DECOM)} & {ICEB} \\ 
\hline
\multicolumn{4}{|p{16cm}|}{Modalidade de Oferta:
[X] presencial \hspace{0.6cm}
[ ] semipresencial \hspace{0.6cm}
[ ] à distância}\\
\hline
\multicolumn{2}{|p{8cm}|}{Carga horária semestral} &
\multicolumn{2}{p{8cm}|}{Carga horária semanal}\\
\hline
\multicolumn{1}{|p{4cm}|}{Total} &
\multicolumn{1}{p{4cm}|}{Extensionista} &
\multicolumn{1}{p{4cm}|}{Teórica} &
\multicolumn{1}{p{4cm}|}{Prática} \\ 
\multicolumn{1}{|p{4cm}|}{60\,horas} &
\multicolumn{1}{p{4cm}|}{0\;horas} &
\multicolumn{1}{p{4cm}|}{4\;horas/aula} &
\multicolumn{1}{p{4cm}|}{0\;horas/aula} \\ 
\hline
\multicolumn{4}{|p{16cm}|}{Ementa:}\\
\multicolumn{4}{|p{16cm}|}{}\\
\multicolumn{4}{|p{\dimexpr 16cm + 6\tabcolsep\relax}|}{Estudo dos componentes básicos de um computador e suas interações. Conceitos de organização e arquitetura de computadores. Sistemas de numeração e codificação. Conjuntos de instruções, ciclo de execução, hierarquia de memória e barramentos.}\\
\multicolumn{4}{|p{16cm}|}{}\\
\hline
\multicolumn{4}{|p{16cm}|}{Conteúdo programático:}\\*
\multicolumn{4}{|p{\dimexpr 16cm + 6\tabcolsep\relax}|}{\begin{enumerate}[start=1]\item Conceitos fundamentais: diferença entre organização e arquitetura de computadores.\end{enumerate}}\\
\multicolumn{4}{|p{\dimexpr 16cm + 6\tabcolsep\relax}|}{\begin{enumerate}[start=2]\item Sistemas de numeração: binário, octal, decimal e hexadecimal.\end{enumerate}}\\
\multicolumn{4}{|p{\dimexpr 16cm + 6\tabcolsep\relax}|}{\begin{enumerate}[start=3]\item Aritmética computacional: operações com números inteiros e de ponto flutuante.\end{enumerate}}\\
\multicolumn{4}{|p{\dimexpr 16cm + 6\tabcolsep\relax}|}{\begin{enumerate}[start=4]\item Unidade Central de Processamento (UCP): registradores, ALU, UC.\end{enumerate}}\\
\multicolumn{4}{|p{\dimexpr 16cm + 6\tabcolsep\relax}|}{\begin{enumerate}[start=5]\item Ciclo de instrução: busca, decodificação e execução.\end{enumerate}}\\
\multicolumn{4}{|p{\dimexpr 16cm + 6\tabcolsep\relax}|}{\begin{enumerate}[start=6]\item Conjuntos de instruções (ISA): instruções RISC e CISC.\end{enumerate}}\\
\multicolumn{4}{|p{\dimexpr 16cm + 6\tabcolsep\relax}|}{\begin{enumerate}[start=7]\item Hierarquia de memória: registradores, cache, RAM, memória secundária.\end{enumerate}}\\
\multicolumn{4}{|p{\dimexpr 16cm + 6\tabcolsep\relax}|}{\begin{enumerate}[start=8]\item Sistemas de entrada e saída: interfaces e dispositivos.\end{enumerate}}\\
\multicolumn{4}{|p{\dimexpr 16cm + 6\tabcolsep\relax}|}{\begin{enumerate}[start=9]\item Barramentos e comunicação entre componentes.\end{enumerate}}\\
\multicolumn{4}{|p{\dimexpr 16cm + 6\tabcolsep\relax}|}{\begin{enumerate}[start=10]\item Noções de pipelines e paralelismo em processadores modernos.\end{enumerate}}\\
\multicolumn{4}{|p{\dimexpr 16cm + 6\tabcolsep\relax}|}{}\\
\hline
\multicolumn{4}{|p{16cm}|}{Bibliografia Básica:}\\*
\multicolumn{4}{|p{\dimexpr 16cm + 6\tabcolsep\relax}|}{\begin{itemize}\item STALLINGS, William. Arquitetura e organização de computadores. 10. ed. São Paulo: Pearson, 2016. E-book. Disponível em: https://plataforma.bvirtual.com.br. Acesso em: jun/2025.\end{itemize}}\\
\multicolumn{4}{|p{\dimexpr 16cm + 6\tabcolsep\relax}|}{\begin{itemize}\item MAK, Ronald. Organização de computadores: arquitetura, interfaces e sistemas operacionais. 1. ed. São Paulo: Cengage Learning, 2020. E-book. Disponível em: https://plataforma.bvirtual.com.br. Acesso em: jun/2025.\end{itemize}}\\
\multicolumn{4}{|p{\dimexpr 16cm + 6\tabcolsep\relax}|}{\begin{itemize}\item TANENBAUM, Andrew S.; AUSTIN, Todd. Estruturas de computadores. 6. ed. São Paulo: Pearson, 2014. E-book. Disponível em: https://plataforma.bvirtual.com.br. Acesso em: jun/2025.\end{itemize}}\\
\multicolumn{4}{|p{16cm}|}{}\\
\hline
\multicolumn{4}{|p{16cm}|}{Bibliografia Complementar:}\\*
\multicolumn{4}{|p{\dimexpr 16cm + 6\tabcolsep\relax}|}{\begin{itemize}\item BRITO, Fábio de Souza. Arquitetura e organização de computadores: teoria e prática. 1. ed. São Paulo: Érica, 2022. E-book. Disponível em: https://plataforma.bvirtual.com.br. Acesso em: jun/2025.\end{itemize}}\\
\multicolumn{4}{|p{\dimexpr 16cm + 6\tabcolsep\relax}|}{\begin{itemize}\item HENNESSY, John L.; PATTERSON, David A. Organização e projeto de computadores: a interface hardware/software. 5. ed. Porto Alegre: Bookman, 2014. E-book. Disponível em: https://plataforma.bvirtual.com.br. Acesso em: jun/2025.\end{itemize}}\\
\multicolumn{4}{|p{\dimexpr 16cm + 6\tabcolsep\relax}|}{\begin{itemize}\item SANTOS, Silas P. Fundamentos de organização de computadores. 2. ed. Curitiba: Intersaberes, 2021. E-book. Disponível em: https://plataforma.bvirtual.com.br. Acesso em: jun/2025.\end{itemize}}\\
\multicolumn{4}{|p{\dimexpr 16cm + 6\tabcolsep\relax}|}{\begin{itemize}\item MACHADO, Fábio; MAIA, Marcos. Arquitetura de computadores moderna. 1. ed. Rio de Janeiro: LTC, 2019. E-book. Disponível em: https://plataforma.bvirtual.com.br. Acesso em: jun/2025.\end{itemize}}\\
\multicolumn{4}{|p{\dimexpr 16cm + 6\tabcolsep\relax}|}{\begin{itemize}\item FERREIRA, Luiz. Sistemas digitais e arquitetura de computadores. 3. ed. São Paulo: Érica, 2018. E-book. Disponível em: https://plataforma.bvirtual.com.br. Acesso em: jun/2025.\end{itemize}}\\
\hline
\end{longtable}
\end{center}

\clearpage
\begin{center}
\begin{tabular}{lccr}
 \multirow{3}{*}{\includegraphics[height=2.7cm]{brasao.png}} &
 \multicolumn{2}{c}{\bfseries UNIVERSIDADE FEDERAL DE OURO PRETO} &
 \ \ \ \ \multirow{3}{*}{\includegraphics[height=2.7cm]{ufop.png}} \\
 & \multicolumn{2}{c}{\bfseries PRÓ-REITORIA DE GRADUAÇÃO} & \\
 & \multicolumn{2}{c}{} & \\
 & \multicolumn{2}{c}{\Large\bfseries PROGRAMA DE DISCIPLINA} & \\
\end{tabular}
\end{center}

\begin{center}
\begin{longtable}{|p{4cm}|p{4cm}|p{4cm}|p{4cm}|}
\hline
\multicolumn{3}{|p{12cm}|}{Nome do Componente Curricular em Português:} &
\multicolumn{1}{p{4cm}|}{Código:} \\ 
\multicolumn{3}{|p{12cm}|}{\textbf{Estruturas de Dados I}} &
\textbf{BCC137}\\ 
\multicolumn{3}{|p{12cm}|}{Nome do Componente Curricular em Inglês:} & \\ 
\multicolumn{3}{|p{12cm}|}{\textbf{Data Structures I}} & \\ 
\hline
\multicolumn{3}{|p{12cm}|}{Nome e Sigla do Departamento} & Unidade Acadêmica: \\ 
\multicolumn{3}{|p{12cm}|}{Departamento de Computação (DECOM)} & {ICEB} \\ 
\hline
\multicolumn{4}{|p{16cm}|}{Modalidade de Oferta:
[X] presencial \hspace{0.6cm}
[ ] semipresencial \hspace{0.6cm}
[ ] à distância}\\
\hline
\multicolumn{2}{|p{8cm}|}{Carga horária semestral} &
\multicolumn{2}{p{8cm}|}{Carga horária semanal}\\
\hline
\multicolumn{1}{|p{4cm}|}{Total} &
\multicolumn{1}{p{4cm}|}{Extensionista} &
\multicolumn{1}{p{4cm}|}{Teórica} &
\multicolumn{1}{p{4cm}|}{Prática} \\ 
\multicolumn{1}{|p{4cm}|}{90\,horas} &
\multicolumn{1}{p{4cm}|}{0\;horas} &
\multicolumn{1}{p{4cm}|}{4\;horas/aula} &
\multicolumn{1}{p{4cm}|}{2\;horas/aula} \\ 
\hline
\multicolumn{4}{|p{16cm}|}{Ementa:}\\
\multicolumn{4}{|p{16cm}|}{}\\
\multicolumn{4}{|p{\dimexpr 16cm + 6\tabcolsep\relax}|}{Recursividade; conceitos básicos de análise assintótica de algoritmos; tipos abstratos de dados; estruturas de dados: listas, pilhas, filas de prioridade e árvores binárias; algoritmos de ordenação por comparação de chaves: seleção, inserção, bolha, shellsort, quicksort, mergesort, heapsort; algoritmos de ordenação em tempo linear: counting sort, radix sort e bucket sort; e algoritmos de pesquisa: simples, binária, árvores binárias de busca, hashing, conjuntos e mapas.}\\
\multicolumn{4}{|p{16cm}|}{}\\
\hline
\multicolumn{4}{|p{16cm}|}{Conteúdo programático:}\\*
\multicolumn{4}{|p{\dimexpr 16cm + 6\tabcolsep\relax}|}{\begin{enumerate}[start=1]\item Revisão de alocação dinâmica de memória\end{enumerate}}\\
\multicolumn{4}{|p{\dimexpr 16cm + 6\tabcolsep\relax}|}{\begin{enumerate}[start=2]\item Recursividade\end{enumerate}}\\
\multicolumn{4}{|p{\dimexpr 16cm + 6\tabcolsep\relax}|}{\begin{enumerate}[start=3]\item Noções de análise de complexidade de algoritmos: conceitos; medidas de avaliação (tempo e espaço); análise assintótica (notação O, Omega e Theta); hierarquia de funções e classes de problemas\end{enumerate}}\\
\multicolumn{4}{|p{\dimexpr 16cm + 6\tabcolsep\relax}|}{\begin{enumerate}[start=4]\item Tipos de dados abstratos\end{enumerate}}\\
\multicolumn{4}{|p{\dimexpr 16cm + 6\tabcolsep\relax}|}{\begin{enumerate}[start=5]\item Estruturas de dados: listas; pilhas; filas; filas de prioridade; árvores; conjuntos; mapas\end{enumerate}}\\
\multicolumn{4}{|p{\dimexpr 16cm + 6\tabcolsep\relax}|}{\begin{enumerate}[start=6]\item Algoritmos: métodos de ordenação por comparação (Selection Sort, Insertion Sort, Bubblesort, Shellsort, Quicksort, Heapsort e Mergesort); métodos de ordenação em tempo linear (Counting Sort, Radix Sort e Bucket Sort); métodos de pesquisa (simples, binária, árvores binárias e AVL e hashing)\end{enumerate}}\\
\multicolumn{4}{|p{\dimexpr 16cm + 6\tabcolsep\relax}|}{}\\
\hline
\multicolumn{4}{|p{16cm}|}{Bibliografia Básica:}\\*
\multicolumn{4}{|p{\dimexpr 16cm + 6\tabcolsep\relax}|}{\begin{itemize}\item ZIVIANI, N. Projeto de algoritmos: com implementações em Pascal e C. 3. ed. rev. e ampl. São Paulo: Cengage Learning, c2011.\end{itemize}}\\
\multicolumn{4}{|p{\dimexpr 16cm + 6\tabcolsep\relax}|}{\begin{itemize}\item CELES, W.; CERQUEIRA, R.; RANGEL, J. L. Introdução a estruturas de dados: com técnicas de programação em C. Rio de Janeiro: Elsevier, 2004.\end{itemize}}\\
\multicolumn{4}{|p{\dimexpr 16cm + 6\tabcolsep\relax}|}{\begin{itemize}\item CORMEN, T. H. Algoritmos: teoria e prática. Rio de Janeiro: Campus, 2002.\end{itemize}}\\
\multicolumn{4}{|p{16cm}|}{}\\
\hline
\multicolumn{4}{|p{16cm}|}{Bibliografia Complementar:}\\*
\multicolumn{4}{|p{\dimexpr 16cm + 6\tabcolsep\relax}|}{\begin{itemize}\item KLEINBERG, J.; TARDOS, E. Algorithm design. Boston: Pearson/Addison-Wesley, c2006.\end{itemize}}\\
\multicolumn{4}{|p{\dimexpr 16cm + 6\tabcolsep\relax}|}{\begin{itemize}\item KNUTH, D. E. The art of computer programming. Upper Saddle River: Addison Wesley, c2005.\end{itemize}}\\
\multicolumn{4}{|p{\dimexpr 16cm + 6\tabcolsep\relax}|}{\begin{itemize}\item GOODRICH, M. T.; TAMASSIA, R.; COPSTEIN, B. Projeto de algoritmos: fundamentos, análise e exemplos da internet. Porto Alegre: Bookman, 2004.\end{itemize}}\\
\multicolumn{4}{|p{\dimexpr 16cm + 6\tabcolsep\relax}|}{\begin{itemize}\item DROZDEK, A. Estrutura de dados e algoritmos em C++. São Paulo: Cengage Learning, 2002.\end{itemize}}\\
\multicolumn{4}{|p{\dimexpr 16cm + 6\tabcolsep\relax}|}{\begin{itemize}\item TENENBAUM, A. M.; LANGSAM, Y.; AUGENSTEIN, M. Estruturas de dados usando C. São Paulo: Makron Books, 1995.\end{itemize}}\\
\hline
\end{longtable}
\end{center}

\clearpage
\begin{center}
\begin{tabular}{lccr}
 \multirow{3}{*}{\includegraphics[height=2.7cm]{brasao.png}} &
 \multicolumn{2}{c}{\bfseries UNIVERSIDADE FEDERAL DE OURO PRETO} &
 \ \ \ \ \multirow{3}{*}{\includegraphics[height=2.7cm]{ufop.png}} \\
 & \multicolumn{2}{c}{\bfseries PRÓ-REITORIA DE GRADUAÇÃO} & \\
 & \multicolumn{2}{c}{} & \\
 & \multicolumn{2}{c}{\Large\bfseries PROGRAMA DE DISCIPLINA} & \\
\end{tabular}
\end{center}

\begin{center}
\begin{longtable}{|p{4cm}|p{4cm}|p{4cm}|p{4cm}|}
\hline
\multicolumn{3}{|p{12cm}|}{Nome do Componente Curricular em Português:} &
\multicolumn{1}{p{4cm}|}{Código:} \\ 
\multicolumn{3}{|p{12cm}|}{\textbf{Estatística e Probabilidade}} &
\textbf{EST202}\\ 
\multicolumn{3}{|p{12cm}|}{Nome do Componente Curricular em Inglês:} & \\ 
\multicolumn{3}{|p{12cm}|}{\textbf{Statistics and Probability}} & \\ 
\hline
\multicolumn{3}{|p{12cm}|}{Nome e Sigla do Departamento} & Unidade Acadêmica: \\ 
\multicolumn{3}{|p{12cm}|}{Departamento de Estatística (DEEST)} & {ICEB} \\ 
\hline
\multicolumn{4}{|p{16cm}|}{Modalidade de Oferta:
[X] presencial \hspace{0.6cm}
[ ] semipresencial \hspace{0.6cm}
[ ] à distância}\\
\hline
\multicolumn{2}{|p{8cm}|}{Carga horária semestral} &
\multicolumn{2}{p{8cm}|}{Carga horária semanal}\\
\hline
\multicolumn{1}{|p{4cm}|}{Total} &
\multicolumn{1}{p{4cm}|}{Extensionista} &
\multicolumn{1}{p{4cm}|}{Teórica} &
\multicolumn{1}{p{4cm}|}{Prática} \\ 
\multicolumn{1}{|p{4cm}|}{60\,horas} &
\multicolumn{1}{p{4cm}|}{0\;horas} &
\multicolumn{1}{p{4cm}|}{4\;horas/aula} &
\multicolumn{1}{p{4cm}|}{0\;horas/aula} \\ 
\hline
\multicolumn{4}{|p{16cm}|}{Ementa:}\\
\multicolumn{4}{|p{16cm}|}{}\\
\multicolumn{4}{|p{\dimexpr 16cm + 6\tabcolsep\relax}|}{Introdução. Técnicas de amostragem. Estatística descritiva. Introdução à probabilidade. Variáveis aleatórias unidimensionais. Modelos de distribuição de probabilidade. Inferência. Regressão linear simples.}\\
\multicolumn{4}{|p{16cm}|}{}\\
\hline
\multicolumn{4}{|p{16cm}|}{Conteúdo programático:}\\*
\multicolumn{4}{|p{\dimexpr 16cm + 6\tabcolsep\relax}|}{\begin{enumerate}[start=1]\item Introdução\end{enumerate}}\\
\multicolumn{4}{|p{\dimexpr 16cm + 6\tabcolsep\relax}|}{\begin{enumerate}[start=2]\item Técnicas de amostragem aleatória: simples; estratificada proporcional; estratificada de igual tamanho; sistemática; por conglomerado\end{enumerate}}\\
\multicolumn{4}{|p{\dimexpr 16cm + 6\tabcolsep\relax}|}{\begin{enumerate}[start=3]\item Estatística descritiva: distribuição de frequência e gráficos; medidas de posição; medidas de dispersão\end{enumerate}}\\
\multicolumn{4}{|p{\dimexpr 16cm + 6\tabcolsep\relax}|}{\begin{enumerate}[start=4]\item Introdução à probabilidade: revisão da teoria dos conjuntos; experimento aleatório, espaço amostral e evento; definição de probabilidade; eventos condicionados; evento interseção e união; independência\end{enumerate}}\\
\multicolumn{4}{|p{\dimexpr 16cm + 6\tabcolsep\relax}|}{\begin{enumerate}[start=5]\item Variáveis aleatórias unidimensionais: função geratriz de variáveis aleatórias; função de probabilidade de variáveis aleatórias discretas e contínuas; média e variância\end{enumerate}}\\
\multicolumn{4}{|p{\dimexpr 16cm + 6\tabcolsep\relax}|}{\begin{enumerate}[start=6]\item Modelos de distribuição discretas: binomial; Poisson; hipergeométrica\end{enumerate}}\\
\multicolumn{4}{|p{\dimexpr 16cm + 6\tabcolsep\relax}|}{\begin{enumerate}[start=7]\item Distribuição normal: definição; propriedades; construção e uso da tabela; combinação linear de variáveis aleatórias normais; distribuição da média amostral\end{enumerate}}\\
\multicolumn{4}{|p{\dimexpr 16cm + 6\tabcolsep\relax}|}{\begin{enumerate}[start=8]\item Inferência: estimação de parâmetros populacionais; intervalo de confiança para média, para proporção e para variância; testes de hipótese para a média\end{enumerate}}\\
\multicolumn{4}{|p{\dimexpr 16cm + 6\tabcolsep\relax}|}{\begin{enumerate}[start=9]\item Regressão linear simples\end{enumerate}}\\
\multicolumn{4}{|p{\dimexpr 16cm + 6\tabcolsep\relax}|}{}\\
\hline
\multicolumn{4}{|p{16cm}|}{Bibliografia Básica:}\\*
\multicolumn{4}{|p{\dimexpr 16cm + 6\tabcolsep\relax}|}{\begin{itemize}\item Mirshawka, Victor. Probabilidade e estatística para a engenharia. v. 1. São Paulo: Livraria Nobel, 1983.\end{itemize}}\\
\multicolumn{4}{|p{\dimexpr 16cm + 6\tabcolsep\relax}|}{\begin{itemize}\item Wonnacott, Thomas H. Introdução à estatística. Rio de Janeiro: Livros Técnicos e Científicos.\end{itemize}}\\
\multicolumn{4}{|p{\dimexpr 16cm + 6\tabcolsep\relax}|}{\begin{itemize}\item Mirshawka, Victor. Probabilidades e estatística para engenharia. São Paulo: Nobel, 1980.\end{itemize}}\\
\multicolumn{4}{|p{\dimexpr 16cm + 6\tabcolsep\relax}|}{\begin{itemize}\item Soares, José Francisco. Introdução à estatística. Rio de Janeiro: Guanabara Koogan, 1991.\end{itemize}}\\
\multicolumn{4}{|p{16cm}|}{}\\
\hline
\multicolumn{4}{|p{16cm}|}{Bibliografia Complementar:}\\*
\multicolumn{4}{|p{\dimexpr 16cm + 6\tabcolsep\relax}|}{\begin{itemize}\item Paiva, Antônio Fabiano. Estatística. v. 1. Belo Horizonte: Imprensa da UFMG, 1981.\end{itemize}}\\
\multicolumn{4}{|p{\dimexpr 16cm + 6\tabcolsep\relax}|}{\begin{itemize}\item Paiva, Antônio Fabiano. Estatística. v. 2. Belo Horizonte: Imprensa da UFMG, 1981.\end{itemize}}\\
\multicolumn{4}{|p{\dimexpr 16cm + 6\tabcolsep\relax}|}{\begin{itemize}\item Meyer, Paul L. Probabilidade: aplicações à estatística. 2. ed. Rio de Janeiro: Livros Técnicos e Científicos, 1983.\end{itemize}}\\
\multicolumn{4}{|p{\dimexpr 16cm + 6\tabcolsep\relax}|}{\begin{itemize}\item Costa Neto, Pedro Luiz de Oliveira. Estatística. São Paulo: E. Blucher, 1977.\end{itemize}}\\
\hline
\end{longtable}
\end{center}

\clearpage
\begin{center}
\begin{tabular}{lccr}
 \multirow{3}{*}{\includegraphics[height=2.7cm]{brasao.png}} &
 \multicolumn{2}{c}{\bfseries UNIVERSIDADE FEDERAL DE OURO PRETO} &
 \ \ \ \ \multirow{3}{*}{\includegraphics[height=2.7cm]{ufop.png}} \\
 & \multicolumn{2}{c}{\bfseries PRÓ-REITORIA DE GRADUAÇÃO} & \\
 & \multicolumn{2}{c}{} & \\
 & \multicolumn{2}{c}{\Large\bfseries PROGRAMA DE DISCIPLINA} & \\
\end{tabular}
\end{center}

\begin{center}
\begin{longtable}{|p{4cm}|p{4cm}|p{4cm}|p{4cm}|}
\hline
\multicolumn{3}{|p{12cm}|}{Nome do Componente Curricular em Português:} &
\multicolumn{1}{p{4cm}|}{Código:} \\ 
\multicolumn{3}{|p{12cm}|}{\textbf{Introdução à Ciência de Dados}} &
\textbf{BIA104}\\ 
\multicolumn{3}{|p{12cm}|}{Nome do Componente Curricular em Inglês:} & \\ 
\multicolumn{3}{|p{12cm}|}{\textbf{Introduction to Data Science}} & \\ 
\hline
\multicolumn{3}{|p{12cm}|}{Nome e Sigla do Departamento} & Unidade Acadêmica: \\ 
\multicolumn{3}{|p{12cm}|}{Departamento de Computação (DECOM)} & {ICEB} \\ 
\hline
\multicolumn{4}{|p{16cm}|}{Modalidade de Oferta:
[X] presencial \hspace{0.6cm}
[ ] semipresencial \hspace{0.6cm}
[ ] à distância}\\
\hline
\multicolumn{2}{|p{8cm}|}{Carga horária semestral} &
\multicolumn{2}{p{8cm}|}{Carga horária semanal}\\
\hline
\multicolumn{1}{|p{4cm}|}{Total} &
\multicolumn{1}{p{4cm}|}{Extensionista} &
\multicolumn{1}{p{4cm}|}{Teórica} &
\multicolumn{1}{p{4cm}|}{Prática} \\ 
\multicolumn{1}{|p{4cm}|}{60\,horas} &
\multicolumn{1}{p{4cm}|}{0\;horas} &
\multicolumn{1}{p{4cm}|}{4\;horas/aula} &
\multicolumn{1}{p{4cm}|}{0\;horas/aula} \\ 
\hline
\multicolumn{4}{|p{16cm}|}{Ementa:}\\
\multicolumn{4}{|p{16cm}|}{}\\
\multicolumn{4}{|p{\dimexpr 16cm + 6\tabcolsep\relax}|}{Introdução à Ciência de Dados. Pré-processamento de dados. Análise exploratória de dados. Processamento de dados em larga escala. Workflow e rastreamento de experimentos. Ciência de Dados Responsável. Tópicos avançados em ciência de dados}\\
\multicolumn{4}{|p{16cm}|}{}\\
\hline
\multicolumn{4}{|p{16cm}|}{Conteúdo programático:}\\*
\multicolumn{4}{|p{\dimexpr 16cm + 6\tabcolsep\relax}|}{\begin{enumerate}[start=1]\item Introdução à Ciência de Dados. \end{enumerate}}\\
\multicolumn{4}{|p{\dimexpr 16cm + 6\tabcolsep\relax}|}{\begin{enumerate}[start=2]\item Tipos de dados.\end{enumerate}}\\
\multicolumn{4}{|p{\dimexpr 16cm + 6\tabcolsep\relax}|}{\begin{enumerate}[start=3]\item Aquisição de dados.\end{enumerate}}\\
\multicolumn{4}{|p{\dimexpr 16cm + 6\tabcolsep\relax}|}{\begin{enumerate}[start=4]\item Preparação de dados. \end{enumerate}}\\
\multicolumn{4}{|p{\dimexpr 16cm + 6\tabcolsep\relax}|}{\begin{enumerate}[start=5]\item Análise exploratória de dados: visualização, estatísticas descritivas e pré-processamento.\end{enumerate}}\\
\multicolumn{4}{|p{\dimexpr 16cm + 6\tabcolsep\relax}|}{\begin{enumerate}[start=6]\item Processamento de dados em larga escala.\end{enumerate}}\\
\multicolumn{4}{|p{\dimexpr 16cm + 6\tabcolsep\relax}|}{\begin{enumerate}[start=7]\item Workflow e rastreamento de experimentos.\end{enumerate}}\\
\multicolumn{4}{|p{\dimexpr 16cm + 6\tabcolsep\relax}|}{\begin{enumerate}[start=8]\item Tópicos avançados em Ciência de Dados.\end{enumerate}}\\
\multicolumn{4}{|p{\dimexpr 16cm + 6\tabcolsep\relax}|}{\begin{enumerate}[start=9]\item Ciência de Dados Ética e Responsável.\end{enumerate}}\\
\multicolumn{4}{|p{\dimexpr 16cm + 6\tabcolsep\relax}|}{}\\
\hline
\multicolumn{4}{|p{16cm}|}{Bibliografia Básica:}\\*
\multicolumn{4}{|p{\dimexpr 16cm + 6\tabcolsep\relax}|}{\begin{itemize}\item CARVALHO, Andre Carlos Ponce de Leon Ferreira de; MENEZES, Angelo Garangau; BONIDIA, Robson Parmezan. Ciência de dados: fundamentos e aplicações. 2024.\end{itemize}}\\
\multicolumn{4}{|p{\dimexpr 16cm + 6\tabcolsep\relax}|}{\begin{itemize}\item TAN, Pang-Ning., Steinbach, Michael., Kumar, Vipin. Introduction to Data Mining. Boston: Pearson Addison Wesley, 2006.\end{itemize}}\\
\multicolumn{4}{|p{\dimexpr 16cm + 6\tabcolsep\relax}|}{\begin{itemize}\item HASTIE, Trevor; TIBSHIRANI, Robert; FRIEDMAN, Jerome. The Elements of Statistical Learning. 2. ed. New York: Springer, 2009.\end{itemize}}\\
\multicolumn{4}{|p{\dimexpr 16cm + 6\tabcolsep\relax}|}{\begin{itemize}\item MURPHY, Kevin P. Probabilistic Machine Learning: An Introduction. Cambridge: MIT Press, 2022.\end{itemize}}\\
\multicolumn{4}{|p{\dimexpr 16cm + 6\tabcolsep\relax}|}{\begin{itemize}\item JAMES, Gareth et al. An Introduction to Statistical Learning. 2. ed. New York: Springer, 2021.\end{itemize}}\\
\multicolumn{4}{|p{16cm}|}{}\\
\hline
\multicolumn{4}{|p{16cm}|}{Bibliografia Complementar:}\\*
\multicolumn{4}{|p{\dimexpr 16cm + 6\tabcolsep\relax}|}{\begin{itemize}\item BRUCE, Peter; BRUCE, Andrew; GEDECK, Peter. Practical statistics for data scientists: 50+ essential concepts using R and Python. O'Reilly Media, 2020.\end{itemize}}\\
\multicolumn{4}{|p{\dimexpr 16cm + 6\tabcolsep\relax}|}{\begin{itemize}\item MÜLLER, Andreas C.; GUIDO, Sarah. Introduction to Machine Learning with Python. 1. ed. O’Reilly, 2016.\end{itemize}}\\
\multicolumn{4}{|p{\dimexpr 16cm + 6\tabcolsep\relax}|}{\begin{itemize}\item BISHOP, Christopher M. Pattern Recognition and Machine Learning. New York: Springer, 2006.\end{itemize}}\\
\multicolumn{4}{|p{\dimexpr 16cm + 6\tabcolsep\relax}|}{\begin{itemize}\item PROVOST, Foster; FAWCETT, Tom. Data science for business: what you need to know about data mining and data-analytic thinking. Sebastopol, CA: O'Reilly Media, 2013.\end{itemize}}\\
\multicolumn{4}{|p{\dimexpr 16cm + 6\tabcolsep\relax}|}{\begin{itemize}\item OLIVEIRA, Hélio C. V. de. Mineração de dados: conceitos, algoritmos, aplicações. 2. ed. Rio de Janeiro: LTC, 2018. E-book. Disponível em: https://plataforma.bvirtual.com.br. Acesso em: jun/2025.\end{itemize}}\\
\hline
\end{longtable}
\end{center}

\clearpage
\begin{center}
\begin{tabular}{lccr}
 \multirow{3}{*}{\includegraphics[height=2.7cm]{brasao.png}} &
 \multicolumn{2}{c}{\bfseries UNIVERSIDADE FEDERAL DE OURO PRETO} &
 \ \ \ \ \multirow{3}{*}{\includegraphics[height=2.7cm]{ufop.png}} \\
 & \multicolumn{2}{c}{\bfseries PRÓ-REITORIA DE GRADUAÇÃO} & \\
 & \multicolumn{2}{c}{} & \\
 & \multicolumn{2}{c}{\Large\bfseries PROGRAMA DE DISCIPLINA} & \\
\end{tabular}
\end{center}

\begin{center}
\begin{longtable}{|p{4cm}|p{4cm}|p{4cm}|p{4cm}|}
\hline
\multicolumn{3}{|p{12cm}|}{Nome do Componente Curricular em Português:} &
\multicolumn{1}{p{4cm}|}{Código:} \\ 
\multicolumn{3}{|p{12cm}|}{\textbf{Aprendizado de Máquina Supervisionado}} &
\textbf{BIA105}\\ 
\multicolumn{3}{|p{12cm}|}{Nome do Componente Curricular em Inglês:} & \\ 
\multicolumn{3}{|p{12cm}|}{\textbf{Supervised Machine Learning}} & \\ 
\hline
\multicolumn{3}{|p{12cm}|}{Nome e Sigla do Departamento} & Unidade Acadêmica: \\ 
\multicolumn{3}{|p{12cm}|}{Departamento de Computação (DECOM)} & {ICEB} \\ 
\hline
\multicolumn{4}{|p{16cm}|}{Modalidade de Oferta:
[X] presencial \hspace{0.6cm}
[ ] semipresencial \hspace{0.6cm}
[ ] à distância}\\
\hline
\multicolumn{2}{|p{8cm}|}{Carga horária semestral} &
\multicolumn{2}{p{8cm}|}{Carga horária semanal}\\
\hline
\multicolumn{1}{|p{4cm}|}{Total} &
\multicolumn{1}{p{4cm}|}{Extensionista} &
\multicolumn{1}{p{4cm}|}{Teórica} &
\multicolumn{1}{p{4cm}|}{Prática} \\ 
\multicolumn{1}{|p{4cm}|}{60\,horas} &
\multicolumn{1}{p{4cm}|}{0\;horas} &
\multicolumn{1}{p{4cm}|}{4\;horas/aula} &
\multicolumn{1}{p{4cm}|}{0\;horas/aula} \\ 
\hline
\multicolumn{4}{|p{16cm}|}{Ementa:}\\
\multicolumn{4}{|p{16cm}|}{}\\
\multicolumn{4}{|p{\dimexpr 16cm + 6\tabcolsep\relax}|}{Conceitos e algoritmos de aprendizado de máquina supervisionado. Classificação, regressão, avaliação de modelos e análise de desempenho. Aplicações em ciência de dados e inteligência artificial. Enfoque prático com experimentação computacional.}\\
\multicolumn{4}{|p{16cm}|}{}\\
\hline
\multicolumn{4}{|p{16cm}|}{Conteúdo programático:}\\*
\multicolumn{4}{|p{\dimexpr 16cm + 6\tabcolsep\relax}|}{\begin{enumerate}[start=1]\item Introdução ao aprendizado de máquina: conceitos, tarefas e tipos de aprendizado.\end{enumerate}}\\
\multicolumn{4}{|p{\dimexpr 16cm + 6\tabcolsep\relax}|}{\begin{enumerate}[start=2]\item Conjuntos de dados supervisionados: características, rótulos e pré-processamento.\end{enumerate}}\\
\multicolumn{4}{|p{\dimexpr 16cm + 6\tabcolsep\relax}|}{\begin{enumerate}[start=3]\item Conceito de função alvo, hipótese e generalização.\end{enumerate}}\\
\multicolumn{4}{|p{\dimexpr 16cm + 6\tabcolsep\relax}|}{\begin{enumerate}[start=4]\item Modelos de classificação: k-NN, Naive Bayes, árvores de decisão, regressão logística, SVM.\end{enumerate}}\\
\multicolumn{4}{|p{\dimexpr 16cm + 6\tabcolsep\relax}|}{\begin{enumerate}[start=5]\item Modelos de regressão: regressão linear simples e múltipla, regularização (Ridge, Lasso).\end{enumerate}}\\
\multicolumn{4}{|p{\dimexpr 16cm + 6\tabcolsep\relax}|}{\begin{enumerate}[start=6]\item Técnicas de ensemble: Random Forest, Gradient Boosting, bagging e boosting.\end{enumerate}}\\
\multicolumn{4}{|p{\dimexpr 16cm + 6\tabcolsep\relax}|}{\begin{enumerate}[start=7]\item Divisão de dados: treino, validação e teste; validação cruzada.\end{enumerate}}\\
\multicolumn{4}{|p{\dimexpr 16cm + 6\tabcolsep\relax}|}{\begin{enumerate}[start=8]\item Métricas de avaliação: acurácia, precisão, recall, F1, curva ROC e AUC.\end{enumerate}}\\
\multicolumn{4}{|p{\dimexpr 16cm + 6\tabcolsep\relax}|}{\begin{enumerate}[start=9]\item Overfitting e underfitting: diagnóstico e controle (bias-variance tradeoff).\end{enumerate}}\\
\multicolumn{4}{|p{\dimexpr 16cm + 6\tabcolsep\relax}|}{\begin{enumerate}[start=10]\item Uso de bibliotecas de software para modelagem supervisionada.\end{enumerate}}\\
\multicolumn{4}{|p{\dimexpr 16cm + 6\tabcolsep\relax}|}{\begin{enumerate}[start=11]\item Impactos Sociais, Viés e Limitações do Aprendizado de Máquina\end{enumerate}}\\
\multicolumn{4}{|p{\dimexpr 16cm + 6\tabcolsep\relax}|}{}\\
\hline
\multicolumn{4}{|p{16cm}|}{Bibliografia Básica:}\\*
\multicolumn{4}{|p{\dimexpr 16cm + 6\tabcolsep\relax}|}{\begin{itemize}\item JAMES, Gareth et al. An introduction to statistical learning. 2. ed. New York: Springer, 2021. Disponível em: <https://www.statlearning.com> . Acesso em: 22 dez. 2025.\end{itemize}}\\
\multicolumn{4}{|p{\dimexpr 16cm + 6\tabcolsep\relax}|}{\begin{itemize}\item GERON, Aurélien. Mãos à obra: aprendizado de máquina com Scikit-Learn, Keras e TensorFlow. 2. ed. São Paulo: Novatec, 2020. E-book. Disponível em: https://plataforma.bvirtual.com.br. Acesso em: jun/2025.\end{itemize}}\\
\multicolumn{4}{|p{\dimexpr 16cm + 6\tabcolsep\relax}|}{\begin{itemize}\item LIU, Yuxi Hayden. Machine learning with Python in practice: unlock machine learning best practices with real-world use cases. 1. ed. São Paulo: Blucher, 2025. E-book. Disponível em: <https://plataforma.bvirtual.com.br> . Acesso em: 22 dez. 2025.\end{itemize}}\\
\multicolumn{4}{|p{16cm}|}{}\\
\hline
\multicolumn{4}{|p{16cm}|}{Bibliografia Complementar:}\\*
\multicolumn{4}{|p{\dimexpr 16cm + 6\tabcolsep\relax}|}{\begin{itemize}\item GOODFELLOW, Ian; BENGIO, Yoshua; COURVILLE, Aaron. Deep Learning. Cambridge: MIT Press, 2016. (Capítulos introdutórios sobre aprendizado supervisionado).\end{itemize}}\\
\multicolumn{4}{|p{\dimexpr 16cm + 6\tabcolsep\relax}|}{\begin{itemize}\item MÜLLER, Andreas C.; GUIDO, Sarah. Introduction to Machine Learning with Python. 1. ed. O'Reilly, 2016.\end{itemize}}\\
\multicolumn{4}{|p{\dimexpr 16cm + 6\tabcolsep\relax}|}{\begin{itemize}\item BISHOP, Christopher M. Pattern Recognition and Machine Learning. New York: Springer, 2006.\end{itemize}}\\
\multicolumn{4}{|p{\dimexpr 16cm + 6\tabcolsep\relax}|}{\begin{itemize}\item RUSSELL, Stuart; NORVIG, Peter. Inteligência artificial. 3. ed. Rio de Janeiro: Elsevier, 2013. (Capítulo sobre aprendizado).\end{itemize}}\\
\multicolumn{4}{|p{\dimexpr 16cm + 6\tabcolsep\relax}|}{\begin{itemize}\item SHALEV-SHWARTZ, Shai; BEN-DAVID, Shai. Understanding machine learning: from theory to algorithms. Cambridge: Cambridge University Press, 2014.\end{itemize}}\\
\hline
\end{longtable}
\end{center}

\clearpage
\begin{center}
\begin{tabular}{lccr}
 \multirow{3}{*}{\includegraphics[height=2.7cm]{brasao.png}} &
 \multicolumn{2}{c}{\bfseries UNIVERSIDADE FEDERAL DE OURO PRETO} &
 \ \ \ \ \multirow{3}{*}{\includegraphics[height=2.7cm]{ufop.png}} \\
 & \multicolumn{2}{c}{\bfseries PRÓ-REITORIA DE GRADUAÇÃO} & \\
 & \multicolumn{2}{c}{} & \\
 & \multicolumn{2}{c}{\Large\bfseries PROGRAMA DE DISCIPLINA} & \\
\end{tabular}
\end{center}

\begin{center}
\begin{longtable}{|p{4cm}|p{4cm}|p{4cm}|p{4cm}|}
\hline
\multicolumn{3}{|p{12cm}|}{Nome do Componente Curricular em Português:} &
\multicolumn{1}{p{4cm}|}{Código:} \\ 
\multicolumn{3}{|p{12cm}|}{\textbf{Aprendizado de Máquina Não Supervisionado}} &
\textbf{BIA106}\\ 
\multicolumn{3}{|p{12cm}|}{Nome do Componente Curricular em Inglês:} & \\ 
\multicolumn{3}{|p{12cm}|}{\textbf{Unsupervised Machine Learning}} & \\ 
\hline
\multicolumn{3}{|p{12cm}|}{Nome e Sigla do Departamento} & Unidade Acadêmica: \\ 
\multicolumn{3}{|p{12cm}|}{Departamento de Computação (DECOM)} & {ICEB} \\ 
\hline
\multicolumn{4}{|p{16cm}|}{Modalidade de Oferta:
[X] presencial \hspace{0.6cm}
[ ] semipresencial \hspace{0.6cm}
[ ] à distância}\\
\hline
\multicolumn{2}{|p{8cm}|}{Carga horária semestral} &
\multicolumn{2}{p{8cm}|}{Carga horária semanal}\\
\hline
\multicolumn{1}{|p{4cm}|}{Total} &
\multicolumn{1}{p{4cm}|}{Extensionista} &
\multicolumn{1}{p{4cm}|}{Teórica} &
\multicolumn{1}{p{4cm}|}{Prática} \\ 
\multicolumn{1}{|p{4cm}|}{60\,horas} &
\multicolumn{1}{p{4cm}|}{0\;horas} &
\multicolumn{1}{p{4cm}|}{4\;horas/aula} &
\multicolumn{1}{p{4cm}|}{0\;horas/aula} \\ 
\hline
\multicolumn{4}{|p{16cm}|}{Ementa:}\\
\multicolumn{4}{|p{16cm}|}{}\\
\multicolumn{4}{|p{\dimexpr 16cm + 6\tabcolsep\relax}|}{Estudo dos principais métodos de aprendizado de máquina não supervisionado. Técnicas de agrupamento, redução de dimensionalidade, extração de características e detecção de anomalias. Aplicações em ciência de dados, visualização e descoberta de padrões.}\\
\multicolumn{4}{|p{16cm}|}{}\\
\hline
\multicolumn{4}{|p{16cm}|}{Conteúdo programático:}\\*
\multicolumn{4}{|p{\dimexpr 16cm + 6\tabcolsep\relax}|}{\begin{enumerate}[start=1]\item Introdução ao aprendizado não supervisionado: definição, desafios e aplicações.\end{enumerate}}\\
\multicolumn{4}{|p{\dimexpr 16cm + 6\tabcolsep\relax}|}{\begin{enumerate}[start=2]\item Análise exploratória de dados: visualização, estatísticas descritivas e pré-processamento.\end{enumerate}}\\
\multicolumn{4}{|p{\dimexpr 16cm + 6\tabcolsep\relax}|}{\begin{enumerate}[start=3]\item Técnicas de agrupamento: k-means, k-medoids, DBSCAN, aglomerativo hierárquico.\end{enumerate}}\\
\multicolumn{4}{|p{\dimexpr 16cm + 6\tabcolsep\relax}|}{\begin{enumerate}[start=4]\item Métricas para avaliação de agrupamentos: silhueta, SSE, Davies-Bouldin.\end{enumerate}}\\
\multicolumn{4}{|p{\dimexpr 16cm + 6\tabcolsep\relax}|}{\begin{enumerate}[start=5]\item Modelos baseados em mistura: Gaussian Mixture Models (GMM).\end{enumerate}}\\
\multicolumn{4}{|p{\dimexpr 16cm + 6\tabcolsep\relax}|}{\begin{enumerate}[start=6]\item Redução de dimensionalidade: PCA (Análise de Componentes Principais), t-SNE, UMAP.\end{enumerate}}\\
\multicolumn{4}{|p{\dimexpr 16cm + 6\tabcolsep\relax}|}{\begin{enumerate}[start=7]\item Extração de características e embeddings para dados complexos.\end{enumerate}}\\
\multicolumn{4}{|p{\dimexpr 16cm + 6\tabcolsep\relax}|}{\begin{enumerate}[start=8]\item Detecção de anomalias: outliers, isolamento de florestas, técnicas baseadas em distância.\end{enumerate}}\\
\multicolumn{4}{|p{\dimexpr 16cm + 6\tabcolsep\relax}|}{\begin{enumerate}[start=9]\item Tópicos avançados: agrupamento em dados de alta dimensão e em grandes volumes.\end{enumerate}}\\
\multicolumn{4}{|p{\dimexpr 16cm + 6\tabcolsep\relax}|}{\begin{enumerate}[start=10]\item Implementações práticas com bibliotecas de software.\end{enumerate}}\\
\multicolumn{4}{|p{\dimexpr 16cm + 6\tabcolsep\relax}|}{\begin{enumerate}[start=11]\item Impactos Sociais, Viés e Limitações do Aprendizado de Máquina\end{enumerate}}\\
\multicolumn{4}{|p{\dimexpr 16cm + 6\tabcolsep\relax}|}{}\\
\hline
\multicolumn{4}{|p{16cm}|}{Bibliografia Básica:}\\*
\multicolumn{4}{|p{\dimexpr 16cm + 6\tabcolsep\relax}|}{\begin{itemize}\item GERON, Aurélien. Mãos à obra: aprendizado de máquina com Scikit-Learn, Keras e TensorFlow. 2. ed. São Paulo: Novatec, 2020. E-book. Disponível em: https://plataforma.bvirtual.com.br. Acesso em: jun/2025.\end{itemize}}\\
\multicolumn{4}{|p{\dimexpr 16cm + 6\tabcolsep\relax}|}{\begin{itemize}\item MURPHY, Kevin P. Probabilistic Machine Learning: An Introduction. Cambridge: MIT Press, 2022.\end{itemize}}\\
\multicolumn{4}{|p{\dimexpr 16cm + 6\tabcolsep\relax}|}{\begin{itemize}\item JAMES, Gareth et al. An introduction to statistical learning. 2. ed. New York: Springer, 2021. Disponível em: <https://www.statlearning.com> . Acesso em: 22 dez. 2025.\end{itemize}}\\
\multicolumn{4}{|p{\dimexpr 16cm + 6\tabcolsep\relax}|}{\begin{itemize}\item LENZ, Maikon L.; NEUMANN, Fabiano B.; SANTARELLI, Rodrigo; et al. Fundamentos de aprendizagem de máquina. Porto Alegre: Grupo A, 2020. E-book. Disponível em: <https://integrada.minhabiblioteca.com.br/reader/books/9786556900902>. Acesso em: 22 dez. 2025.\end{itemize}}\\
\multicolumn{4}{|p{16cm}|}{}\\
\hline
\multicolumn{4}{|p{16cm}|}{Bibliografia Complementar:}\\*
\multicolumn{4}{|p{\dimexpr 16cm + 6\tabcolsep\relax}|}{\begin{itemize}\item HASTIE, Trevor; TIBSHIRANI, Robert; FRIEDMAN, Jerome. The Elements of Statistical Learning. 2. ed. New York: Springer, 2009.\end{itemize}}\\
\multicolumn{4}{|p{\dimexpr 16cm + 6\tabcolsep\relax}|}{\begin{itemize}\item MÜLLER, Andreas C.; GUIDO, Sarah. Introduction to Machine Learning with Python. 1. ed. O’Reilly, 2016.\end{itemize}}\\
\multicolumn{4}{|p{\dimexpr 16cm + 6\tabcolsep\relax}|}{\begin{itemize}\item BISHOP, Christopher M. Pattern Recognition and Machine Learning. New York: Springer, 2006.\end{itemize}}\\
\multicolumn{4}{|p{\dimexpr 16cm + 6\tabcolsep\relax}|}{\begin{itemize}\item AGGARWAL, Charu C.; REDDY, Chandan K. (ed.). Data clustering: algorithms and applications. Boca Raton: Chapman \& Hall/CRC, 2014.\end{itemize}}\\
\multicolumn{4}{|p{\dimexpr 16cm + 6\tabcolsep\relax}|}{\begin{itemize}\item OLIVEIRA, Hélio C. V. de. Mineração de dados: conceitos, algoritmos, aplicações. 2. ed. Rio de Janeiro: LTC, 2018. E-book. Disponível em: https://plataforma.bvirtual.com.br. Acesso em: jun/2025.\end{itemize}}\\
\hline
\end{longtable}
\end{center}

\clearpage
\begin{center}
\begin{tabular}{lccr}
 \multirow{3}{*}{\includegraphics[height=2.7cm]{brasao.png}} &
 \multicolumn{2}{c}{\bfseries UNIVERSIDADE FEDERAL DE OURO PRETO} &
 \ \ \ \ \multirow{3}{*}{\includegraphics[height=2.7cm]{ufop.png}} \\
 & \multicolumn{2}{c}{\bfseries PRÓ-REITORIA DE GRADUAÇÃO} & \\
 & \multicolumn{2}{c}{} & \\
 & \multicolumn{2}{c}{\Large\bfseries PROGRAMA DE DISCIPLINA} & \\
\end{tabular}
\end{center}

\begin{center}
\begin{longtable}{|p{4cm}|p{4cm}|p{4cm}|p{4cm}|}
\hline
\multicolumn{3}{|p{12cm}|}{Nome do Componente Curricular em Português:} &
\multicolumn{1}{p{4cm}|}{Código:} \\ 
\multicolumn{3}{|p{12cm}|}{\textbf{Programação Orientada a Objetos}} &
\textbf{BCC138}\\ 
\multicolumn{3}{|p{12cm}|}{Nome do Componente Curricular em Inglês:} & \\ 
\multicolumn{3}{|p{12cm}|}{\textbf{Object Oriented Programming}} & \\ 
\hline
\multicolumn{3}{|p{12cm}|}{Nome e Sigla do Departamento} & Unidade Acadêmica: \\ 
\multicolumn{3}{|p{12cm}|}{Departamento de Computação (DECOM)} & {ICEB} \\ 
\hline
\multicolumn{4}{|p{16cm}|}{Modalidade de Oferta:
[X] presencial \hspace{0.6cm}
[ ] semipresencial \hspace{0.6cm}
[ ] à distância}\\
\hline
\multicolumn{2}{|p{8cm}|}{Carga horária semestral} &
\multicolumn{2}{p{8cm}|}{Carga horária semanal}\\
\hline
\multicolumn{1}{|p{4cm}|}{Total} &
\multicolumn{1}{p{4cm}|}{Extensionista} &
\multicolumn{1}{p{4cm}|}{Teórica} &
\multicolumn{1}{p{4cm}|}{Prática} \\ 
\multicolumn{1}{|p{4cm}|}{60\,horas} &
\multicolumn{1}{p{4cm}|}{0\;horas} &
\multicolumn{1}{p{4cm}|}{4\;horas/aula} &
\multicolumn{1}{p{4cm}|}{0\;horas/aula} \\ 
\hline
\multicolumn{4}{|p{16cm}|}{Ementa:}\\
\multicolumn{4}{|p{16cm}|}{}\\
\multicolumn{4}{|p{\dimexpr 16cm + 6\tabcolsep\relax}|}{Conceitos básicos de orientação a objetos, Classe, Objeto, Mensagem, Encapsulamento, Herança, Polimorfismo, Ligação dinâmica, Tratamento de exceções, Genéricos, Coleções, Modelagem UML, Interface gráfica em ambientes orientados a objetos. Objetos persistentes.}\\
\multicolumn{4}{|p{16cm}|}{}\\
\hline
\multicolumn{4}{|p{16cm}|}{Conteúdo programático:}\\*
\multicolumn{4}{|p{\dimexpr 16cm + 6\tabcolsep\relax}|}{\begin{enumerate}[start=1]\item Visão geral do paradigma de programação orientada a objetos\end{enumerate}}\\
\multicolumn{4}{|p{\dimexpr 16cm + 6\tabcolsep\relax}|}{\begin{enumerate}[start=2]\item Modelagem UML\end{enumerate}}\\
\multicolumn{4}{|p{\dimexpr 16cm + 6\tabcolsep\relax}|}{\begin{enumerate}[start=3]\item Classes, objetos, mensagens\end{enumerate}}\\
\multicolumn{4}{|p{\dimexpr 16cm + 6\tabcolsep\relax}|}{\begin{enumerate}[start=4]\item Herança\end{enumerate}}\\
\multicolumn{4}{|p{\dimexpr 16cm + 6\tabcolsep\relax}|}{\begin{enumerate}[start=5]\item Polimorfismo\end{enumerate}}\\
\multicolumn{4}{|p{\dimexpr 16cm + 6\tabcolsep\relax}|}{\begin{enumerate}[start=6]\item Ligação dinâmica\end{enumerate}}\\
\multicolumn{4}{|p{\dimexpr 16cm + 6\tabcolsep\relax}|}{\begin{enumerate}[start=7]\item Tratamento de exceções\end{enumerate}}\\
\multicolumn{4}{|p{\dimexpr 16cm + 6\tabcolsep\relax}|}{\begin{enumerate}[start=8]\item Genéricos\end{enumerate}}\\
\multicolumn{4}{|p{\dimexpr 16cm + 6\tabcolsep\relax}|}{\begin{enumerate}[start=9]\item Coleções\end{enumerate}}\\
\multicolumn{4}{|p{\dimexpr 16cm + 6\tabcolsep\relax}|}{\begin{enumerate}[start=10]\item Objetos persistentes\end{enumerate}}\\
\multicolumn{4}{|p{\dimexpr 16cm + 6\tabcolsep\relax}|}{\begin{enumerate}[start=11]\item Interfaces gráficas\end{enumerate}}\\
\multicolumn{4}{|p{\dimexpr 16cm + 6\tabcolsep\relax}|}{\begin{enumerate}[start=12]\item Programação em C++ e Java\end{enumerate}}\\
\multicolumn{4}{|p{\dimexpr 16cm + 6\tabcolsep\relax}|}{}\\
\hline
\multicolumn{4}{|p{16cm}|}{Bibliografia Básica:}\\*
\multicolumn{4}{|p{\dimexpr 16cm + 6\tabcolsep\relax}|}{\begin{itemize}\item DEITEL, Harvey M.; DEITEL, P. J. C++ como programar. 5. ed. São Paulo: Pearson Prentice Hall, 2006. Disponível em: https://plataforma.bvirtual.com.br/Acervo/Publicacao/338\end{itemize}}\\
\multicolumn{4}{|p{\dimexpr 16cm + 6\tabcolsep\relax}|}{\begin{itemize}\item DEITEL, Paul; DEITEL, Harvey. Java: como programar. 8. ed. São Paulo: Prentice-Hall, 2010. Disponível em: https://plataforma.bvirtual.com.br/Acervo/Publicacao/1142\end{itemize}}\\
\multicolumn{4}{|p{\dimexpr 16cm + 6\tabcolsep\relax}|}{\begin{itemize}\item FOWLER, Martin. UML essencial: um breve guia para a linguagem-padrão de modelagem de objetos. 3. ed. Bookman, 2005. Disponível em: https://integrada.minhabiblioteca.com.br/\#/books/9788560031382/cfi/6/2!/4/14/6/10@0:58.0\end{itemize}}\\
\multicolumn{4}{|p{\dimexpr 16cm + 6\tabcolsep\relax}|}{\begin{itemize}\item MEYER, Bertrand. Object-oriented software construction. 2. ed. Upper Saddle River, NJ: Prentice-Hall PTR, 1997.\end{itemize}}\\
\multicolumn{4}{|p{\dimexpr 16cm + 6\tabcolsep\relax}|}{\begin{itemize}\item BOOCH, Grady. Object-oriented analysis and design with applications. 3. ed. New Delhi: Pearson, 2009.\end{itemize}}\\
\multicolumn{4}{|p{16cm}|}{}\\
\hline
\multicolumn{4}{|p{16cm}|}{Bibliografia Complementar:}\\*
\multicolumn{4}{|p{\dimexpr 16cm + 6\tabcolsep\relax}|}{\begin{itemize}\item LEE, Richard C.; TEPFENHART, William M. UML e C++: guia prático de desenvolvimento orientado a objeto. Pearson, 2001. Disponível em: https://plataforma.bvirtual.com.br/Acervo/Publicacao/40\end{itemize}}\\
\multicolumn{4}{|p{\dimexpr 16cm + 6\tabcolsep\relax}|}{\begin{itemize}\item PAGE-JONES, Meilir. Fundamentos do desenho orientado a objeto com UML. São Paulo: Makron Books, 2001. Disponível em: https://plataforma.bvirtual.com.br/Acervo/Publicacao/33\end{itemize}}\\
\multicolumn{4}{|p{\dimexpr 16cm + 6\tabcolsep\relax}|}{\begin{itemize}\item FELIX, Rafael (org.). Programação orientada a objetos. Pearson, 2017. Disponível em: https://plataforma.bvirtual.com.br/Acervo/Publicacao/128217\end{itemize}}\\
\multicolumn{4}{|p{\dimexpr 16cm + 6\tabcolsep\relax}|}{\begin{itemize}\item SINTES, Anthony. Aprenda programação orientada a objetos em 21 dias. 5. ed. Pearson, 2014. Disponível em: https://plataforma.bvirtual.com.br/Acervo/Publicacao/8\end{itemize}}\\
\multicolumn{4}{|p{\dimexpr 16cm + 6\tabcolsep\relax}|}{\begin{itemize}\item SANTOS, Rafael. Introdução à programação orientada a objetos usando JAVA. 9. ed. Rio de Janeiro: Elsevier, 2003.\end{itemize}}\\
\multicolumn{4}{|p{\dimexpr 16cm + 6\tabcolsep\relax}|}{\begin{itemize}\item BARNES, David J.; KÖLLING, Michael. Programação orientada a objetos com Java: uma introdução prática usando o BlueJ. 4. ed. São Paulo: Prentice-Hall, 2009.\end{itemize}}\\
\multicolumn{4}{|p{\dimexpr 16cm + 6\tabcolsep\relax}|}{\begin{itemize}\item LARMAN, Craig. Utilizando UML e padrões: uma introdução à análise e ao projeto orientados a objetos. 2. ed. Porto Alegre: Bookman, 2004.\end{itemize}}\\
\hline
\end{longtable}
\end{center}

\clearpage
\begin{center}
\begin{tabular}{lccr}
 \multirow{3}{*}{\includegraphics[height=2.7cm]{brasao.png}} &
 \multicolumn{2}{c}{\bfseries UNIVERSIDADE FEDERAL DE OURO PRETO} &
 \ \ \ \ \multirow{3}{*}{\includegraphics[height=2.7cm]{ufop.png}} \\
 & \multicolumn{2}{c}{\bfseries PRÓ-REITORIA DE GRADUAÇÃO} & \\
 & \multicolumn{2}{c}{} & \\
 & \multicolumn{2}{c}{\Large\bfseries PROGRAMA DE DISCIPLINA} & \\
\end{tabular}
\end{center}

\begin{center}
\begin{longtable}{|p{4cm}|p{4cm}|p{4cm}|p{4cm}|}
\hline
\multicolumn{3}{|p{12cm}|}{Nome do Componente Curricular em Português:} &
\multicolumn{1}{p{4cm}|}{Código:} \\ 
\multicolumn{3}{|p{12cm}|}{\textbf{Projeto e Análise de Algoritmos}} &
\textbf{BCC241}\\ 
\multicolumn{3}{|p{12cm}|}{Nome do Componente Curricular em Inglês:} & \\ 
\multicolumn{3}{|p{12cm}|}{\textbf{Design and Analysis of Algorithms}} & \\ 
\hline
\multicolumn{3}{|p{12cm}|}{Nome e Sigla do Departamento} & Unidade Acadêmica: \\ 
\multicolumn{3}{|p{12cm}|}{Departamento de Computação (DECOM)} & {ICEB} \\ 
\hline
\multicolumn{4}{|p{16cm}|}{Modalidade de Oferta:
[X] presencial \hspace{0.6cm}
[ ] semipresencial \hspace{0.6cm}
[ ] à distância}\\
\hline
\multicolumn{2}{|p{8cm}|}{Carga horária semestral} &
\multicolumn{2}{p{8cm}|}{Carga horária semanal}\\
\hline
\multicolumn{1}{|p{4cm}|}{Total} &
\multicolumn{1}{p{4cm}|}{Extensionista} &
\multicolumn{1}{p{4cm}|}{Teórica} &
\multicolumn{1}{p{4cm}|}{Prática} \\ 
\multicolumn{1}{|p{4cm}|}{60\,horas} &
\multicolumn{1}{p{4cm}|}{0\;horas} &
\multicolumn{1}{p{4cm}|}{4\;horas/aula} &
\multicolumn{1}{p{4cm}|}{0\;horas/aula} \\ 
\hline
\multicolumn{4}{|p{16cm}|}{Ementa:}\\
\multicolumn{4}{|p{16cm}|}{}\\
\multicolumn{4}{|p{\dimexpr 16cm + 6\tabcolsep\relax}|}{Medidas de complexidade, análise assintótica de complexidade e notações Big O, Little o, Omega e Theta. Análise de algoritmos iterativos e recursivos. Medidas empíricas de performance. Estratégias de projeto de algoritmos: divisão e conquista, método guloso, programação dinâmica, backtracking, branch and bound, probabilístico e aproximado. Classes de complexidade: P, NP, NP-Completo e NP-Difícil.}\\
\multicolumn{4}{|p{16cm}|}{}\\
\hline
\multicolumn{4}{|p{16cm}|}{Conteúdo programático:}\\*
\multicolumn{4}{|p{\dimexpr 16cm + 6\tabcolsep\relax}|}{\begin{enumerate}[start=1]\item Medidas de complexidade, análise assintótica de complexidade e notação Big O, Little o, Omega e Theta: panorama e conceitos básicos; medidas de complexidade (tempo, espaço); análise assintótica\end{enumerate}}\\
\multicolumn{4}{|p{\dimexpr 16cm + 6\tabcolsep\relax}|}{\begin{enumerate}[start=2]\item Análise de algoritmos iterativos e recursivos: Teorema Mestre\end{enumerate}}\\
\multicolumn{4}{|p{\dimexpr 16cm + 6\tabcolsep\relax}|}{\begin{enumerate}[start=3]\item Medidas empíricas de performance\end{enumerate}}\\
\multicolumn{4}{|p{\dimexpr 16cm + 6\tabcolsep\relax}|}{\begin{enumerate}[start=4]\item Estratégias de projeto de algoritmos: divisão e conquista (MergeSort, Medianas, QuickSort e Exponencial); método guloso (conceitos, árvores geradoras mínimas - Prim e Kruskal, código de Huffman, cláusula de Horn, problema da mochila e seleção de atividades); programação dinâmica (conceitos, maior sequência crescente, distância de edição, problema da mochila e multiplicação de cadeia de matrizes); backtracking; branch and bound; probabilístico; aproximado\end{enumerate}}\\
\multicolumn{4}{|p{\dimexpr 16cm + 6\tabcolsep\relax}|}{\begin{enumerate}[start=5]\item Classes de complexidade: problemas de busca (decisão e otimização); classe P; classe NP; classe NP-Completo; NP-Difícil; redução de problemas\end{enumerate}}\\
\multicolumn{4}{|p{\dimexpr 16cm + 6\tabcolsep\relax}|}{}\\
\hline
\multicolumn{4}{|p{16cm}|}{Bibliografia Básica:}\\*
\multicolumn{4}{|p{\dimexpr 16cm + 6\tabcolsep\relax}|}{\begin{itemize}\item Dasgupta, Sanjoy. Algoritmos. São Paulo: McGraw-Hill, 2009.\end{itemize}}\\
\multicolumn{4}{|p{\dimexpr 16cm + 6\tabcolsep\relax}|}{\begin{itemize}\item Sedgewick, Robert. Algorithms. 4. ed. Upper Saddle River: Addison Wesley, 2011.\end{itemize}}\\
\multicolumn{4}{|p{\dimexpr 16cm + 6\tabcolsep\relax}|}{\begin{itemize}\item Cormen, Thomas H. et al.; tradução Vandenberg D. de Souza. Algoritmos: teoria e prática. Rio de Janeiro: Campus, 2002.\end{itemize}}\\
\multicolumn{4}{|p{16cm}|}{}\\
\hline
\multicolumn{4}{|p{16cm}|}{Bibliografia Complementar:}\\*
\multicolumn{4}{|p{\dimexpr 16cm + 6\tabcolsep\relax}|}{\begin{itemize}\item Ziviani, Nivio; Botelho, Fabiano Cupertino. Projeto de algoritmos: com implementações em Java e C++. São Paulo: Cengage Learning, 2006.\end{itemize}}\\
\multicolumn{4}{|p{\dimexpr 16cm + 6\tabcolsep\relax}|}{\begin{itemize}\item Knuth, Donald Ervin. The art of computer programming. Upper Saddle River: Addison Wesley, c2005.\end{itemize}}\\
\multicolumn{4}{|p{\dimexpr 16cm + 6\tabcolsep\relax}|}{\begin{itemize}\item Goodrich, Michael T.; Tamassia, Roberto; tradução Bernardo Copstein. Projeto de algoritmos: fundamentos, análise e exemplos da internet. Porto Alegre: Bookman, 2004.\end{itemize}}\\
\multicolumn{4}{|p{\dimexpr 16cm + 6\tabcolsep\relax}|}{\begin{itemize}\item Horowitz, Ellis. Fundamentals of computer algorithms. New Delhi: Galgotia, 1990.\end{itemize}}\\
\multicolumn{4}{|p{\dimexpr 16cm + 6\tabcolsep\relax}|}{\begin{itemize}\item Wirth, Niklaus. Algoritmos e estruturas de dados. Livros Técnicos e Científicos.\end{itemize}}\\
\hline
\end{longtable}
\end{center}

\clearpage
\begin{center}
\begin{tabular}{lccr}
 \multirow{3}{*}{\includegraphics[height=2.7cm]{brasao.png}} &
 \multicolumn{2}{c}{\bfseries UNIVERSIDADE FEDERAL DE OURO PRETO} &
 \ \ \ \ \multirow{3}{*}{\includegraphics[height=2.7cm]{ufop.png}} \\
 & \multicolumn{2}{c}{\bfseries PRÓ-REITORIA DE GRADUAÇÃO} & \\
 & \multicolumn{2}{c}{} & \\
 & \multicolumn{2}{c}{\Large\bfseries PROGRAMA DE DISCIPLINA} & \\
\end{tabular}
\end{center}

\begin{center}
\begin{longtable}{|p{4cm}|p{4cm}|p{4cm}|p{4cm}|}
\hline
\multicolumn{3}{|p{12cm}|}{Nome do Componente Curricular em Português:} &
\multicolumn{1}{p{4cm}|}{Código:} \\ 
\multicolumn{3}{|p{12cm}|}{\textbf{Inteligência de Negócios e Armazenamento de Dados}} &
\textbf{BIA107}\\ 
\multicolumn{3}{|p{12cm}|}{Nome do Componente Curricular em Inglês:} & \\ 
\multicolumn{3}{|p{12cm}|}{\textbf{Business Intelligence and Data}} & \\ 
\hline
\multicolumn{3}{|p{12cm}|}{Nome e Sigla do Departamento} & Unidade Acadêmica: \\ 
\multicolumn{3}{|p{12cm}|}{Departamento de Computação (DECOM)} & {ICEB} \\ 
\hline
\multicolumn{4}{|p{16cm}|}{Modalidade de Oferta:
[X] presencial \hspace{0.6cm}
[ ] semipresencial \hspace{0.6cm}
[ ] à distância}\\
\hline
\multicolumn{2}{|p{8cm}|}{Carga horária semestral} &
\multicolumn{2}{p{8cm}|}{Carga horária semanal}\\
\hline
\multicolumn{1}{|p{4cm}|}{Total} &
\multicolumn{1}{p{4cm}|}{Extensionista} &
\multicolumn{1}{p{4cm}|}{Teórica} &
\multicolumn{1}{p{4cm}|}{Prática} \\ 
\multicolumn{1}{|p{4cm}|}{60\,horas} &
\multicolumn{1}{p{4cm}|}{0\;horas} &
\multicolumn{1}{p{4cm}|}{4\;horas/aula} &
\multicolumn{1}{p{4cm}|}{0\;horas/aula} \\ 
\hline
\multicolumn{4}{|p{16cm}|}{Ementa:}\\
\multicolumn{4}{|p{16cm}|}{}\\
\multicolumn{4}{|p{\dimexpr 16cm + 6\tabcolsep\relax}|}{Conceitos e tecnologias para organização, modelagem e análise de dados em apoio à tomada de decisão. Banco de dados relacionais, não relacionais e multidimensionais. Processos de ETL, modelagem dimensional, OLAP e fundamentos de inteligência de negócios.}\\
\multicolumn{4}{|p{16cm}|}{}\\
\hline
\multicolumn{4}{|p{16cm}|}{Conteúdo programático:}\\*
\multicolumn{4}{|p{\dimexpr 16cm + 6\tabcolsep\relax}|}{\begin{enumerate}[start=1]\item Introdução à inteligência de negócios (BI): conceitos, objetivos e aplicações.\end{enumerate}}\\
\multicolumn{4}{|p{\dimexpr 16cm + 6\tabcolsep\relax}|}{\begin{enumerate}[start=2]\item Modelagem de dados: modelo relacional, MER e DER.\end{enumerate}}\\
\multicolumn{4}{|p{\dimexpr 16cm + 6\tabcolsep\relax}|}{\begin{enumerate}[start=3]\item Banco de dados relacionais: SQL, normalização, integridade e consultas.\end{enumerate}}\\
\multicolumn{4}{|p{\dimexpr 16cm + 6\tabcolsep\relax}|}{\begin{enumerate}[start=4]\item NoSQL e bancos de dados não estruturados: documentos, grafos e chave-valor.\end{enumerate}}\\
\multicolumn{4}{|p{\dimexpr 16cm + 6\tabcolsep\relax}|}{\begin{enumerate}[start=5]\item Modelagem dimensional: fatos, dimensões, estrelas e flocos de neve.\end{enumerate}}\\
\multicolumn{4}{|p{\dimexpr 16cm + 6\tabcolsep\relax}|}{\begin{enumerate}[start=6]\item Processos ETL: extração, transformação e carga de dados.\end{enumerate}}\\
\multicolumn{4}{|p{\dimexpr 16cm + 6\tabcolsep\relax}|}{\begin{enumerate}[start=7]\item Armazenamento de dados: data warehouse e data lake.\end{enumerate}}\\
\multicolumn{4}{|p{\dimexpr 16cm + 6\tabcolsep\relax}|}{\begin{enumerate}[start=8]\item Consultas analíticas e OLAP: operações roll-up, drill-down, slice e dice.\end{enumerate}}\\
\multicolumn{4}{|p{\dimexpr 16cm + 6\tabcolsep\relax}|}{\begin{enumerate}[start=9]\item Ferramentas de visualização e dashboards: princípios de design de relatórios.\end{enumerate}}\\
\multicolumn{4}{|p{\dimexpr 16cm + 6\tabcolsep\relax}|}{\begin{enumerate}[start=10]\item Estudos de caso: uso de dados em decisões empresariais e operacionais.\end{enumerate}}\\
\multicolumn{4}{|p{\dimexpr 16cm + 6\tabcolsep\relax}|}{}\\
\hline
\multicolumn{4}{|p{16cm}|}{Bibliografia Básica:}\\*
\multicolumn{4}{|p{\dimexpr 16cm + 6\tabcolsep\relax}|}{\begin{itemize}\item SCHAEDLER, Andrew; MENDES, Giselly Santos. Business intelligence. 1. ed. Curitiba: Intersaberes, 2021. E-book. Disponível em: https://plataforma.bvirtual.com.br. Acesso em: 22 dez. 2025.\end{itemize}}\\
\multicolumn{4}{|p{\dimexpr 16cm + 6\tabcolsep\relax}|}{\begin{itemize}\item SILBERSCHATZ, Abraham; KORTH, Henry F.; SUDARSHAN, S. Sistemas de bancos de dados. 6. ed. São Paulo: Pearson, 2013. E-book. Disponível em: https://plataforma.bvirtual.com.br. Acesso em: jun/2025.\end{itemize}}\\
\multicolumn{4}{|p{\dimexpr 16cm + 6\tabcolsep\relax}|}{\begin{itemize}\item VIDA, Edinilson da Silva; ALVES, Nicolli S. Rios; FERREIRA, Rafael G. Coimbra; et al. Data warehouse. Porto Alegre: Grupo A, 2021. E-book. Disponível em: <https://integrada.minhabiblioteca.com.br/reader/books/9786556901916>. Acesso em: 22 dez. 2025\end{itemize}}\\
\multicolumn{4}{|p{16cm}|}{}\\
\hline
\multicolumn{4}{|p{16cm}|}{Bibliografia Complementar:}\\*
\multicolumn{4}{|p{\dimexpr 16cm + 6\tabcolsep\relax}|}{\begin{itemize}\item DATE, C. J. Introdução a sistemas de bancos de dados. 8. ed. Rio de Janeiro: Campus, 2004. E-book. Disponível em: https://plataforma.bvirtual.com.br. Acesso em: jun/2025.\end{itemize}}\\
\multicolumn{4}{|p{\dimexpr 16cm + 6\tabcolsep\relax}|}{\begin{itemize}\item HELMERS, Shawn; LARSON, Baya Dewitt. Fundamentos de BI: inteligência de negócios orientada por dados. 2. ed. São Paulo: Alta Books, 2022. E-book. Disponível em: https://plataforma.bvirtual.com.br. Acesso em: jun/2025.\end{itemize}}\\
\multicolumn{4}{|p{\dimexpr 16cm + 6\tabcolsep\relax}|}{\begin{itemize}\item KRUG, Michael. Bancos de dados NoSQL: fundamentos e aplicações. 1. ed. São Paulo: Érica, 2021. E-book. Disponível em: https://plataforma.bvirtual.com.br. Acesso em: jun/2025.\end{itemize}}\\
\multicolumn{4}{|p{\dimexpr 16cm + 6\tabcolsep\relax}|}{\begin{itemize}\item PANIZ, David. NoSQL: como armazenar os dados de uma aplicação moderna. São Paulo, SP: Casa do Código, 2016. E-book. Disponível em: https://plataforma.bvirtual.com.br. Acesso em: 22 dez. 2025.\end{itemize}}\\
\multicolumn{4}{|p{\dimexpr 16cm + 6\tabcolsep\relax}|}{\begin{itemize}\item KIMBALL, Ralph; ROSS, Margy. The Data Warehouse Toolkit: The Definitive Guide to Dimensional Modeling. 3. ed. Indianapolis: Wiley, 2013.\end{itemize}}\\
\multicolumn{4}{|p{\dimexpr 16cm + 6\tabcolsep\relax}|}{\begin{itemize}\item PONNIAH, Paulraj. Data Warehousing Fundamentals. 1. ed. New York: Wiley, 2001.\end{itemize}}\\
\multicolumn{4}{|p{\dimexpr 16cm + 6\tabcolsep\relax}|}{\begin{itemize}\item FERREIRA, Rodrigo Siqueira. Banco de dados: teoria e prática. 1. ed. São Paulo: Novatec, 2020. E-book. Disponível em: https://plataforma.bvirtual.com.br. Acesso em: jun/2025.\end{itemize}}\\
\hline
\end{longtable}
\end{center}

\clearpage
\begin{center}
\begin{tabular}{lccr}
 \multirow{3}{*}{\includegraphics[height=2.7cm]{brasao.png}} &
 \multicolumn{2}{c}{\bfseries UNIVERSIDADE FEDERAL DE OURO PRETO} &
 \ \ \ \ \multirow{3}{*}{\includegraphics[height=2.7cm]{ufop.png}} \\
 & \multicolumn{2}{c}{\bfseries PRÓ-REITORIA DE GRADUAÇÃO} & \\
 & \multicolumn{2}{c}{} & \\
 & \multicolumn{2}{c}{\Large\bfseries PROGRAMA DE DISCIPLINA} & \\
\end{tabular}
\end{center}

\begin{center}
\begin{longtable}{|p{4cm}|p{4cm}|p{4cm}|p{4cm}|}
\hline
\multicolumn{3}{|p{12cm}|}{Nome do Componente Curricular em Português:} &
\multicolumn{1}{p{4cm}|}{Código:} \\ 
\multicolumn{3}{|p{12cm}|}{\textbf{Engenharia de Software I}} &
\textbf{BCC112}\\ 
\multicolumn{3}{|p{12cm}|}{Nome do Componente Curricular em Inglês:} & \\ 
\multicolumn{3}{|p{12cm}|}{\textbf{Software Engineering I}} & \\ 
\hline
\multicolumn{3}{|p{12cm}|}{Nome e Sigla do Departamento} & Unidade Acadêmica: \\ 
\multicolumn{3}{|p{12cm}|}{Departamento de Computação (DECOM)} & {ICEB} \\ 
\hline
\multicolumn{4}{|p{16cm}|}{Modalidade de Oferta:
[X] presencial \hspace{0.6cm}
[ ] semipresencial \hspace{0.6cm}
[ ] à distância}\\
\hline
\multicolumn{2}{|p{8cm}|}{Carga horária semestral} &
\multicolumn{2}{p{8cm}|}{Carga horária semanal}\\
\hline
\multicolumn{1}{|p{4cm}|}{Total} &
\multicolumn{1}{p{4cm}|}{Extensionista} &
\multicolumn{1}{p{4cm}|}{Teórica} &
\multicolumn{1}{p{4cm}|}{Prática} \\ 
\multicolumn{1}{|p{4cm}|}{60\,horas} &
\multicolumn{1}{p{4cm}|}{0\;horas} &
\multicolumn{1}{p{4cm}|}{4\;horas/aula} &
\multicolumn{1}{p{4cm}|}{0\;horas/aula} \\ 
\hline
\multicolumn{4}{|p{16cm}|}{Ementa:}\\
\multicolumn{4}{|p{16cm}|}{}\\
\multicolumn{4}{|p{\dimexpr 16cm + 6\tabcolsep\relax}|}{Introdução à Engenharia de Software. Modelagem Orientada por Objetos. Princípios de Projeto de Software. Padrões de Projeto de Software. Arquiteturas de Software. Testes de Software.}\\
\multicolumn{4}{|p{16cm}|}{}\\
\hline
\multicolumn{4}{|p{16cm}|}{Conteúdo programático:}\\*
\multicolumn{4}{|p{\dimexpr 16cm + 6\tabcolsep\relax}|}{\begin{enumerate}[start=1]\item Introdução à Engenharia de Software: definições, contexto e história; visão geral dos conceitos envolvidos no desenvolvimento de software (engenharia de requisitos; projeto de software arquitetural e dirigido por API; construção do software; teste de software; manutenção e evolução; gerência de configuração; gerência de projetos; modelos de software; processos de desenvolvimento de software; qualidade de software; aspectos econômicos; tipos de software)\end{enumerate}}\\
\multicolumn{4}{|p{\dimexpr 16cm + 6\tabcolsep\relax}|}{\begin{enumerate}[start=2]\item Modelagem Orientada por Objetos: Linguagem Unificada de Modelagem (UML): diagrama de estados; diagrama de classes; diagrama de pacotes; diagrama de sequência; diagrama de atividades; diagrama de componentes\end{enumerate}}\\
\multicolumn{4}{|p{\dimexpr 16cm + 6\tabcolsep\relax}|}{\begin{enumerate}[start=3]\item Princípios de Projeto: princípios básicos (integridade conceitual; ocultamento de informação; alta coesão; fraco acoplamento); projeto dirigido por API; princípios da inversão de controle e da injeção de dependência; princípios SOLID e outros (responsabilidade única; segregação de interfaces; inversão de dependências; prefira composição à herança; Demeter; aberto/fechado; substituição de Liskov); métricas de código fonte: tamanho (LOC - Lines of Code), coesão (LCOM - Lack of Cohesion Between Methods), acoplamento (CBO - Coupling Between Objects) e complexidade ciclomática (CC)\end{enumerate}}\\
\multicolumn{4}{|p{\dimexpr 16cm + 6\tabcolsep\relax}|}{\begin{enumerate}[start=4]\item Padrões de Projeto: contexto histórico e definição; tipos de padrões: criacionais (Singleton, Método de Fábrica, Fábrica Abstrata, Builder), estruturais (Proxy, Adapter, Facade, Decorator) e comportamentais (Strategy, Observer, Template Method, Visitor, Iterator); críticas ao uso de padrões de projeto\end{enumerate}}\\
\multicolumn{4}{|p{\dimexpr 16cm + 6\tabcolsep\relax}|}{\begin{enumerate}[start=5]\item Arquitetura: definição; desenvolvimento dirigido por componentes (CDD); padrões arquiteturais pioneiros (monolítico, MicroKernel, pipes e filtros, cliente/servidor, peer-to-peer, em camadas incluindo Clean); padrões arquiteturais recentes (Model-View-Controller (MVC), microsserviços incluindo SOA, orientada a mensagens, publish-subscribe)\end{enumerate}}\\
\multicolumn{4}{|p{\dimexpr 16cm + 6\tabcolsep\relax}|}{\begin{enumerate}[start=6]\item Testes: definições (teste, caso de teste (BDD), suíte de teste); tipos de testes: caixa-preta e caixa-branca; pirâmide de testes (unitários, funcionais, de integração, de sistema, de aceitação alfa e beta); princípios FIRST (rápidos (Fast), independentes, determinísticos (Repeatable), auto-verificáveis (Self-checking), escritos o quanto antes (Timely)); test smells (obscuros, com lógica condicional, de código duplicado); cobertura de teste; mocks e stubs de testes; testes de requisitos não-funcionais; desenvolvimento dirigido por testes (TDD) e integração contínua; ferramentas para automação de teste e integração contínua\end{enumerate}}\\
\multicolumn{4}{|p{\dimexpr 16cm + 6\tabcolsep\relax}|}{}\\
\hline
\multicolumn{4}{|p{16cm}|}{Bibliografia Básica:}\\*
\multicolumn{4}{|p{\dimexpr 16cm + 6\tabcolsep\relax}|}{\begin{itemize}\item VALENTE, M. T. Engenharia de software moderna: princípios e práticas para desenvolvimento de software com produtividade. 1. ed. UmLivro, 2020. ISBN 978-6500019506. https://engsoftmoderna.info/\end{itemize}}\\
\multicolumn{4}{|p{\dimexpr 16cm + 6\tabcolsep\relax}|}{\begin{itemize}\item PATTERSON, D. A.; FOX, A. Construindo software como serviço (SaaS): uma abordagem ágil usando computação em nuvem. 1st ed. Strawberry Canyon LLC, 2015. ISBN 9780984881277. http://www.saasbook.info/translations/br ou http://www.saasbook.info/first-edition\end{itemize}}\\
\multicolumn{4}{|p{\dimexpr 16cm + 6\tabcolsep\relax}|}{\begin{itemize}\item GAMMA, E. et al. Padrões de projetos: soluções reutilizáveis de software orientados a objetos. Reading: Addison Wesley, 1995. 395 p. (Addison-Wesley Professional Computing Series). ISBN 9788573076103.\end{itemize}}\\
\multicolumn{4}{|p{16cm}|}{}\\
\hline
\multicolumn{4}{|p{16cm}|}{Bibliografia Complementar:}\\*
\multicolumn{4}{|p{\dimexpr 16cm + 6\tabcolsep\relax}|}{\begin{itemize}\item SOMMERVILLE, I. Engenharia de software. 10. ed. São Paulo: Pearson, 2019. ISBN 9788543024974. https://plataforma.bvirtual.com.br/Leitor/Loader/168127/pdf/6?keep=True\end{itemize}}\\
\multicolumn{4}{|p{\dimexpr 16cm + 6\tabcolsep\relax}|}{\begin{itemize}\item FOWLER, M. UML distilled: a brief guide to the standard object modeling language. 3. ed. Boston: Addison Wesley, c2004. xxx, 175 p. (Object technology series). ISBN 9780321193681. https://integrada.minhabiblioteca.com.br/\#/books/9788560031382/cfi/6/2!/4/18/52@0:76.8\end{itemize}}\\
\multicolumn{4}{|p{\dimexpr 16cm + 6\tabcolsep\relax}|}{\begin{itemize}\item SUMMERFIELD, Mark. Advanced Qt programming: creating great software with C++ and Qt4. New York: Addison Wesley, 2010. 536 p. ISBN 9780321635907. http://www.qtrac.eu/aqpbook.html\end{itemize}}\\
\multicolumn{4}{|p{\dimexpr 16cm + 6\tabcolsep\relax}|}{\begin{itemize}\item BRAGA, Henrique Cacique (org.). Teste de software. 1. ed. São Paulo: Pearson Education do Brasil, 2016. 139 p. ISBN 9788543020211. https://plataforma.bvirtual.com.br/Acervo/Publicacao/150962\end{itemize}}\\
\multicolumn{4}{|p{\dimexpr 16cm + 6\tabcolsep\relax}|}{\begin{itemize}\item CRISPIN, Lisa; GREGORY, Janet. Agile testing: a practical guide for testers and agile teams. Upper Saddle River, NJ: Addison Wesley, c2009. xli, 533 p. (The Addison-Wesley signature series). ISBN 9780321534460.\end{itemize}}\\
\hline
\end{longtable}
\end{center}

\clearpage
\begin{center}
\begin{tabular}{lccr}
 \multirow{3}{*}{\includegraphics[height=2.7cm]{brasao.png}} &
 \multicolumn{2}{c}{\bfseries UNIVERSIDADE FEDERAL DE OURO PRETO} &
 \ \ \ \ \multirow{3}{*}{\includegraphics[height=2.7cm]{ufop.png}} \\
 & \multicolumn{2}{c}{\bfseries PRÓ-REITORIA DE GRADUAÇÃO} & \\
 & \multicolumn{2}{c}{} & \\
 & \multicolumn{2}{c}{\Large\bfseries PROGRAMA DE DISCIPLINA} & \\
\end{tabular}
\end{center}

\begin{center}
\begin{longtable}{|p{4cm}|p{4cm}|p{4cm}|p{4cm}|}
\hline
\multicolumn{3}{|p{12cm}|}{Nome do Componente Curricular em Português:} &
\multicolumn{1}{p{4cm}|}{Código:} \\ 
\multicolumn{3}{|p{12cm}|}{\textbf{Redes Neurais e Aprendizagem em Profundidade}} &
\textbf{BCC129}\\ 
\multicolumn{3}{|p{12cm}|}{Nome do Componente Curricular em Inglês:} & \\ 
\multicolumn{3}{|p{12cm}|}{\textbf{Neural Networks and Deep Learning}} & \\ 
\hline
\multicolumn{3}{|p{12cm}|}{Nome e Sigla do Departamento} & Unidade Acadêmica: \\ 
\multicolumn{3}{|p{12cm}|}{Departamento de Computação (DECOM)} & {ICEB} \\ 
\hline
\multicolumn{4}{|p{16cm}|}{Modalidade de Oferta:
[X] presencial \hspace{0.6cm}
[ ] semipresencial \hspace{0.6cm}
[ ] à distância}\\
\hline
\multicolumn{2}{|p{8cm}|}{Carga horária semestral} &
\multicolumn{2}{p{8cm}|}{Carga horária semanal}\\
\hline
\multicolumn{1}{|p{4cm}|}{Total} &
\multicolumn{1}{p{4cm}|}{Extensionista} &
\multicolumn{1}{p{4cm}|}{Teórica} &
\multicolumn{1}{p{4cm}|}{Prática} \\ 
\multicolumn{1}{|p{4cm}|}{60\,horas} &
\multicolumn{1}{p{4cm}|}{0\;horas} &
\multicolumn{1}{p{4cm}|}{2\;horas/aula} &
\multicolumn{1}{p{4cm}|}{2\;horas/aula} \\ 
\hline
\multicolumn{4}{|p{16cm}|}{Ementa:}\\
\multicolumn{4}{|p{16cm}|}{}\\
\multicolumn{4}{|p{\dimexpr 16cm + 6\tabcolsep\relax}|}{Regressão Linear e Regressão Logística; Otimização por descida do gradiente; Retro-propagação; Neurônio artificial; Rede neural de múltiplas camadas; Redes convolucionais; Arquiteturas de redes de convolução para classificação, detecção e segmentação; Modelos sequenciais; Mecanismo de atenção; Modelos Generativos; Estudo de caso.}\\
\multicolumn{4}{|p{16cm}|}{}\\
\hline
\multicolumn{4}{|p{16cm}|}{Conteúdo programático:}\\*
\multicolumn{4}{|p{\dimexpr 16cm + 6\tabcolsep\relax}|}{\begin{enumerate}[start=1]\item Introdução à aprendizagem em profundidade: o que é uma rede neural; breve histórico; por que a aprendizagem em profundidade está em alta\end{enumerate}}\\
\multicolumn{4}{|p{\dimexpr 16cm + 6\tabcolsep\relax}|}{\begin{enumerate}[start=2]\item Noções básicas de redes neurais: regressão linear; regressão logística, função Softmax, função de custo; descida do gradiente; vetorização; perceptron e visão geral de um neurônio artificial; representação de uma rede neural; computando a saída da rede neural; funções de ativação; descida de gradiente para redes neurais; intuição de retro-propagação; importância da inicialização dos pesos\end{enumerate}}\\
\multicolumn{4}{|p{\dimexpr 16cm + 6\tabcolsep\relax}|}{\begin{enumerate}[start=3]\item Redes neurais profundas: rede neural de múltiplas camadas; propagação em uma rede profunda; por que representações profundas; construindo arquiteturas de redes neurais profundas; propagação e retro-propagação; parâmetros vs hiperparâmetros\end{enumerate}}\\
\multicolumn{4}{|p{\dimexpr 16cm + 6\tabcolsep\relax}|}{\begin{enumerate}[start=4]\item Fundamentos das redes neurais convolucionais: visão computacional; redes neurais convolucionais; padding; stride; pooling; dropout; exemplo de CNN; por que usar operações de convolução\end{enumerate}}\\
\multicolumn{4}{|p{\dimexpr 16cm + 6\tabcolsep\relax}|}{\begin{enumerate}[start=5]\item Modelos convolutivos profundos: redes convolucionais clássicas; arquiteturas convolucionais modernas; normalização por lote; transferência de aprendizagem e ajuste fino; aumento de dados; detecção e segmentação de objetos\end{enumerate}}\\
\multicolumn{4}{|p{\dimexpr 16cm + 6\tabcolsep\relax}|}{\begin{enumerate}[start=6]\item Modelos sequenciais: processando sequências; processamento de linguagem natural; redes recorrentes e unidade recorrente simples (GRU); LSTMs; redes recorrentes profundas; redes recorrentes bidirecionais; redes recorrentes modernas\end{enumerate}}\\
\multicolumn{4}{|p{\dimexpr 16cm + 6\tabcolsep\relax}|}{\begin{enumerate}[start=7]\item Mecanismo de atenção: o que é o mecanismo de atenção; modelo com atenção em múltiplas cabeças; auto-atenção; Transformer\end{enumerate}}\\
\multicolumn{4}{|p{\dimexpr 16cm + 6\tabcolsep\relax}|}{\begin{enumerate}[start=8]\item Modelos generativos: arquiteturas encoder-decoder; auto-encoders; redes generativas adversariais\end{enumerate}}\\
\multicolumn{4}{|p{\dimexpr 16cm + 6\tabcolsep\relax}|}{}\\
\hline
\multicolumn{4}{|p{16cm}|}{Bibliografia Básica:}\\*
\multicolumn{4}{|p{\dimexpr 16cm + 6\tabcolsep\relax}|}{\begin{itemize}\item ZHANG, A.; LIPTON, Z. C.; LI, M.; SMOLA, A. J. Dive into deep learning. Unpublished draft, 2019. Disponível em: https://d2l.ai/\end{itemize}}\\
\multicolumn{4}{|p{\dimexpr 16cm + 6\tabcolsep\relax}|}{\begin{itemize}\item GOODFELLOW, Ian; BENGIO, Yoshua; COURVILLE, Aaron. Deep learning. MIT Press, 2016. http://www.deeplearningbook.org\end{itemize}}\\
\multicolumn{4}{|p{\dimexpr 16cm + 6\tabcolsep\relax}|}{\begin{itemize}\item BISHOP, Christopher M. Pattern recognition and machine learning. New York: Springer, 2006.\end{itemize}}\\
\multicolumn{4}{|p{\dimexpr 16cm + 6\tabcolsep\relax}|}{\begin{itemize}\item POOLE, David L.; MACKWORTH, Alan K. Artificial intelligence: foundations of computational agents. Cambridge University Press, 2010. Disponível em: https://artint.info/\end{itemize}}\\
\multicolumn{4}{|p{\dimexpr 16cm + 6\tabcolsep\relax}|}{\begin{itemize}\item CASTRO, Leandro Nunes de. Fundamentals of natural computing: basic concepts, algorithms, and applications. Boca Raton: Chapman \& Hall/CRC, 2006.\end{itemize}}\\
\multicolumn{4}{|p{16cm}|}{}\\
\hline
\multicolumn{4}{|p{16cm}|}{Bibliografia Complementar:}\\*
\multicolumn{4}{|p{\dimexpr 16cm + 6\tabcolsep\relax}|}{\begin{itemize}\item BISHOP, Christopher M. Neural networks for pattern recognition. Oxford: Oxford University Press, 2010.\end{itemize}}\\
\multicolumn{4}{|p{\dimexpr 16cm + 6\tabcolsep\relax}|}{\begin{itemize}\item BRAGA, Antonio de Padua; CARVALHO, Andre Ponce de Leon F. de; LUDERMIR, Teresa Bernarda. Redes neurais artificiais: teoria e aplicações. Rio de Janeiro: LTC, 2000.\end{itemize}}\\
\multicolumn{4}{|p{\dimexpr 16cm + 6\tabcolsep\relax}|}{\begin{itemize}\item HAYKIN, Simon. Redes neurais: princípios e prática. Tradução de Paulo Martins Engel. 2. ed. Porto Alegre: Bookman, 2001.\end{itemize}}\\
\multicolumn{4}{|p{\dimexpr 16cm + 6\tabcolsep\relax}|}{\begin{itemize}\item THEODORIDIS, Sergios; PIKRAKIS, Aggelos; KOUTROUMBAS, Konstantinos; CAVOURAS, Dionisis. An introduction to pattern recognition: a Matlab approach. Academic Press, 2010.\end{itemize}}\\
\multicolumn{4}{|p{\dimexpr 16cm + 6\tabcolsep\relax}|}{\begin{itemize}\item BRATKO, I. Prolog programming for artificial intelligence. 3. ed. Addison-Wesley, 2000.\end{itemize}}\\
\hline
\end{longtable}
\end{center}

\clearpage
\begin{center}
\begin{tabular}{lccr}
 \multirow{3}{*}{\includegraphics[height=2.7cm]{brasao.png}} &
 \multicolumn{2}{c}{\bfseries UNIVERSIDADE FEDERAL DE OURO PRETO} &
 \ \ \ \ \multirow{3}{*}{\includegraphics[height=2.7cm]{ufop.png}} \\
 & \multicolumn{2}{c}{\bfseries PRÓ-REITORIA DE GRADUAÇÃO} & \\
 & \multicolumn{2}{c}{} & \\
 & \multicolumn{2}{c}{\Large\bfseries PROGRAMA DE DISCIPLINA} & \\
\end{tabular}
\end{center}

\begin{center}
\begin{longtable}{|p{4cm}|p{4cm}|p{4cm}|p{4cm}|}
\hline
\multicolumn{3}{|p{12cm}|}{Nome do Componente Curricular em Português:} &
\multicolumn{1}{p{4cm}|}{Código:} \\ 
\multicolumn{3}{|p{12cm}|}{\textbf{Introdução à Otimização}} &
\textbf{BCC107}\\ 
\multicolumn{3}{|p{12cm}|}{Nome do Componente Curricular em Inglês:} & \\ 
\multicolumn{3}{|p{12cm}|}{\textbf{Introduction to Optimization}} & \\ 
\hline
\multicolumn{3}{|p{12cm}|}{Nome e Sigla do Departamento} & Unidade Acadêmica: \\ 
\multicolumn{3}{|p{12cm}|}{Departamento de Computação (DECOM)} & {ICEB} \\ 
\hline
\multicolumn{4}{|p{16cm}|}{Modalidade de Oferta:
[X] presencial \hspace{0.6cm}
[ ] semipresencial \hspace{0.6cm}
[ ] à distância}\\
\hline
\multicolumn{2}{|p{8cm}|}{Carga horária semestral} &
\multicolumn{2}{p{8cm}|}{Carga horária semanal}\\
\hline
\multicolumn{1}{|p{4cm}|}{Total} &
\multicolumn{1}{p{4cm}|}{Extensionista} &
\multicolumn{1}{p{4cm}|}{Teórica} &
\multicolumn{1}{p{4cm}|}{Prática} \\ 
\multicolumn{1}{|p{4cm}|}{60\,horas} &
\multicolumn{1}{p{4cm}|}{0\;horas} &
\multicolumn{1}{p{4cm}|}{4\;horas/aula} &
\multicolumn{1}{p{4cm}|}{0\;horas/aula} \\ 
\hline
\multicolumn{4}{|p{16cm}|}{Ementa:}\\
\multicolumn{4}{|p{16cm}|}{}\\
\multicolumn{4}{|p{\dimexpr 16cm + 6\tabcolsep\relax}|}{Programação linear e inteira: formulação, modelagem, algoritmo Simplex, planos de corte, métodos de enumeração implícita; programação não linear: conceitos básicos e condições de otimalidade, algoritmos diretos e algoritmos indiretos.}\\
\multicolumn{4}{|p{16cm}|}{}\\
\hline
\multicolumn{4}{|p{16cm}|}{Conteúdo programático:}\\*
\multicolumn{4}{|p{\dimexpr 16cm + 6\tabcolsep\relax}|}{\begin{enumerate}[start=1]\item Otimização: introdução\end{enumerate}}\\
\multicolumn{4}{|p{\dimexpr 16cm + 6\tabcolsep\relax}|}{\begin{enumerate}[start=2]\item Programação Linear e Inteira: modelagem em programação linear e inteira; forma-padrão de um Problema de Programação Linear (PPL); solução gráfica de um PPL; fundamentação teórica do método SIMPLEX (introdução, caracterização do conjunto de soluções viáveis e vértice do politopo); o algoritmo SIMPLEX (geração de soluções básicas viáveis, método das duas fases e interpretação geométrica); planos de corte; enumeração implícita; otimização inteira em redes\end{enumerate}}\\
\multicolumn{4}{|p{\dimexpr 16cm + 6\tabcolsep\relax}|}{\begin{enumerate}[start=3]\item Programação não linear: conceitos básicos; condições de otimalidade para problemas irrestritos e restritos; algoritmos de busca diretos; algoritmos de busca indiretos\end{enumerate}}\\
\multicolumn{4}{|p{\dimexpr 16cm + 6\tabcolsep\relax}|}{}\\
\hline
\multicolumn{4}{|p{16cm}|}{Bibliografia Básica:}\\*
\multicolumn{4}{|p{\dimexpr 16cm + 6\tabcolsep\relax}|}{\begin{itemize}\item GOLDBARG, M. C.; LUNA, H. P. L. Otimização combinatória e programação linear: modelos e algoritmos. Rio de Janeiro: Campus, 2000.\end{itemize}}\\
\multicolumn{4}{|p{\dimexpr 16cm + 6\tabcolsep\relax}|}{\begin{itemize}\item RAO, S. S. Engineering optimization: theory and practice. 3. ed. New York: Wiley-Interscience, 1996.\end{itemize}}\\
\multicolumn{4}{|p{\dimexpr 16cm + 6\tabcolsep\relax}|}{\begin{itemize}\item TAHA, H. A. Pesquisa operacional. 8. ed. São Paulo: Prentice-Hall, 2008.\end{itemize}}\\
\multicolumn{4}{|p{16cm}|}{}\\
\hline
\multicolumn{4}{|p{16cm}|}{Bibliografia Complementar:}\\*
\multicolumn{4}{|p{\dimexpr 16cm + 6\tabcolsep\relax}|}{\begin{itemize}\item LACHTERMACHER, G. Pesquisa operacional na tomada de decisões: modelagem em Excel. 3. ed. Rio de Janeiro: Elsevier, 2007.\end{itemize}}\\
\multicolumn{4}{|p{\dimexpr 16cm + 6\tabcolsep\relax}|}{\begin{itemize}\item WINSTON, W. L. Operations research: applications and algorithms. 4. ed. Belmont: Thomson Brooks Cole, 2004.\end{itemize}}\\
\multicolumn{4}{|p{\dimexpr 16cm + 6\tabcolsep\relax}|}{\begin{itemize}\item BAZARAA, M. S.; SHERALI, H. D.; SHETTY, C. M. Nonlinear programming: theory and algorithms. 3. ed. Hoboken, NJ: Wiley-Interscience, 2006.\end{itemize}}\\
\multicolumn{4}{|p{\dimexpr 16cm + 6\tabcolsep\relax}|}{\begin{itemize}\item BAZARAA, M. S.; JARVIS, J. J.; SHERALI, H. D. Linear programming and network flows. 4. ed. New York: J. Wiley, 2010.\end{itemize}}\\
\multicolumn{4}{|p{\dimexpr 16cm + 6\tabcolsep\relax}|}{\begin{itemize}\item GONZALEZ, T. F. Handbook of approximation algorithms and metaheuristics. New York: Chapman \& Hall/CRC, 2007.\end{itemize}}\\
\hline
\end{longtable}
\end{center}

\clearpage
\begin{center}
\begin{tabular}{lccr}
 \multirow{3}{*}{\includegraphics[height=2.7cm]{brasao.png}} &
 \multicolumn{2}{c}{\bfseries UNIVERSIDADE FEDERAL DE OURO PRETO} &
 \ \ \ \ \multirow{3}{*}{\includegraphics[height=2.7cm]{ufop.png}} \\
 & \multicolumn{2}{c}{\bfseries PRÓ-REITORIA DE GRADUAÇÃO} & \\
 & \multicolumn{2}{c}{} & \\
 & \multicolumn{2}{c}{\Large\bfseries PROGRAMA DE DISCIPLINA} & \\
\end{tabular}
\end{center}

\begin{center}
\begin{longtable}{|p{4cm}|p{4cm}|p{4cm}|p{4cm}|}
\hline
\multicolumn{3}{|p{12cm}|}{Nome do Componente Curricular em Português:} &
\multicolumn{1}{p{4cm}|}{Código:} \\ 
\multicolumn{3}{|p{12cm}|}{\textbf{Aprendizado por Reforço}} &
\textbf{BCC131}\\ 
\multicolumn{3}{|p{12cm}|}{Nome do Componente Curricular em Inglês:} & \\ 
\multicolumn{3}{|p{12cm}|}{\textbf{Reinforcement Learning}} & \\ 
\hline
\multicolumn{3}{|p{12cm}|}{Nome e Sigla do Departamento} & Unidade Acadêmica: \\ 
\multicolumn{3}{|p{12cm}|}{Departamento de Computação (DECOM)} & {ICEB} \\ 
\hline
\multicolumn{4}{|p{16cm}|}{Modalidade de Oferta:
[X] presencial \hspace{0.6cm}
[ ] semipresencial \hspace{0.6cm}
[ ] à distância}\\
\hline
\multicolumn{2}{|p{8cm}|}{Carga horária semestral} &
\multicolumn{2}{p{8cm}|}{Carga horária semanal}\\
\hline
\multicolumn{1}{|p{4cm}|}{Total} &
\multicolumn{1}{p{4cm}|}{Extensionista} &
\multicolumn{1}{p{4cm}|}{Teórica} &
\multicolumn{1}{p{4cm}|}{Prática} \\ 
\multicolumn{1}{|p{4cm}|}{60\,horas} &
\multicolumn{1}{p{4cm}|}{0\;horas} &
\multicolumn{1}{p{4cm}|}{4\;horas/aula} &
\multicolumn{1}{p{4cm}|}{0\;horas/aula} \\ 
\hline
\multicolumn{4}{|p{16cm}|}{Ementa:}\\
\multicolumn{4}{|p{16cm}|}{}\\
\multicolumn{4}{|p{\dimexpr 16cm + 6\tabcolsep\relax}|}{Aprendizado por reforço: formulação, modelagem, métodos baseados em valores, métodos baseados em políticas, métodos baseados em modelos, aprendizado por imitação, meta-learning.}\\
\multicolumn{4}{|p{16cm}|}{}\\
\hline
\multicolumn{4}{|p{16cm}|}{Conteúdo programático:}\\*
\multicolumn{4}{|p{\dimexpr 16cm + 6\tabcolsep\relax}|}{\begin{enumerate}[start=1]\item Formulação de problemas de aprendizado por reforço (espaço de ações, espaço de estados, ambiente)\end{enumerate}}\\
\multicolumn{4}{|p{\dimexpr 16cm + 6\tabcolsep\relax}|}{\begin{enumerate}[start=2]\item Métodos tabulares (Programação Dinâmica, Monte Carlo, Diferença Temporal)\end{enumerate}}\\
\multicolumn{4}{|p{\dimexpr 16cm + 6\tabcolsep\relax}|}{\begin{enumerate}[start=3]\item Aproximação de funções (Deep Q-networks)\end{enumerate}}\\
\multicolumn{4}{|p{\dimexpr 16cm + 6\tabcolsep\relax}|}{\begin{enumerate}[start=4]\item Gradientes de políticas (REINFORCE)\end{enumerate}}\\
\multicolumn{4}{|p{\dimexpr 16cm + 6\tabcolsep\relax}|}{\begin{enumerate}[start=5]\item Aprendizado por reforço baseado em modelos\end{enumerate}}\\
\multicolumn{4}{|p{\dimexpr 16cm + 6\tabcolsep\relax}|}{\begin{enumerate}[start=6]\item Aprendizado por imitação\end{enumerate}}\\
\multicolumn{4}{|p{\dimexpr 16cm + 6\tabcolsep\relax}|}{\begin{enumerate}[start=7]\item Meta-learning\end{enumerate}}\\
\multicolumn{4}{|p{\dimexpr 16cm + 6\tabcolsep\relax}|}{}\\
\hline
\multicolumn{4}{|p{16cm}|}{Bibliografia Básica:}\\*
\multicolumn{4}{|p{\dimexpr 16cm + 6\tabcolsep\relax}|}{\begin{itemize}\item SUTTON, Richard S.; BARTO, Andrew G. Reinforcement learning: an introduction. 2. ed. Cambridge: MIT Press, 2018.\end{itemize}}\\
\multicolumn{4}{|p{\dimexpr 16cm + 6\tabcolsep\relax}|}{\begin{itemize}\item GRAESSER, Laura; KENG, Wah Loon. Foundations of deep reinforcement learning: theory and practice in Python. Boston: Addison-Wesley Professional, 2019.\end{itemize}}\\
\multicolumn{4}{|p{\dimexpr 16cm + 6\tabcolsep\relax}|}{\begin{itemize}\item LUGER, George F. Inteligência artificial. 6. ed. São Paulo: Pearson, 2014.\end{itemize}}\\
\multicolumn{4}{|p{16cm}|}{}\\
\hline
\multicolumn{4}{|p{16cm}|}{Bibliografia Complementar:}\\*
\multicolumn{4}{|p{\dimexpr 16cm + 6\tabcolsep\relax}|}{\begin{itemize}\item RUSSELL, Stuart J.; NORVIG, Peter. Artificial intelligence: a modern approach. 2. ed. Upper Saddle River: Prentice Hall, 2003.\end{itemize}}\\
\multicolumn{4}{|p{\dimexpr 16cm + 6\tabcolsep\relax}|}{\begin{itemize}\item POOLE, David L.; MACKWORTH, Alan K. Artificial intelligence: foundations of computational agents. 2. ed. Cambridge: Cambridge University Press, 2017. Disponível em: https://artint.info/2e/html/ArtInt2e.html.\end{itemize}}\\
\multicolumn{4}{|p{\dimexpr 16cm + 6\tabcolsep\relax}|}{\begin{itemize}\item FACELI, Katti; LORENA, Ana C.; GAMA, João; et al. Inteligência artificial: uma abordagem de aprendizado de máquina. Porto Alegre: Grupo GEN, 2021. Disponível em: https://integrada.minhabiblioteca.com.br/\#/books/9788521637509.\end{itemize}}\\
\hline
\end{longtable}
\end{center}

\clearpage
\begin{center}
\begin{tabular}{lccr}
 \multirow{3}{*}{\includegraphics[height=2.7cm]{brasao.png}} &
 \multicolumn{2}{c}{\bfseries UNIVERSIDADE FEDERAL DE OURO PRETO} &
 \ \ \ \ \multirow{3}{*}{\includegraphics[height=2.7cm]{ufop.png}} \\
 & \multicolumn{2}{c}{\bfseries PRÓ-REITORIA DE GRADUAÇÃO} & \\
 & \multicolumn{2}{c}{} & \\
 & \multicolumn{2}{c}{\Large\bfseries PROGRAMA DE DISCIPLINA} & \\
\end{tabular}
\end{center}

\begin{center}
\begin{longtable}{|p{4cm}|p{4cm}|p{4cm}|p{4cm}|}
\hline
\multicolumn{3}{|p{12cm}|}{Nome do Componente Curricular em Português:} &
\multicolumn{1}{p{4cm}|}{Código:} \\ 
\multicolumn{3}{|p{12cm}|}{\textbf{Fundamentos da Computação Paralela e Distribuída para Inteligência Artificial}} &
\textbf{BIA109}\\ 
\multicolumn{3}{|p{12cm}|}{Nome do Componente Curricular em Inglês:} & \\ 
\multicolumn{3}{|p{12cm}|}{\textbf{Fundamentals of Parallel and Distributed Computing for Artificial Intelligence}} & \\ 
\hline
\multicolumn{3}{|p{12cm}|}{Nome e Sigla do Departamento} & Unidade Acadêmica: \\ 
\multicolumn{3}{|p{12cm}|}{Departamento de Computação (DECOM)} & {ICEB} \\ 
\hline
\multicolumn{4}{|p{16cm}|}{Modalidade de Oferta:
[X] presencial \hspace{0.6cm}
[ ] semipresencial \hspace{0.6cm}
[ ] à distância}\\
\hline
\multicolumn{2}{|p{8cm}|}{Carga horária semestral} &
\multicolumn{2}{p{8cm}|}{Carga horária semanal}\\
\hline
\multicolumn{1}{|p{4cm}|}{Total} &
\multicolumn{1}{p{4cm}|}{Extensionista} &
\multicolumn{1}{p{4cm}|}{Teórica} &
\multicolumn{1}{p{4cm}|}{Prática} \\ 
\multicolumn{1}{|p{4cm}|}{60\,horas} &
\multicolumn{1}{p{4cm}|}{0\;horas} &
\multicolumn{1}{p{4cm}|}{4\;horas/aula} &
\multicolumn{1}{p{4cm}|}{0\;horas/aula} \\ 
\hline
\multicolumn{4}{|p{16cm}|}{Ementa:}\\
\multicolumn{4}{|p{16cm}|}{}\\
\multicolumn{4}{|p{\dimexpr 16cm + 6\tabcolsep\relax}|}{Estudo dos fundamentos computacionais para a execução de aplicações de Inteligência Artificial. Modelo de Von Neumann e suas limitações, concorrência, paralelismo e distribuição. Processos, threads, sincronização e paralelismo de tarefas e de dados. Métricas de desempenho e escalabilidade. Arquiteturas paralelas e sistemas distribuídos, computação em nuvem para IA e técnicas de treinamento paralelo e distribuído de modelos. Estudo de caso integrador em aprendizado dividido, distribuído e federado.}\\
\multicolumn{4}{|p{16cm}|}{}\\
\hline
\multicolumn{4}{|p{16cm}|}{Conteúdo programático:}\\*
\multicolumn{4}{|p{\dimexpr 16cm + 6\tabcolsep\relax}|}{\begin{enumerate}[start=1]\item Modelo de Von Neumann, gargalos de comunicação e limitações para workloads de Inteligência Artificial\end{enumerate}}\\
\multicolumn{4}{|p{\dimexpr 16cm + 6\tabcolsep\relax}|}{\begin{enumerate}[start=2]\item Evolução da infraestrutura de computação para Inteligência Artificial\end{enumerate}}\\
\multicolumn{4}{|p{\dimexpr 16cm + 6\tabcolsep\relax}|}{\begin{enumerate}[start=3]\item Concorrência, paralelismo e distribuição: definições e distinções conceituais\end{enumerate}}\\
\multicolumn{4}{|p{\dimexpr 16cm + 6\tabcolsep\relax}|}{\begin{enumerate}[start=4]\item Processos, threads e execução concorrente em memória compartilhada\end{enumerate}}\\
\multicolumn{4}{|p{\dimexpr 16cm + 6\tabcolsep\relax}|}{\begin{enumerate}[start=5]\item Condições de corrida e princípios básicos de sincronização\end{enumerate}}\\
\multicolumn{4}{|p{\dimexpr 16cm + 6\tabcolsep\relax}|}{\begin{enumerate}[start=6]\item Paralelismo de tarefas e de dados\end{enumerate}}\\
\multicolumn{4}{|p{\dimexpr 16cm + 6\tabcolsep\relax}|}{\begin{enumerate}[start=7]\item Métricas de desempenho e escalabilidade: speedup, eficiência, Leis de Amdahl e Gustafson\end{enumerate}}\\
\multicolumn{4}{|p{\dimexpr 16cm + 6\tabcolsep\relax}|}{\begin{enumerate}[start=8]\item Taxonomia de Flynn e arquiteturas paralelas e distribuídas modernas\end{enumerate}}\\
\multicolumn{4}{|p{\dimexpr 16cm + 6\tabcolsep\relax}|}{\begin{enumerate}[start=9]\item Fundamentos de sistemas distribuídos: memória distribuída, comunicação entre processos, sistemas de nomes e sincronização\end{enumerate}}\\
\multicolumn{4}{|p{\dimexpr 16cm + 6\tabcolsep\relax}|}{\begin{enumerate}[start=10]\item Computação em nuvem para Inteligência Artificial: IaaS, PaaS, SaaS, elasticidade, escalabilidade e custos computacionais\end{enumerate}}\\
\multicolumn{4}{|p{\dimexpr 16cm + 6\tabcolsep\relax}|}{\begin{enumerate}[start=11]\item Treinamento paralelo de modelos de Inteligência Artificial\end{enumerate}}\\
\multicolumn{4}{|p{\dimexpr 16cm + 6\tabcolsep\relax}|}{\begin{enumerate}[start=12]\item Treinamento distribuído de modelos de Inteligência Artificial\end{enumerate}}\\
\multicolumn{4}{|p{\dimexpr 16cm + 6\tabcolsep\relax}|}{\begin{enumerate}[start=13]\item Estudo de caso: Introdução ao aprendizado dividido (split learning), distribuído (distributed learning) e federado (federated learning)\end{enumerate}}\\
\multicolumn{4}{|p{\dimexpr 16cm + 6\tabcolsep\relax}|}{}\\
\hline
\multicolumn{4}{|p{16cm}|}{Bibliografia Básica:}\\*
\multicolumn{4}{|p{\dimexpr 16cm + 6\tabcolsep\relax}|}{\begin{itemize}\item SILBERSCHATZ, Abraham; GALVIN, Peter B.; GAGNE, Greg. Fundamentos de sistemas operacionais. Nova edição. São Paulo: LTC, 2015. ISBN 978-85-216-2321-2. Disponível em: <https://integrada.minhabiblioteca.com.br/\#/books/978-85-216-3001-2>. Acesso em: jun. 2025.\end{itemize}}\\
\multicolumn{4}{|p{\dimexpr 16cm + 6\tabcolsep\relax}|}{\begin{itemize}\item TANENBAUM, Andrew S.; STEEN, Maarten van. Sistemas distribuídos: princípios e paradigmas. 2. ed. São Paulo: Pearson Prentice Hall, 2007. Disponível em: <https://plataforma.bvirtual.com.br/Acervo/Publicacao/411>. Acesso em: jun. 2025.\end{itemize}}\\
\multicolumn{4}{|p{\dimexpr 16cm + 6\tabcolsep\relax}|}{\begin{itemize}\item GRAMA, Ananth; GUPTA, Anshul; KARYPIS, George; KUMAR, Vipin. Introduction to parallel computing. 2. ed. Upper Saddle River: Pearson, 2003.\end{itemize}}\\
\multicolumn{4}{|p{16cm}|}{}\\
\hline
\multicolumn{4}{|p{16cm}|}{Bibliografia Complementar:}\\*
\multicolumn{4}{|p{\dimexpr 16cm + 6\tabcolsep\relax}|}{\begin{itemize}\item TANENBAUM, Andrew S. Sistemas operacionais modernos. 4. ed. São Paulo: Pearson, 2009. ISBN 978-85-430-0567-6. Disponível em: <https://plataforma.bvirtual.com.br/Acervo/Publicacao/36876>. Acesso em: jun. 2025.\end{itemize}}\\
\multicolumn{4}{|p{\dimexpr 16cm + 6\tabcolsep\relax}|}{\begin{itemize}\item DEAN, Jeffrey; GHEMAWAT, Sanjay. MapReduce: simplified data processing on large clusters. Communications of the ACM, New York, v. 51, n. 1, p. 107–113, 2008.\end{itemize}}\\
\multicolumn{4}{|p{\dimexpr 16cm + 6\tabcolsep\relax}|}{\begin{itemize}\item PACHECO, Peter S. An introduction to parallel programming. 2. ed. Waltham: Morgan Kaufmann, 2011.\end{itemize}}\\
\multicolumn{4}{|p{\dimexpr 16cm + 6\tabcolsep\relax}|}{\begin{itemize}\item HENNESSY, John L.; PATTERSON, David A. Computer architecture: a quantitative approach. 4. ed. Amsterdam: Morgan Kaufmann, 2007. ISBN 978-0-12-370490-0.\end{itemize}}\\
\multicolumn{4}{|p{\dimexpr 16cm + 6\tabcolsep\relax}|}{\begin{itemize}\item COULOURIS, George; DOLLIMORE, Jean; KINDBERG, Tim; BLAIR, Gordon. Sistemas distribuídos: conceitos e projeto. Tradução de Alexandre Carissimi e João Tortello. 5. ed. Porto Alegre: Bookman, 2013. Disponível em: <https://integrada.minhabiblioteca.com.br/\#/books/9788582600542>. Acesso em: jun. 2025.\end{itemize}}\\
\hline
\end{longtable}
\end{center}

\clearpage
\begin{center}
\begin{tabular}{lccr}
 \multirow{3}{*}{\includegraphics[height=2.7cm]{brasao.png}} &
 \multicolumn{2}{c}{\bfseries UNIVERSIDADE FEDERAL DE OURO PRETO} &
 \ \ \ \ \multirow{3}{*}{\includegraphics[height=2.7cm]{ufop.png}} \\
 & \multicolumn{2}{c}{\bfseries PRÓ-REITORIA DE GRADUAÇÃO} & \\
 & \multicolumn{2}{c}{} & \\
 & \multicolumn{2}{c}{\Large\bfseries PROGRAMA DE DISCIPLINA} & \\
\end{tabular}
\end{center}

\begin{center}
\begin{longtable}{|p{4cm}|p{4cm}|p{4cm}|p{4cm}|}
\hline
\multicolumn{3}{|p{12cm}|}{Nome do Componente Curricular em Português:} &
\multicolumn{1}{p{4cm}|}{Código:} \\ 
\multicolumn{3}{|p{12cm}|}{\textbf{Processamento de Linguagem Natural}} &
\textbf{BIA110}\\ 
\multicolumn{3}{|p{12cm}|}{Nome do Componente Curricular em Inglês:} & \\ 
\multicolumn{3}{|p{12cm}|}{\textbf{Natural Language Processing}} & \\ 
\hline
\multicolumn{3}{|p{12cm}|}{Nome e Sigla do Departamento} & Unidade Acadêmica: \\ 
\multicolumn{3}{|p{12cm}|}{Departamento de Computação (DECOM)} & {ICEB} \\ 
\hline
\multicolumn{4}{|p{16cm}|}{Modalidade de Oferta:
[X] presencial \hspace{0.6cm}
[ ] semipresencial \hspace{0.6cm}
[ ] à distância}\\
\hline
\multicolumn{2}{|p{8cm}|}{Carga horária semestral} &
\multicolumn{2}{p{8cm}|}{Carga horária semanal}\\
\hline
\multicolumn{1}{|p{4cm}|}{Total} &
\multicolumn{1}{p{4cm}|}{Extensionista} &
\multicolumn{1}{p{4cm}|}{Teórica} &
\multicolumn{1}{p{4cm}|}{Prática} \\ 
\multicolumn{1}{|p{4cm}|}{60\,horas} &
\multicolumn{1}{p{4cm}|}{0\;horas} &
\multicolumn{1}{p{4cm}|}{4\;horas/aula} &
\multicolumn{1}{p{4cm}|}{0\;horas/aula} \\ 
\hline
\multicolumn{4}{|p{16cm}|}{Ementa:}\\
\multicolumn{4}{|p{16cm}|}{}\\
\multicolumn{4}{|p{\dimexpr 16cm + 6\tabcolsep\relax}|}{Estudo de técnicas avançadas em Processamento de Linguagem Natural (PLN) e modelos de linguagem de larga escala (LLMs). Arquiteturas baseadas em transformadores, pré-treinamento, fine-tuning, geração de texto, aplicações e desafios éticos.}\\
\multicolumn{4}{|p{16cm}|}{}\\
\hline
\multicolumn{4}{|p{16cm}|}{Conteúdo programático:}\\*
\multicolumn{4}{|p{\dimexpr 16cm + 6\tabcolsep\relax}|}{\begin{enumerate}[start=1]\item Introdução ao Processamento de Linguagem Natural: história e desafios.\end{enumerate}}\\
\multicolumn{4}{|p{\dimexpr 16cm + 6\tabcolsep\relax}|}{\begin{enumerate}[start=2]\item Representação de texto: embeddings tradicionais (TF-IDF, word2vec, GloVe) e contextualizados (ELMo, BERT).\end{enumerate}}\\
\multicolumn{4}{|p{\dimexpr 16cm + 6\tabcolsep\relax}|}{\begin{enumerate}[start=3]\item Modelos de linguagem clássicos: n-gramas, Markov e RNNs.\end{enumerate}}\\
\multicolumn{4}{|p{\dimexpr 16cm + 6\tabcolsep\relax}|}{\begin{enumerate}[start=4]\item Arquitetura Transformer: atenção, codificadores e decodificadores.\end{enumerate}}\\
\multicolumn{4}{|p{\dimexpr 16cm + 6\tabcolsep\relax}|}{\begin{enumerate}[start=5]\item Modelos de linguagem de larga escala: GPT, BERT, T5 e variantes.\end{enumerate}}\\
\multicolumn{4}{|p{\dimexpr 16cm + 6\tabcolsep\relax}|}{\begin{enumerate}[start=6]\item Pré-treinamento e fine-tuning: técnicas, datasets e estratégias.\end{enumerate}}\\
\multicolumn{4}{|p{\dimexpr 16cm + 6\tabcolsep\relax}|}{\begin{enumerate}[start=7]\item Avaliação de modelos de linguagem: perplexidade, BLEU, ROUGE e outras métricas.\end{enumerate}}\\
\multicolumn{4}{|p{\dimexpr 16cm + 6\tabcolsep\relax}|}{\begin{enumerate}[start=8]\item Aplicações práticas: chatbots, tradução automática, sumarização e análise de sentimento.\end{enumerate}}\\
\multicolumn{4}{|p{\dimexpr 16cm + 6\tabcolsep\relax}|}{\begin{enumerate}[start=9]\item Desafios e limitações: vieses, ética, consumo computacional e explicabilidade.\end{enumerate}}\\
\multicolumn{4}{|p{\dimexpr 16cm + 6\tabcolsep\relax}|}{\begin{enumerate}[start=10]\item Tendências atuais e pesquisas emergentes em PLN e LLMs.\end{enumerate}}\\
\multicolumn{4}{|p{\dimexpr 16cm + 6\tabcolsep\relax}|}{\begin{enumerate}[start=11]\item Ética e Vieses e Limites de LLMs\end{enumerate}}\\
\multicolumn{4}{|p{\dimexpr 16cm + 6\tabcolsep\relax}|}{}\\
\hline
\multicolumn{4}{|p{16cm}|}{Bibliografia Básica:}\\*
\multicolumn{4}{|p{\dimexpr 16cm + 6\tabcolsep\relax}|}{\begin{itemize}\item JURAFSKY, Daniel; MARTIN, James H. Speech and Language Processing. 3. ed. Draft. 2023. Disponível em: https://web.stanford.edu/~jurafsky/slp3/. Acesso em: jun/2025.\end{itemize}}\\
\multicolumn{4}{|p{\dimexpr 16cm + 6\tabcolsep\relax}|}{\begin{itemize}\item GOODFELLOW, Ian; BENGIO, Yoshua; COURVILLE, Aaron. Deep Learning. Cambridge: MIT Press, 2016.\end{itemize}}\\
\multicolumn{4}{|p{\dimexpr 16cm + 6\tabcolsep\relax}|}{\begin{itemize}\item CASELI, Helena de Medeiros; NUNES, Maria das Graças Volpe (org.). Processamento de linguagem natural: conceitos, técnicas e aplicações em português. 3. ed. São Carlos: BPLN, 2024. Disponível em: <https://brasileiraspln.com/livro-pln>. Acesso em: 22 dez. 2025.\end{itemize}}\\
\multicolumn{4}{|p{\dimexpr 16cm + 6\tabcolsep\relax}|}{\begin{itemize}\item EISENSTEIN, Jacob. Introduction to natural language processing. Cambridge; London: The MIT Press, 2019. ISBN 978-0-262-04284-0. 536 p.\end{itemize}}\\
\multicolumn{4}{|p{\dimexpr 16cm + 6\tabcolsep\relax}|}{\begin{itemize}\item MARTINS, Júlio Serafim; LENZ, Maikon Lucian; SILVA, Michel Bernardo Fernandes da; et al. Processamentos de linguagem natural. Porto Alegre: Grupo A, 2020. E-book. Disponível em: <https://integrada.minhabiblioteca.com.br/reader/books/9786556900575>. Acesso em: 22 dez. 2025.\end{itemize}}\\
\multicolumn{4}{|p{16cm}|}{}\\
\hline
\multicolumn{4}{|p{16cm}|}{Bibliografia Complementar:}\\*
\multicolumn{4}{|p{\dimexpr 16cm + 6\tabcolsep\relax}|}{\begin{itemize}\item VASWANI, Ashish et al. Attention is All You Need. In: Advances in Neural Information Processing Systems (NeurIPS), 2017.\end{itemize}}\\
\multicolumn{4}{|p{\dimexpr 16cm + 6\tabcolsep\relax}|}{\begin{itemize}\item BENDER, Emily M.; GEBRU, Timnit; McMILLAN-MAJOR, Angelina; SHMITCHELL, Shmargaret. On the dangers of stochastic parrots: can language models be too big? In: ACM CONFERENCE ON FAIRNESS, ACCOUNTABILITY, AND TRANSPARENCY (FAccT), 2021, New York. Proceedings… New York: Association for Computing Machinery, 2021. p. 610–623. DOI: https://doi.org/10.1145/3442188.3445922.\end{itemize}}\\
\multicolumn{4}{|p{\dimexpr 16cm + 6\tabcolsep\relax}|}{\begin{itemize}\item RADFORD, Alec et al. Language models are unsupervised multitask learners. OpenAI Technical Report, 2019.\end{itemize}}\\
\multicolumn{4}{|p{\dimexpr 16cm + 6\tabcolsep\relax}|}{\begin{itemize}\item DEVLIN, Jacob et al. BERT: Pre-training of Deep Bidirectional Transformers for Language Understanding. 2019. Disponível em: https://arxiv.org/abs/1810.04805.\end{itemize}}\\
\multicolumn{4}{|p{\dimexpr 16cm + 6\tabcolsep\relax}|}{\begin{itemize}\item CHOWDHURY, Gaurav. Natural Language Processing. 1. ed. Springer, 2021.\end{itemize}}\\
\multicolumn{4}{|p{\dimexpr 16cm + 6\tabcolsep\relax}|}{\begin{itemize}\item MANNING, Christopher D.; SCHÜTZE, Hinrich. Foundations of Statistical Natural Language Processing. Cambridge: MIT Press, 1999.\end{itemize}}\\
\multicolumn{4}{|p{\dimexpr 16cm + 6\tabcolsep\relax}|}{\begin{itemize}\item BROWN, Tom B. et al. Language models are few-shot learners. In: Advances in Neural Information Processing Systems (NeurIPS), v. 33, 2020. Disponível em: https://arxiv.org/abs/2005.14165.\end{itemize}}\\
\hline
\end{longtable}
\end{center}

\clearpage
\begin{center}
\begin{tabular}{lccr}
 \multirow{3}{*}{\includegraphics[height=2.7cm]{brasao.png}} &
 \multicolumn{2}{c}{\bfseries UNIVERSIDADE FEDERAL DE OURO PRETO} &
 \ \ \ \ \multirow{3}{*}{\includegraphics[height=2.7cm]{ufop.png}} \\
 & \multicolumn{2}{c}{\bfseries PRÓ-REITORIA DE GRADUAÇÃO} & \\
 & \multicolumn{2}{c}{} & \\
 & \multicolumn{2}{c}{\Large\bfseries PROGRAMA DE DISCIPLINA} & \\
\end{tabular}
\end{center}

\begin{center}
\begin{longtable}{|p{4cm}|p{4cm}|p{4cm}|p{4cm}|}
\hline
\multicolumn{3}{|p{12cm}|}{Nome do Componente Curricular em Português:} &
\multicolumn{1}{p{4cm}|}{Código:} \\ 
\multicolumn{3}{|p{12cm}|}{\textbf{Visão Computacional e Processamento de Imagem}} &
\textbf{BIA111}\\ 
\multicolumn{3}{|p{12cm}|}{Nome do Componente Curricular em Inglês:} & \\ 
\multicolumn{3}{|p{12cm}|}{\textbf{Computer Vision and Image Processing}} & \\ 
\hline
\multicolumn{3}{|p{12cm}|}{Nome e Sigla do Departamento} & Unidade Acadêmica: \\ 
\multicolumn{3}{|p{12cm}|}{Departamento de Computação (DECOM)} & {ICEB} \\ 
\hline
\multicolumn{4}{|p{16cm}|}{Modalidade de Oferta:
[X] presencial \hspace{0.6cm}
[ ] semipresencial \hspace{0.6cm}
[ ] à distância}\\
\hline
\multicolumn{2}{|p{8cm}|}{Carga horária semestral} &
\multicolumn{2}{p{8cm}|}{Carga horária semanal}\\
\hline
\multicolumn{1}{|p{4cm}|}{Total} &
\multicolumn{1}{p{4cm}|}{Extensionista} &
\multicolumn{1}{p{4cm}|}{Teórica} &
\multicolumn{1}{p{4cm}|}{Prática} \\ 
\multicolumn{1}{|p{4cm}|}{60\,horas} &
\multicolumn{1}{p{4cm}|}{0\;horas} &
\multicolumn{1}{p{4cm}|}{4\;horas/aula} &
\multicolumn{1}{p{4cm}|}{0\;horas/aula} \\ 
\hline
\multicolumn{4}{|p{16cm}|}{Ementa:}\\
\multicolumn{4}{|p{16cm}|}{}\\
\multicolumn{4}{|p{\dimexpr 16cm + 6\tabcolsep\relax}|}{Introdução aos fundamentos da visão computacional e processamento de imagem digital. Técnicas para análise, transformação e interpretação de imagens. Aplicações em reconhecimento, segmentação e análise visual.}\\
\multicolumn{4}{|p{16cm}|}{}\\
\hline
\multicolumn{4}{|p{16cm}|}{Conteúdo programático:}\\*
\multicolumn{4}{|p{\dimexpr 16cm + 6\tabcolsep\relax}|}{\begin{enumerate}[start=1]\item Fundamentos da formação de imagens e percepção visual.\end{enumerate}}\\
\multicolumn{4}{|p{\dimexpr 16cm + 6\tabcolsep\relax}|}{\begin{enumerate}[start=2]\item Representação digital de imagens: pixels, cores e formatos.\end{enumerate}}\\
\multicolumn{4}{|p{\dimexpr 16cm + 6\tabcolsep\relax}|}{\begin{enumerate}[start=3]\item Operações básicas de processamento de imagens: filtragem, realce e transformação.\end{enumerate}}\\
\multicolumn{4}{|p{\dimexpr 16cm + 6\tabcolsep\relax}|}{\begin{enumerate}[start=4]\item Detecção de bordas, contornos e extração de características.\end{enumerate}}\\
\multicolumn{4}{|p{\dimexpr 16cm + 6\tabcolsep\relax}|}{\begin{enumerate}[start=5]\item Segmentação de imagens: métodos baseados em limiarização, region growing e clustering.\end{enumerate}}\\
\multicolumn{4}{|p{\dimexpr 16cm + 6\tabcolsep\relax}|}{\begin{enumerate}[start=6]\item Morfologia matemática e processamento de formas.\end{enumerate}}\\
\multicolumn{4}{|p{\dimexpr 16cm + 6\tabcolsep\relax}|}{\begin{enumerate}[start=7]\item Reconhecimento de padrões e classificação de imagens.\end{enumerate}}\\
\multicolumn{4}{|p{\dimexpr 16cm + 6\tabcolsep\relax}|}{\begin{enumerate}[start=8]\item Visão computacional baseada em aprendizado: redes neurais convolucionais (CNNs).\end{enumerate}}\\
\multicolumn{4}{|p{\dimexpr 16cm + 6\tabcolsep\relax}|}{\begin{enumerate}[start=9]\item Detecção e rastreamento de objetos em vídeo.\end{enumerate}}\\
\multicolumn{4}{|p{\dimexpr 16cm + 6\tabcolsep\relax}|}{\begin{enumerate}[start=10]\item Aplicações práticas: reconhecimento facial, visão para robótica, análise médica.\end{enumerate}}\\
\multicolumn{4}{|p{\dimexpr 16cm + 6\tabcolsep\relax}|}{}\\
\hline
\multicolumn{4}{|p{16cm}|}{Bibliografia Básica:}\\*
\multicolumn{4}{|p{\dimexpr 16cm + 6\tabcolsep\relax}|}{\begin{itemize}\item GONZALEZ, Rafael C.; WOODS, Richard E. Digital image processing. 3. ed. São Paulo: Pearson, 2009. E-book. Disponível em: https://plataforma.bvirtual.com.br. Acesso em: 22 dez. 2025.\end{itemize}}\\
\multicolumn{4}{|p{\dimexpr 16cm + 6\tabcolsep\relax}|}{\begin{itemize}\item BARELLI, Felipe. Introdução à visão computacional: uma abordagem prática com Python e OpenCV. São Paulo: Casa do Código, 2018. E-book. Disponível em: https://plataforma.bvirtual.com.br. Acesso em: 22 dez. 2025.\end{itemize}}\\
\multicolumn{4}{|p{\dimexpr 16cm + 6\tabcolsep\relax}|}{\begin{itemize}\item SZELISKI, Richard. Computer vision: algorithms and applications. 2. ed. Cham: Springer, 2022. E-book. DOI: https://doi.org/10.1007/978-3-030-34372-9.\end{itemize}}\\
\multicolumn{4}{|p{16cm}|}{}\\
\hline
\multicolumn{4}{|p{16cm}|}{Bibliografia Complementar:}\\*
\multicolumn{4}{|p{\dimexpr 16cm + 6\tabcolsep\relax}|}{\begin{itemize}\item SZAJDAK, Marek. Introdução à visão computacional. 1. ed. Rio de Janeiro: LTC, 2021.\end{itemize}}\\
\multicolumn{4}{|p{\dimexpr 16cm + 6\tabcolsep\relax}|}{\begin{itemize}\item GOODFELLOW, Ian; BENGIO, Yoshua; COURVILLE, Aaron. Deep learning. Cambridge: MIT Press, 2016.\end{itemize}}\\
\multicolumn{4}{|p{\dimexpr 16cm + 6\tabcolsep\relax}|}{\begin{itemize}\item RUSSELL, Stuart; NORVIG, Peter. Inteligência artificial. 3. ed. Rio de Janeiro: Elsevier, 2013.\end{itemize}}\\
\multicolumn{4}{|p{\dimexpr 16cm + 6\tabcolsep\relax}|}{\begin{itemize}\item FORSYTH, David A.; PONCE, Jean. Computer vision: a modern approach. 2. ed. Boston: Pearson, 2011.\end{itemize}}\\
\multicolumn{4}{|p{\dimexpr 16cm + 6\tabcolsep\relax}|}{\begin{itemize}\item BALLARD, Dana H.; BROWN, Christopher M. Computer vision. Englewood Cliffs: Prentice Hall, 1982.\end{itemize}}\\
\multicolumn{4}{|p{\dimexpr 16cm + 6\tabcolsep\relax}|}{\begin{itemize}\item VOULODIMOS, Athanasios; DOULAMIS, Nikolaos; DOULAMIS, Anastasios; PROTOPAPADAKIS, Eftychios. Deep learning for computer vision: a brief review. Computational Intelligence and Neuroscience, v. 2018, art. 7068349, 2018. DOI: 10.1155/2018/7068349.\end{itemize}}\\
\hline
\end{longtable}
\end{center}

\clearpage
\begin{center}
\begin{tabular}{lccr}
 \multirow{3}{*}{\includegraphics[height=2.7cm]{brasao.png}} &
 \multicolumn{2}{c}{\bfseries UNIVERSIDADE FEDERAL DE OURO PRETO} &
 \ \ \ \ \multirow{3}{*}{\includegraphics[height=2.7cm]{ufop.png}} \\
 & \multicolumn{2}{c}{\bfseries PRÓ-REITORIA DE GRADUAÇÃO} & \\
 & \multicolumn{2}{c}{} & \\
 & \multicolumn{2}{c}{\Large\bfseries PROGRAMA DE DISCIPLINA} & \\
\end{tabular}
\end{center}

\begin{center}
\begin{longtable}{|p{4cm}|p{4cm}|p{4cm}|p{4cm}|}
\hline
\multicolumn{3}{|p{12cm}|}{Nome do Componente Curricular em Português:} &
\multicolumn{1}{p{4cm}|}{Código:} \\ 
\multicolumn{3}{|p{12cm}|}{\textbf{Metodologia Científica para Inteligência Artificial}} &
\textbf{BIA112}\\ 
\multicolumn{3}{|p{12cm}|}{Nome do Componente Curricular em Inglês:} & \\ 
\multicolumn{3}{|p{12cm}|}{\textbf{Scientific Methodology for Artificial Intelligence}} & \\ 
\hline
\multicolumn{3}{|p{12cm}|}{Nome e Sigla do Departamento} & Unidade Acadêmica: \\ 
\multicolumn{3}{|p{12cm}|}{Departamento de Computação (DECOM)} & {ICEB} \\ 
\hline
\multicolumn{4}{|p{16cm}|}{Modalidade de Oferta:
[X] presencial \hspace{0.6cm}
[ ] semipresencial \hspace{0.6cm}
[ ] à distância}\\
\hline
\multicolumn{2}{|p{8cm}|}{Carga horária semestral} &
\multicolumn{2}{p{8cm}|}{Carga horária semanal}\\
\hline
\multicolumn{1}{|p{4cm}|}{Total} &
\multicolumn{1}{p{4cm}|}{Extensionista} &
\multicolumn{1}{p{4cm}|}{Teórica} &
\multicolumn{1}{p{4cm}|}{Prática} \\ 
\multicolumn{1}{|p{4cm}|}{30\,horas} &
\multicolumn{1}{p{4cm}|}{0\;horas} &
\multicolumn{1}{p{4cm}|}{2\;horas/aula} &
\multicolumn{1}{p{4cm}|}{0\;horas/aula} \\ 
\hline
\multicolumn{4}{|p{16cm}|}{Ementa:}\\
\multicolumn{4}{|p{16cm}|}{}\\
\multicolumn{4}{|p{\dimexpr 16cm + 6\tabcolsep\relax}|}{Estudo dos fundamentos epistemológicos da ciência e de seus métodos, com ênfase na filosofia da ciência. Análise dos elementos da pesquisa científica em ciência da computação, incluindo formulação de problemas, hipóteses e desenho experimental. Abordagem da metodologia científica própria da pesquisa em inteligência artificial, com foco em avaliação empírica, validade, reprodutibilidade, interpretação de resultados experimentais e ética}\\
\multicolumn{4}{|p{16cm}|}{}\\
\hline
\multicolumn{4}{|p{16cm}|}{Conteúdo programático:}\\*
\multicolumn{4}{|p{\dimexpr 16cm + 6\tabcolsep\relax}|}{\begin{enumerate}[start=1]\item Ciência e Demarcação Científica: critérios de cientificidade, distinção entre ciência, pseudociência e tecnologia, e limites do conhecimento científico.\end{enumerate}}\\
\multicolumn{4}{|p{\dimexpr 16cm + 6\tabcolsep\relax}|}{\begin{enumerate}[start=2]\item Observação, Indução e Teoria: papel da observação na ciência, crítica ao indutivismo e influência das teorias na produção do conhecimento científico.\end{enumerate}}\\
\multicolumn{4}{|p{\dimexpr 16cm + 6\tabcolsep\relax}|}{\begin{enumerate}[start=3]\item Progresso Científico e Mudança de Paradigmas: falsificacionismo, paradigmas científicos, programas de pesquisa e pluralismo metodológico.\end{enumerate}}\\
\multicolumn{4}{|p{\dimexpr 16cm + 6\tabcolsep\relax}|}{\begin{enumerate}[start=4]\item Pesquisa em Ciência da Computação e Inteligência Artificial: natureza da pesquisa em computação, distinção entre pesquisa científica, experimental e desenvolvimento tecnológico.\end{enumerate}}\\
\multicolumn{4}{|p{\dimexpr 16cm + 6\tabcolsep\relax}|}{\begin{enumerate}[start=5]\item Busca e avaliação de fontes bibliográficas: bases de dados, periódicos e conferências em IA.\end{enumerate}}\\
\multicolumn{4}{|p{\dimexpr 16cm + 6\tabcolsep\relax}|}{\begin{enumerate}[start=6]\item Formulação de Problemas e Hipóteses: definição de problemas científicos, abstrações, modelos, hipóteses e objetivos de pesquisa.\end{enumerate}}\\
\multicolumn{4}{|p{\dimexpr 16cm + 6\tabcolsep\relax}|}{\begin{enumerate}[start=7]\item Pesquisa Empírica em Inteligência Artificial: características da investigação empírica em IA e aprendizado de máquina, estudos exploratórios e confirmatórios.\end{enumerate}}\\
\multicolumn{4}{|p{\dimexpr 16cm + 6\tabcolsep\relax}|}{\begin{enumerate}[start=8]\item Desenho Experimental e Avaliação em Inteligência Artificial: experimentos controlados, comparação de métodos, métricas e interpretação de resultados.\end{enumerate}}\\
\multicolumn{4}{|p{\dimexpr 16cm + 6\tabcolsep\relax}|}{\begin{enumerate}[start=9]\item Rigor Científico, Validade e Reprodutibilidade em IA: validade interna e externa, vieses experimentais, reprodutibilidade e boas práticas em pesquisa em IA.\end{enumerate}}\\
\multicolumn{4}{|p{\dimexpr 16cm + 6\tabcolsep\relax}|}{\begin{enumerate}[start=10]\item Avaliação Empírica e Inferência em Sistemas de IA: métricas de desempenho, testes estatísticos, comparação de modelos e interpretação científica de evidências.\end{enumerate}}\\
\multicolumn{4}{|p{\dimexpr 16cm + 6\tabcolsep\relax}|}{\begin{enumerate}[start=11]\item Ética em pesquisa científica: plágio, integridade, viés algorítmico e responsabilidade social.\end{enumerate}}\\
\multicolumn{4}{|p{\dimexpr 16cm + 6\tabcolsep\relax}|}{}\\
\hline
\multicolumn{4}{|p{16cm}|}{Bibliografia Básica:}\\*
\multicolumn{4}{|p{\dimexpr 16cm + 6\tabcolsep\relax}|}{\begin{itemize}\item CHALMERS, Alan F. O que é ciência, afinal? São Paulo: Brasiliense, 1993.\end{itemize}}\\
\multicolumn{4}{|p{\dimexpr 16cm + 6\tabcolsep\relax}|}{\begin{itemize}\item WAZLAWICK, Raul Sidnei. Metodologia de pesquisa para ciência da computação. 2. ed. Rio de Janeiro: Elsevier, 2014.\end{itemize}}\\
\multicolumn{4}{|p{\dimexpr 16cm + 6\tabcolsep\relax}|}{\begin{itemize}\item PINEAU, Joelle et al. Improving reproducibility in machine learning research. Journal of Machine Learning Research, v. 22, n. 164, p. 1–20, 2021.\end{itemize}}\\
\multicolumn{4}{|p{\dimexpr 16cm + 6\tabcolsep\relax}|}{\begin{itemize}\item RAINIO, O.; TEUHO, J.; KLÉN, R. Evaluation metrics and statistical tests for machine learning. Scientific Reports, London, v. 14, p. 6086, 2024. DOI: 10.1038/s41598-024-56706-x. Disponível em: <https://doi.org/10.1038/s41598-024-56706-x>. Acesso em: dez. 2025.\end{itemize}}\\
\multicolumn{4}{|p{16cm}|}{}\\
\hline
\multicolumn{4}{|p{16cm}|}{Bibliografia Complementar:}\\*
\multicolumn{4}{|p{\dimexpr 16cm + 6\tabcolsep\relax}|}{\begin{itemize}\item NEWELL, Allen; SIMON, Herbert A. Computer science as empirical inquiry: symbols and search. Communications of the ACM, v. 19, n. 3, p. 113–126, 1976.\end{itemize}}\\
\multicolumn{4}{|p{\dimexpr 16cm + 6\tabcolsep\relax}|}{\begin{itemize}\item COHEN, Paul R. Empirical methods for artificial intelligence. AI Magazine, v. 16, n. 1, p. 63–73, 1995.\end{itemize}}\\
\multicolumn{4}{|p{\dimexpr 16cm + 6\tabcolsep\relax}|}{\begin{itemize}\item KAPLAN, David M. (ed.). Readings in the philosophy of technology. 2. ed. Lanham, MD: Rowman \& Littlefield, 2009.\end{itemize}}\\
\multicolumn{4}{|p{\dimexpr 16cm + 6\tabcolsep\relax}|}{\begin{itemize}\item ASSOCIAÇÃO BRASILEIRA DE NORMAS TÉCNICAS. NBR 6023, NBR 10520 e demais normas para trabalhos acadêmicos. Disponível em: https://www.abnt.org.br. Acesso em: jun/2025.\end{itemize}}\\
\multicolumn{4}{|p{\dimexpr 16cm + 6\tabcolsep\relax}|}{\begin{itemize}\item LIPTON, Zachary C.; STEINHARDT, Jacob. Troubling trends in machine learning scholarship. arXiv preprint, arXiv:1807.03341, 2018.\end{itemize}}\\
\multicolumn{4}{|p{\dimexpr 16cm + 6\tabcolsep\relax}|}{\begin{itemize}\item OLTEANU, Alexandra; BLODGETT, Su Lin; BALAYN, Agathe; WANG, Angelina; DIAZ, Fernando; CALMON, Flavio du Pin; MITCHELL, Margaret; EKSTRAND, Michael; BINNS, Reuben; BAROCAS, Solon. Rigor in AI: doing rigorous AI work requires a broader, responsible AI-informed conception of rigor. arXiv preprint, arXiv:2506.14652, 2025. Disponível em: <https://arxiv.org/abs/2506.14652>. Acesso em: dez. 2025.\end{itemize}}\\
\multicolumn{4}{|p{\dimexpr 16cm + 6\tabcolsep\relax}|}{\begin{itemize}\item HERRMANN, Moritz; LANGE, F. Julian D.; EGGENSPERGER, Katharina; CASALICCHIO, Giuseppe; WEVER, Marcel; FEURER, Matthias; RÜGAMER, David; HÜLLERMEIER, Eyke; BOULESTEIX, Anne-Laure; BISCHL, Bernd. Position: why we must rethink empirical research in machine learning. arXiv preprint, arXiv:2405.02200, 2024. Disponível em: <https://arxiv.org/abs/2405.02200>. Acesso em: dez. 2025.\end{itemize}}\\
\hline
\end{longtable}
\end{center}

\clearpage
\begin{center}
\begin{tabular}{lccr}
 \multirow{3}{*}{\includegraphics[height=2.7cm]{brasao.png}} &
 \multicolumn{2}{c}{\bfseries UNIVERSIDADE FEDERAL DE OURO PRETO} &
 \ \ \ \ \multirow{3}{*}{\includegraphics[height=2.7cm]{ufop.png}} \\
 & \multicolumn{2}{c}{\bfseries PRÓ-REITORIA DE GRADUAÇÃO} & \\
 & \multicolumn{2}{c}{} & \\
 & \multicolumn{2}{c}{\Large\bfseries PROGRAMA DE DISCIPLINA} & \\
\end{tabular}
\end{center}

\begin{center}
\begin{longtable}{|p{4cm}|p{4cm}|p{4cm}|p{4cm}|}
\hline
\multicolumn{3}{|p{12cm}|}{Nome do Componente Curricular em Português:} &
\multicolumn{1}{p{4cm}|}{Código:} \\ 
\multicolumn{3}{|p{12cm}|}{\textbf{Segurança, Ética e Sociedade em Inteligência Artificial}} &
\textbf{BIA113}\\ 
\multicolumn{3}{|p{12cm}|}{Nome do Componente Curricular em Inglês:} & \\ 
\multicolumn{3}{|p{12cm}|}{\textbf{AI Ethics, Safety and Risk}} & \\ 
\hline
\multicolumn{3}{|p{12cm}|}{Nome e Sigla do Departamento} & Unidade Acadêmica: \\ 
\multicolumn{3}{|p{12cm}|}{Departamento de Computação (DECOM)} & {ICEB} \\ 
\hline
\multicolumn{4}{|p{16cm}|}{Modalidade de Oferta:
[X] presencial \hspace{0.6cm}
[ ] semipresencial \hspace{0.6cm}
[ ] à distância}\\
\hline
\multicolumn{2}{|p{8cm}|}{Carga horária semestral} &
\multicolumn{2}{p{8cm}|}{Carga horária semanal}\\
\hline
\multicolumn{1}{|p{4cm}|}{Total} &
\multicolumn{1}{p{4cm}|}{Extensionista} &
\multicolumn{1}{p{4cm}|}{Teórica} &
\multicolumn{1}{p{4cm}|}{Prática} \\ 
\multicolumn{1}{|p{4cm}|}{60\,horas} &
\multicolumn{1}{p{4cm}|}{0\;horas} &
\multicolumn{1}{p{4cm}|}{4\;horas/aula} &
\multicolumn{1}{p{4cm}|}{0\;horas/aula} \\ 
\hline
\multicolumn{4}{|p{16cm}|}{Ementa:}\\
\multicolumn{4}{|p{16cm}|}{}\\
\multicolumn{4}{|p{\dimexpr 16cm + 6\tabcolsep\relax}|}{Estudo dos riscos associados a sistemas de Inteligência Artificial e dos fundamentos técnicos que sustentam tais riscos. Análise de riscos catastróficos, uso malicioso e dinâmicas competitivas no desenvolvimento de IA. Fundamentos de Inteligência Artificial relevantes para segurança, incluindo aprendizado de máquina, deep learning, leis de escala e velocidade de desenvolvimento. Segurança de agentes individuais, abordando monitoramento, robustez, alinhamento e capacidades emergentes. Engenharia de segurança aplicada a sistemas de IA, incluindo confiabilidade, decomposição de riscos, falhas de componentes, fatores sistêmicos e eventos extremos. Análise da IA como sistema complexo e sociotécnico.}\\
\multicolumn{4}{|p{16cm}|}{}\\
\hline
\multicolumn{4}{|p{16cm}|}{Conteúdo programático:}\\*
\multicolumn{4}{|p{\dimexpr 16cm + 6\tabcolsep\relax}|}{\begin{enumerate}[start=1]\item Visão Geral dos Riscos Catastróficos da Inteligência Artificial: principais riscos existenciais e sistêmicos da IA, incluindo uso malicioso, corrida tecnológica e falhas organizacionais.\end{enumerate}}\\
\multicolumn{4}{|p{\dimexpr 16cm + 6\tabcolsep\relax}|}{\begin{enumerate}[start=2]\item Fundamentos de Inteligência Artificial: conceitos básicos de IA, aprendizado de máquina e aprendizado profundo, com ênfase em leis de escala e ritmo de desenvolvimento.\end{enumerate}}\\
\multicolumn{4}{|p{\dimexpr 16cm + 6\tabcolsep\relax}|}{\begin{enumerate}[start=3]\item Segurança em Sistemas de Agente Único: problemas de alinhamento, robustez, monitoramento e segurança em agentes individuais de IA.\end{enumerate}}\\
\multicolumn{4}{|p{\dimexpr 16cm + 6\tabcolsep\relax}|}{\begin{enumerate}[start=4]\item Engenharia de Segurança: princípios e métodos de engenharia de segurança aplicados a sistemas de IA, incluindo confiabilidade e análise de falhas.\end{enumerate}}\\
\multicolumn{4}{|p{\dimexpr 16cm + 6\tabcolsep\relax}|}{\begin{enumerate}[start=5]\item Sistemas Complexos: noções de sistemas complexos e comportamentos emergentes relevantes para a segurança em IA.\end{enumerate}}\\
\multicolumn{4}{|p{\dimexpr 16cm + 6\tabcolsep\relax}|}{\begin{enumerate}[start=6]\item Inteligência Artificial Benéfica e Ética das Máquinas: fundamentos éticos, legais e sociais para o desenvolvimento de sistemas de IA alinhados a valores humanos.\end{enumerate}}\\
\multicolumn{4}{|p{\dimexpr 16cm + 6\tabcolsep\relax}|}{\begin{enumerate}[start=7]\item Problemas de Ação Coletiva: desafios de cooperação, conflito e incentivos no desenvolvimento e uso da IA.\end{enumerate}}\\
\multicolumn{4}{|p{\dimexpr 16cm + 6\tabcolsep\relax}|}{\begin{enumerate}[start=8]\item Governança: estruturas e políticas de governança da IA em níveis corporativo, nacional e internacional.\end{enumerate}}\\
\multicolumn{4}{|p{\dimexpr 16cm + 6\tabcolsep\relax}|}{}\\
\hline
\multicolumn{4}{|p{16cm}|}{Bibliografia Básica:}\\*
\multicolumn{4}{|p{\dimexpr 16cm + 6\tabcolsep\relax}|}{\begin{itemize}\item HENDRYCKS, Dan. Introduction to AI Safety, Ethics and Society. Boca Raton: Taylor \& Francis, 2024. ISBN 978-1-032-79802-8. Disponível em: <https://www.aisafetybook.com>. Acesso em: 17 dez. 2025.\end{itemize}}\\
\multicolumn{4}{|p{\dimexpr 16cm + 6\tabcolsep\relax}|}{\begin{itemize}\item RUSSELL, Stuart. Inteligência artificial a nosso favor: como manter o controle sobre a tecnologia. Tradução de Berilo Vargas. São Paulo: Companhia das Letras, 2021.\end{itemize}}\\
\multicolumn{4}{|p{\dimexpr 16cm + 6\tabcolsep\relax}|}{\begin{itemize}\item AMODEI, Dario et al. Concrete problems in AI safety. arXiv preprint, arXiv:1606.06565, 2016. Disponível em: <https://arxiv.org/abs/1606.06565>. Acesso em: 17 dez. 2025.\end{itemize}}\\
\multicolumn{4}{|p{16cm}|}{}\\
\hline
\multicolumn{4}{|p{16cm}|}{Bibliografia Complementar:}\\*
\multicolumn{4}{|p{\dimexpr 16cm + 6\tabcolsep\relax}|}{\begin{itemize}\item BOSTROM, Nick. Superinteligência: caminhos, perigos e estratégias. Rio de Janeiro: DarkSide Books, 2018.\end{itemize}}\\
\multicolumn{4}{|p{\dimexpr 16cm + 6\tabcolsep\relax}|}{\begin{itemize}\item PERROW, Charles. Normal accidents: living with high-risk technologies. Updated ed. Princeton: Princeton University Press, 1999.\end{itemize}}\\
\multicolumn{4}{|p{\dimexpr 16cm + 6\tabcolsep\relax}|}{\begin{itemize}\item DOBBE, Roel I. J. System safety and artificial intelligence. arXiv preprint, arXiv:2202.09292, 2022. Disponível em: <https://arxiv.org/abs/2202.09292>.\end{itemize}}\\
\multicolumn{4}{|p{\dimexpr 16cm + 6\tabcolsep\relax}|}{\begin{itemize}\item CHRISTIAN, Brian. The alignment problem: machine learning and human values. New York: W. W. Norton, 2020.\end{itemize}}\\
\multicolumn{4}{|p{\dimexpr 16cm + 6\tabcolsep\relax}|}{\begin{itemize}\item BRASIL. Ministério da Ciência, Tecnologia e Inovações. Estratégia brasileira de inteligência artificial (EBIA). Brasília, DF: MCTI, 2021. Disponível em: <https://www.gov.br/mcti>. Acesso em: 17 dez. 2025.\end{itemize}}\\
\hline
\end{longtable}
\end{center}

\clearpage
\begin{center}
\begin{tabular}{lccr}
 \multirow{3}{*}{\includegraphics[height=2.7cm]{brasao.png}} &
 \multicolumn{2}{c}{\bfseries UNIVERSIDADE FEDERAL DE OURO PRETO} &
 \ \ \ \ \multirow{3}{*}{\includegraphics[height=2.7cm]{ufop.png}} \\
 & \multicolumn{2}{c}{\bfseries PRÓ-REITORIA DE GRADUAÇÃO} & \\
 & \multicolumn{2}{c}{} & \\
 & \multicolumn{2}{c}{\Large\bfseries PROGRAMA DE DISCIPLINA} & \\
\end{tabular}
\end{center}

\begin{center}
\begin{longtable}{|p{4cm}|p{4cm}|p{4cm}|p{4cm}|}
\hline
\multicolumn{3}{|p{12cm}|}{Nome do Componente Curricular em Português:} &
\multicolumn{1}{p{4cm}|}{Código:} \\ 
\multicolumn{3}{|p{12cm}|}{\textbf{Projeto 1: Resolução de Problemas do Mundo Real}} &
\textbf{BIA201}\\ 
\multicolumn{3}{|p{12cm}|}{Nome do Componente Curricular em Inglês:} & \\ 
\multicolumn{3}{|p{12cm}|}{\textbf{Project 1: Real-World Problem Solving}} & \\ 
\hline
\multicolumn{3}{|p{12cm}|}{Nome e Sigla do Departamento} & Unidade Acadêmica: \\ 
\multicolumn{3}{|p{12cm}|}{Departamento de Computação (DECOM)} & {ICEB} \\ 
\hline
\multicolumn{4}{|p{16cm}|}{Modalidade de Oferta:
[X] presencial \hspace{0.6cm}
[ ] semipresencial \hspace{0.6cm}
[ ] à distância}\\
\hline
\multicolumn{2}{|p{8cm}|}{Carga horária semestral} &
\multicolumn{2}{p{8cm}|}{Carga horária semanal}\\
\hline
\multicolumn{1}{|p{4cm}|}{Total} &
\multicolumn{1}{p{4cm}|}{Extensionista} &
\multicolumn{1}{p{4cm}|}{Teórica} &
\multicolumn{1}{p{4cm}|}{Prática} \\ 
\multicolumn{1}{|p{4cm}|}{120\,horas} &
\multicolumn{1}{p{4cm}|}{8\;horas} &
\multicolumn{1}{p{4cm}|}{0\;horas/aula} &
\multicolumn{1}{p{4cm}|}{0\;horas/aula} \\ 
\hline
\multicolumn{4}{|p{16cm}|}{Ementa:}\\
\multicolumn{4}{|p{16cm}|}{}\\
\multicolumn{4}{|p{\dimexpr 16cm + 6\tabcolsep\relax}|}{Desenvolvimento de projetos extensionistas voltados à identificação e resolução de problemas reais com aplicação de técnicas de inteligência artificial e computação. Enfoque em trabalho interdisciplinar, colaboração com a comunidade e inovação social.}\\
\multicolumn{4}{|p{16cm}|}{}\\
\hline
\multicolumn{4}{|p{16cm}|}{Conteúdo programático:}\\*
\multicolumn{4}{|p{\dimexpr 16cm + 6\tabcolsep\relax}|}{\begin{enumerate}[start=1]\item Identificação e definição de problemas reais em diferentes contextos sociais e econômicos.\end{enumerate}}\\
\multicolumn{4}{|p{\dimexpr 16cm + 6\tabcolsep\relax}|}{\begin{enumerate}[start=2]\item Levantamento e análise das necessidades dos stakeholders.\end{enumerate}}\\
\multicolumn{4}{|p{\dimexpr 16cm + 6\tabcolsep\relax}|}{\begin{enumerate}[start=3]\item Metodologias de pesquisa aplicada e extensão universitária.\end{enumerate}}\\
\multicolumn{4}{|p{\dimexpr 16cm + 6\tabcolsep\relax}|}{\begin{enumerate}[start=4]\item Planejamento e gestão de projetos colaborativos.\end{enumerate}}\\
\multicolumn{4}{|p{\dimexpr 16cm + 6\tabcolsep\relax}|}{\begin{enumerate}[start=5]\item Aplicação de técnicas de inteligência artificial e computação para desenvolvimento de soluções.\end{enumerate}}\\
\multicolumn{4}{|p{\dimexpr 16cm + 6\tabcolsep\relax}|}{\begin{enumerate}[start=6]\item Desenvolvimento, implementação e prototipagem de projetos.\end{enumerate}}\\
\multicolumn{4}{|p{\dimexpr 16cm + 6\tabcolsep\relax}|}{\begin{enumerate}[start=7]\item Avaliação de impacto social e técnico dos projetos.\end{enumerate}}\\
\multicolumn{4}{|p{\dimexpr 16cm + 6\tabcolsep\relax}|}{\begin{enumerate}[start=8]\item Documentação e apresentação dos resultados.\end{enumerate}}\\
\multicolumn{4}{|p{\dimexpr 16cm + 6\tabcolsep\relax}|}{\begin{enumerate}[start=9]\item Trabalho em equipe e comunicação interdisciplinar.\end{enumerate}}\\
\multicolumn{4}{|p{\dimexpr 16cm + 6\tabcolsep\relax}|}{\begin{enumerate}[start=10]\item Ética, responsabilidade social e sustentabilidade em projetos tecnológicos.\end{enumerate}}\\
\multicolumn{4}{|p{\dimexpr 16cm + 6\tabcolsep\relax}|}{\begin{enumerate}[start=11]\item \textbf{Perfil da Comunidade:} A disciplina atua em parceria com comunidades locais, órgãos públicos, escolas, ONGs, associações comunitárias e pequenos negócios que enfrentam desafios reais relacionados a tecnologia, gestão ou inclusão digital. Os projetos são desenvolvidos de forma colaborativa, com participação da comunidade na identificação das necessidades, definição das soluções e avaliação dos resultados, priorizando contextos com demandas sociais claras e potencial de impacto.\end{enumerate}}\\
\multicolumn{4}{|p{\dimexpr 16cm + 6\tabcolsep\relax}|}{\begin{enumerate}[start=12]\item \textbf{Objetivos Extencionistas:} A disciplina visa (i) identificar e co-construir soluções tecnológicas para problemas reais em conjunto com comunidades e instituições locais, (ii) aplicar técnicas de computação e inteligência artificial de modo a gerar impacto social e apoiar demandas comunitárias, (iii) promover a inclusão digital e a democratização do conhecimento aproximando a universidade da sociedade, (iv) desenvolver competências sociotécnicas nos discentes, como comunicação, ética, cooperação e responsabilidade social, e (v) fortalecer o papel da UFOP como agente de transformação social por meio de projetos colaborativos e tecnologicamente orientados.\end{enumerate}}\\
\multicolumn{4}{|p{\dimexpr 16cm + 6\tabcolsep\relax}|}{}\\
\hline
\multicolumn{4}{|p{16cm}|}{Bibliografia Básica:}\\*
\multicolumn{4}{|p{\dimexpr 16cm + 6\tabcolsep\relax}|}{\begin{itemize}\item BRASIL. Ministério da Educação. Conselho Nacional de Educação. Câmara de Educação Superior. Resolução nº 7, de 18 de dezembro de 2018: estabelece as Diretrizes para a Extensão na Educação Superior Brasileira. Diário Oficial da União, Brasília, DF, 19 dez. 2018, Seção 1, p. 49-50.\end{itemize}}\\
\multicolumn{4}{|p{\dimexpr 16cm + 6\tabcolsep\relax}|}{\begin{itemize}\item MORAIS, Regis de (org.). Filosofia da ciência e da tecnologia: introdução metodológica e crítica. Campinas: Papirus, 2013.\end{itemize}}\\
\multicolumn{4}{|p{\dimexpr 16cm + 6\tabcolsep\relax}|}{\begin{itemize}\item BENDER, William N. Aprendizagem baseada em projetos: educação diferenciada para o século XXI. Tradução de Fernando de Siqueira Rodrigues. Porto Alegre: Penso, 2014.\end{itemize}}\\
\multicolumn{4}{|p{16cm}|}{}\\
\hline
\multicolumn{4}{|p{16cm}|}{Bibliografia Complementar:}\\*
\multicolumn{4}{|p{\dimexpr 16cm + 6\tabcolsep\relax}|}{\begin{itemize}\item DAGNINO, Renato (org.). Tecnologia social: ferramenta para construir outra sociedade. 2. ed. rev. e ampl. Campinas: Komedi, 2010.\end{itemize}}\\
\multicolumn{4}{|p{\dimexpr 16cm + 6\tabcolsep\relax}|}{\begin{itemize}\item BROWN, Tim. Change by Design: How Design Thinking Creates New Alternatives for Business and Society. HarperBusiness, 2009.\end{itemize}}\\
\multicolumn{4}{|p{\dimexpr 16cm + 6\tabcolsep\relax}|}{\begin{itemize}\item AXELOS. Managing successful projects with PRINCE2. 6. ed. Norwich: TSO (The Stationery Office), 2017.\end{itemize}}\\
\multicolumn{4}{|p{\dimexpr 16cm + 6\tabcolsep\relax}|}{\begin{itemize}\item RUSSELL, Stuart; NORVIG, Peter. Inteligência Artificial. 3. ed. Rio de Janeiro: Elsevier, 2013. (Capítulo sobre aplicações práticas).\end{itemize}}\\
\multicolumn{4}{|p{\dimexpr 16cm + 6\tabcolsep\relax}|}{\begin{itemize}\item MUNIZ, Antonio et al. Inteligência artificial: entenda como a IA pode impactar no mercado de trabalho e na sociedade. Brasport, 2024.\end{itemize}}\\
\hline
\end{longtable}
\end{center}

\clearpage
\begin{center}
\begin{tabular}{lccr}
 \multirow{3}{*}{\includegraphics[height=2.7cm]{brasao.png}} &
 \multicolumn{2}{c}{\bfseries UNIVERSIDADE FEDERAL DE OURO PRETO} &
 \ \ \ \ \multirow{3}{*}{\includegraphics[height=2.7cm]{ufop.png}} \\
 & \multicolumn{2}{c}{\bfseries PRÓ-REITORIA DE GRADUAÇÃO} & \\
 & \multicolumn{2}{c}{} & \\
 & \multicolumn{2}{c}{\Large\bfseries PROGRAMA DE DISCIPLINA} & \\
\end{tabular}
\end{center}

\begin{center}
\begin{longtable}{|p{4cm}|p{4cm}|p{4cm}|p{4cm}|}
\hline
\multicolumn{3}{|p{12cm}|}{Nome do Componente Curricular em Português:} &
\multicolumn{1}{p{4cm}|}{Código:} \\ 
\multicolumn{3}{|p{12cm}|}{\textbf{Projeto 2: Resolução de Problemas do Mundo Real}} &
\textbf{BIA202}\\ 
\multicolumn{3}{|p{12cm}|}{Nome do Componente Curricular em Inglês:} & \\ 
\multicolumn{3}{|p{12cm}|}{\textbf{Project 2: Real-World Problem Solving}} & \\ 
\hline
\multicolumn{3}{|p{12cm}|}{Nome e Sigla do Departamento} & Unidade Acadêmica: \\ 
\multicolumn{3}{|p{12cm}|}{Departamento de Computação (DECOM)} & {ICEB} \\ 
\hline
\multicolumn{4}{|p{16cm}|}{Modalidade de Oferta:
[X] presencial \hspace{0.6cm}
[ ] semipresencial \hspace{0.6cm}
[ ] à distância}\\
\hline
\multicolumn{2}{|p{8cm}|}{Carga horária semestral} &
\multicolumn{2}{p{8cm}|}{Carga horária semanal}\\
\hline
\multicolumn{1}{|p{4cm}|}{Total} &
\multicolumn{1}{p{4cm}|}{Extensionista} &
\multicolumn{1}{p{4cm}|}{Teórica} &
\multicolumn{1}{p{4cm}|}{Prática} \\ 
\multicolumn{1}{|p{4cm}|}{120\,horas} &
\multicolumn{1}{p{4cm}|}{8\;horas} &
\multicolumn{1}{p{4cm}|}{0\;horas/aula} &
\multicolumn{1}{p{4cm}|}{0\;horas/aula} \\ 
\hline
\multicolumn{4}{|p{16cm}|}{Ementa:}\\
\multicolumn{4}{|p{16cm}|}{}\\
\multicolumn{4}{|p{\dimexpr 16cm + 6\tabcolsep\relax}|}{Continuação do desenvolvimento de projetos extensionistas focados na solução de problemas reais, com aprofundamento técnico e interdisciplinar. Ênfase em prototipagem avançada, validação e avaliação de impacto social e tecnológico.}\\
\multicolumn{4}{|p{16cm}|}{}\\
\hline
\multicolumn{4}{|p{16cm}|}{Conteúdo programático:}\\*
\multicolumn{4}{|p{\dimexpr 16cm + 6\tabcolsep\relax}|}{\begin{enumerate}[start=1]\item Revisão e aprimoramento dos projetos iniciados em Projeto 1.\end{enumerate}}\\
\multicolumn{4}{|p{\dimexpr 16cm + 6\tabcolsep\relax}|}{\begin{enumerate}[start=2]\item Técnicas avançadas para desenvolvimento e otimização de soluções.\end{enumerate}}\\
\multicolumn{4}{|p{\dimexpr 16cm + 6\tabcolsep\relax}|}{\begin{enumerate}[start=3]\item Validação e testes com usuários e stakeholders.\end{enumerate}}\\
\multicolumn{4}{|p{\dimexpr 16cm + 6\tabcolsep\relax}|}{\begin{enumerate}[start=4]\item Metodologias para avaliação de impacto social e técnico.\end{enumerate}}\\
\multicolumn{4}{|p{\dimexpr 16cm + 6\tabcolsep\relax}|}{\begin{enumerate}[start=5]\item Documentação técnica e relatórios de progresso.\end{enumerate}}\\
\multicolumn{4}{|p{\dimexpr 16cm + 6\tabcolsep\relax}|}{\begin{enumerate}[start=6]\item Comunicação científica e apresentação de resultados.\end{enumerate}}\\
\multicolumn{4}{|p{\dimexpr 16cm + 6\tabcolsep\relax}|}{\begin{enumerate}[start=7]\item Gestão de projetos e planejamento de etapas finais.\end{enumerate}}\\
\multicolumn{4}{|p{\dimexpr 16cm + 6\tabcolsep\relax}|}{\begin{enumerate}[start=8]\item Ética e responsabilidade social no desenvolvimento tecnológico.\end{enumerate}}\\
\multicolumn{4}{|p{\dimexpr 16cm + 6\tabcolsep\relax}|}{\begin{enumerate}[start=9]\item Trabalho colaborativo interdisciplinar e com a comunidade.\end{enumerate}}\\
\multicolumn{4}{|p{\dimexpr 16cm + 6\tabcolsep\relax}|}{\begin{enumerate}[start=10]\item Preparação para continuidade e escalabilidade dos projetos.\end{enumerate}}\\
\multicolumn{4}{|p{\dimexpr 16cm + 6\tabcolsep\relax}|}{\begin{enumerate}[start=11]\item \textbf{Perfil da Comunidade:} A disciplina atua em parceria com comunidades locais, órgãos públicos, escolas, ONGs, associações comunitárias e pequenos negócios que enfrentam desafios reais relacionados a tecnologia, gestão ou inclusão digital. Os projetos são desenvolvidos de forma colaborativa, com participação da comunidade na identificação das necessidades, definição das soluções e avaliação dos resultados, priorizando contextos com demandas sociais claras e potencial de impacto.\end{enumerate}}\\
\multicolumn{4}{|p{\dimexpr 16cm + 6\tabcolsep\relax}|}{\begin{enumerate}[start=12]\item \textbf{Objetivos Extencionistas:} A disciplina visa (i) identificar e co-construir soluções tecnológicas para problemas reais em conjunto com comunidades e instituições locais, (ii) aplicar técnicas de computação e inteligência artificial de modo a gerar impacto social e apoiar demandas comunitárias, (iii) promover a inclusão digital e a democratização do conhecimento aproximando a universidade da sociedade, (iv) desenvolver competências sociotécnicas nos discentes, como comunicação, ética, cooperação e responsabilidade social, e (v) fortalecer o papel da UFOP como agente de transformação social por meio de projetos colaborativos e tecnologicamente orientados.\end{enumerate}}\\
\multicolumn{4}{|p{\dimexpr 16cm + 6\tabcolsep\relax}|}{}\\
\hline
\multicolumn{4}{|p{16cm}|}{Bibliografia Básica:}\\*
\multicolumn{4}{|p{\dimexpr 16cm + 6\tabcolsep\relax}|}{\begin{itemize}\item BRASIL. Ministério da Educação. Conselho Nacional de Educação. Câmara de Educação Superior. Resolução nº 7, de 18 de dezembro de 2018: estabelece as Diretrizes para a Extensão na Educação Superior Brasileira. Diário Oficial da União, Brasília, DF, 19 dez. 2018, Seção 1, p. 49-50.\end{itemize}}\\
\multicolumn{4}{|p{\dimexpr 16cm + 6\tabcolsep\relax}|}{\begin{itemize}\item MORAIS, Regis de (org.). Filosofia da ciência e da tecnologia: introdução metodológica e crítica. Campinas: Papirus, 2013.\end{itemize}}\\
\multicolumn{4}{|p{\dimexpr 16cm + 6\tabcolsep\relax}|}{\begin{itemize}\item BENDER, William N. Aprendizagem baseada em projetos: educação diferenciada para o século XXI. Tradução de Fernando de Siqueira Rodrigues. Porto Alegre: Penso, 2014.\end{itemize}}\\
\multicolumn{4}{|p{16cm}|}{}\\
\hline
\multicolumn{4}{|p{16cm}|}{Bibliografia Complementar:}\\*
\multicolumn{4}{|p{\dimexpr 16cm + 6\tabcolsep\relax}|}{\begin{itemize}\item DAGNINO, Renato (org.). Tecnologia social: ferramenta para construir outra sociedade. 2. ed. rev. e ampl. Campinas: Komedi, 2010.\end{itemize}}\\
\multicolumn{4}{|p{\dimexpr 16cm + 6\tabcolsep\relax}|}{\begin{itemize}\item BROWN, Tim. Change by Design: How Design Thinking Creates New Alternatives for Business and Society. HarperBusiness, 2009.\end{itemize}}\\
\multicolumn{4}{|p{\dimexpr 16cm + 6\tabcolsep\relax}|}{\begin{itemize}\item AXELOS. Managing successful projects with PRINCE2. 6. ed. Norwich: TSO (The Stationery Office), 2017.\end{itemize}}\\
\multicolumn{4}{|p{\dimexpr 16cm + 6\tabcolsep\relax}|}{\begin{itemize}\item RUSSELL, Stuart; NORVIG, Peter. Inteligência Artificial. 3. ed. Rio de Janeiro: Elsevier, 2013. (Capítulo sobre aplicações práticas).\end{itemize}}\\
\multicolumn{4}{|p{\dimexpr 16cm + 6\tabcolsep\relax}|}{\begin{itemize}\item MUNIZ, Antonio et al. Inteligência artificial: entenda como a IA pode impactar no mercado de trabalho e na sociedade. Brasport, 2024.\end{itemize}}\\
\hline
\end{longtable}
\end{center}

\clearpage
\begin{center}
\begin{tabular}{lccr}
 \multirow{3}{*}{\includegraphics[height=2.7cm]{brasao.png}} &
 \multicolumn{2}{c}{\bfseries UNIVERSIDADE FEDERAL DE OURO PRETO} &
 \ \ \ \ \multirow{3}{*}{\includegraphics[height=2.7cm]{ufop.png}} \\
 & \multicolumn{2}{c}{\bfseries PRÓ-REITORIA DE GRADUAÇÃO} & \\
 & \multicolumn{2}{c}{} & \\
 & \multicolumn{2}{c}{\Large\bfseries PROGRAMA DE DISCIPLINA} & \\
\end{tabular}
\end{center}

\begin{center}
\begin{longtable}{|p{4cm}|p{4cm}|p{4cm}|p{4cm}|}
\hline
\multicolumn{3}{|p{12cm}|}{Nome do Componente Curricular em Português:} &
\multicolumn{1}{p{4cm}|}{Código:} \\ 
\multicolumn{3}{|p{12cm}|}{\textbf{Projeto 3: Resolução de Problemas do Mundo Real}} &
\textbf{BIA203}\\ 
\multicolumn{3}{|p{12cm}|}{Nome do Componente Curricular em Inglês:} & \\ 
\multicolumn{3}{|p{12cm}|}{\textbf{Project 3: Real-World Problem Solving}} & \\ 
\hline
\multicolumn{3}{|p{12cm}|}{Nome e Sigla do Departamento} & Unidade Acadêmica: \\ 
\multicolumn{3}{|p{12cm}|}{Departamento de Computação (DECOM)} & {ICEB} \\ 
\hline
\multicolumn{4}{|p{16cm}|}{Modalidade de Oferta:
[X] presencial \hspace{0.6cm}
[ ] semipresencial \hspace{0.6cm}
[ ] à distância}\\
\hline
\multicolumn{2}{|p{8cm}|}{Carga horária semestral} &
\multicolumn{2}{p{8cm}|}{Carga horária semanal}\\
\hline
\multicolumn{1}{|p{4cm}|}{Total} &
\multicolumn{1}{p{4cm}|}{Extensionista} &
\multicolumn{1}{p{4cm}|}{Teórica} &
\multicolumn{1}{p{4cm}|}{Prática} \\ 
\multicolumn{1}{|p{4cm}|}{120\,horas} &
\multicolumn{1}{p{4cm}|}{8\;horas} &
\multicolumn{1}{p{4cm}|}{0\;horas/aula} &
\multicolumn{1}{p{4cm}|}{0\;horas/aula} \\ 
\hline
\multicolumn{4}{|p{16cm}|}{Ementa:}\\
\multicolumn{4}{|p{16cm}|}{}\\
\multicolumn{4}{|p{\dimexpr 16cm + 6\tabcolsep\relax}|}{Conclusão dos projetos extensionistas iniciados nas disciplinas anteriores, com foco na entrega final, avaliação de resultados, documentação completa e disseminação das soluções desenvolvidas para problemas reais.}\\
\multicolumn{4}{|p{16cm}|}{}\\
\hline
\multicolumn{4}{|p{16cm}|}{Conteúdo programático:}\\*
\multicolumn{4}{|p{\dimexpr 16cm + 6\tabcolsep\relax}|}{\begin{enumerate}[start=1]\item Finalização e refinamento das soluções desenvolvidas.\end{enumerate}}\\
\multicolumn{4}{|p{\dimexpr 16cm + 6\tabcolsep\relax}|}{\begin{enumerate}[start=2]\item Validação definitiva com usuários e stakeholders.\end{enumerate}}\\
\multicolumn{4}{|p{\dimexpr 16cm + 6\tabcolsep\relax}|}{\begin{enumerate}[start=3]\item Avaliação do impacto social, econômico e tecnológico.\end{enumerate}}\\
\multicolumn{4}{|p{\dimexpr 16cm + 6\tabcolsep\relax}|}{\begin{enumerate}[start=4]\item Documentação técnica, científica e de extensão.\end{enumerate}}\\
\multicolumn{4}{|p{\dimexpr 16cm + 6\tabcolsep\relax}|}{\begin{enumerate}[start=5]\item Preparação e realização de apresentações e defesas dos projetos.\end{enumerate}}\\
\multicolumn{4}{|p{\dimexpr 16cm + 6\tabcolsep\relax}|}{\begin{enumerate}[start=6]\item Divulgação e publicação dos resultados em eventos e mídias acadêmicas e comunitárias.\end{enumerate}}\\
\multicolumn{4}{|p{\dimexpr 16cm + 6\tabcolsep\relax}|}{\begin{enumerate}[start=7]\item Planejamento para continuidade, manutenção e escalabilidade das soluções.\end{enumerate}}\\
\multicolumn{4}{|p{\dimexpr 16cm + 6\tabcolsep\relax}|}{\begin{enumerate}[start=8]\item Ética e responsabilidade na aplicação dos projetos.\end{enumerate}}\\
\multicolumn{4}{|p{\dimexpr 16cm + 6\tabcolsep\relax}|}{\begin{enumerate}[start=9]\item Trabalho colaborativo interdisciplinar e comunitário.\end{enumerate}}\\
\multicolumn{4}{|p{\dimexpr 16cm + 6\tabcolsep\relax}|}{\begin{enumerate}[start=10]\item Reflexão crítica sobre o processo de extensão e inovação social.\end{enumerate}}\\
\multicolumn{4}{|p{\dimexpr 16cm + 6\tabcolsep\relax}|}{\begin{enumerate}[start=11]\item \textbf{Perfil da Comunidade:} A disciplina atua em parceria com comunidades locais, órgãos públicos, escolas, ONGs, associações comunitárias e pequenos negócios que enfrentam desafios reais relacionados a tecnologia, gestão ou inclusão digital. Os projetos são desenvolvidos de forma colaborativa, com participação da comunidade na identificação das necessidades, definição das soluções e avaliação dos resultados, priorizando contextos com demandas sociais claras e potencial de impacto.\end{enumerate}}\\
\multicolumn{4}{|p{\dimexpr 16cm + 6\tabcolsep\relax}|}{\begin{enumerate}[start=12]\item \textbf{Objetivos Extencionistas:} A disciplina visa (i) identificar e co-construir soluções tecnológicas para problemas reais em conjunto com comunidades e instituições locais, (ii) aplicar técnicas de computação e inteligência artificial de modo a gerar impacto social e apoiar demandas comunitárias, (iii) promover a inclusão digital e a democratização do conhecimento aproximando a universidade da sociedade, (iv) desenvolver competências sociotécnicas nos discentes, como comunicação, ética, cooperação e responsabilidade social, e (v) fortalecer o papel da UFOP como agente de transformação social por meio de projetos colaborativos e tecnologicamente orientados.\end{enumerate}}\\
\multicolumn{4}{|p{\dimexpr 16cm + 6\tabcolsep\relax}|}{}\\
\hline
\multicolumn{4}{|p{16cm}|}{Bibliografia Básica:}\\*
\multicolumn{4}{|p{\dimexpr 16cm + 6\tabcolsep\relax}|}{\begin{itemize}\item BRASIL. Ministério da Educação. Conselho Nacional de Educação. Câmara de Educação Superior. Resolução nº 7, de 18 de dezembro de 2018: estabelece as Diretrizes para a Extensão na Educação Superior Brasileira. Diário Oficial da União, Brasília, DF, 19 dez. 2018, Seção 1, p. 49-50.\end{itemize}}\\
\multicolumn{4}{|p{\dimexpr 16cm + 6\tabcolsep\relax}|}{\begin{itemize}\item MORAIS, Regis de (org.). Filosofia da ciência e da tecnologia: introdução metodológica e crítica. Campinas: Papirus, 2013.\end{itemize}}\\
\multicolumn{4}{|p{\dimexpr 16cm + 6\tabcolsep\relax}|}{\begin{itemize}\item BENDER, William N. Aprendizagem baseada em projetos: educação diferenciada para o século XXI. Tradução de Fernando de Siqueira Rodrigues. Porto Alegre: Penso, 2014.\end{itemize}}\\
\multicolumn{4}{|p{16cm}|}{}\\
\hline
\multicolumn{4}{|p{16cm}|}{Bibliografia Complementar:}\\*
\multicolumn{4}{|p{\dimexpr 16cm + 6\tabcolsep\relax}|}{\begin{itemize}\item DAGNINO, Renato (org.). Tecnologia social: ferramenta para construir outra sociedade. 2. ed. rev. e ampl. Campinas: Komedi, 2010.\end{itemize}}\\
\multicolumn{4}{|p{\dimexpr 16cm + 6\tabcolsep\relax}|}{\begin{itemize}\item BROWN, Tim. Change by Design: How Design Thinking Creates New Alternatives for Business and Society. HarperBusiness, 2009.\end{itemize}}\\
\multicolumn{4}{|p{\dimexpr 16cm + 6\tabcolsep\relax}|}{\begin{itemize}\item AXELOS. Managing successful projects with PRINCE2. 6. ed. Norwich: TSO (The Stationery Office), 2017.\end{itemize}}\\
\multicolumn{4}{|p{\dimexpr 16cm + 6\tabcolsep\relax}|}{\begin{itemize}\item RUSSELL, Stuart; NORVIG, Peter. Inteligência Artificial. 3. ed. Rio de Janeiro: Elsevier, 2013. (Capítulo sobre aplicações práticas).\end{itemize}}\\
\multicolumn{4}{|p{\dimexpr 16cm + 6\tabcolsep\relax}|}{\begin{itemize}\item MUNIZ, Antonio et al. Inteligência artificial: entenda como a IA pode impactar no mercado de trabalho e na sociedade. Brasport, 2024.\end{itemize}}\\
\hline
\end{longtable}
\end{center}

\clearpage
\begin{center}
\begin{tabular}{lccr}
 \multirow{3}{*}{\includegraphics[height=2.7cm]{brasao.png}} &
 \multicolumn{2}{c}{\bfseries UNIVERSIDADE FEDERAL DE OURO PRETO} &
 \ \ \ \ \multirow{3}{*}{\includegraphics[height=2.7cm]{ufop.png}} \\
 & \multicolumn{2}{c}{\bfseries PRÓ-REITORIA DE GRADUAÇÃO} & \\
 & \multicolumn{2}{c}{} & \\
 & \multicolumn{2}{c}{\Large\bfseries PROGRAMA DE DISCIPLINA} & \\
\end{tabular}
\end{center}

\begin{center}
\begin{longtable}{|p{4cm}|p{4cm}|p{4cm}|p{4cm}|}
\hline
\multicolumn{3}{|p{12cm}|}{Nome do Componente Curricular em Português:} &
\multicolumn{1}{p{4cm}|}{Código:} \\ 
\multicolumn{3}{|p{12cm}|}{\textbf{Projeto 4: Resolução de Problemas do Mundo Real}} &
\textbf{BIA204}\\ 
\multicolumn{3}{|p{12cm}|}{Nome do Componente Curricular em Inglês:} & \\ 
\multicolumn{3}{|p{12cm}|}{\textbf{Project 4: Real-World Problem Solving}} & \\ 
\hline
\multicolumn{3}{|p{12cm}|}{Nome e Sigla do Departamento} & Unidade Acadêmica: \\ 
\multicolumn{3}{|p{12cm}|}{Departamento de Computação (DECOM)} & {ICEB} \\ 
\hline
\multicolumn{4}{|p{16cm}|}{Modalidade de Oferta:
[X] presencial \hspace{0.6cm}
[ ] semipresencial \hspace{0.6cm}
[ ] à distância}\\
\hline
\multicolumn{2}{|p{8cm}|}{Carga horária semestral} &
\multicolumn{2}{p{8cm}|}{Carga horária semanal}\\
\hline
\multicolumn{1}{|p{4cm}|}{Total} &
\multicolumn{1}{p{4cm}|}{Extensionista} &
\multicolumn{1}{p{4cm}|}{Teórica} &
\multicolumn{1}{p{4cm}|}{Prática} \\ 
\multicolumn{1}{|p{4cm}|}{120\,horas} &
\multicolumn{1}{p{4cm}|}{8\;horas} &
\multicolumn{1}{p{4cm}|}{0\;horas/aula} &
\multicolumn{1}{p{4cm}|}{0\;horas/aula} \\ 
\hline
\multicolumn{4}{|p{16cm}|}{Ementa:}\\
\multicolumn{4}{|p{16cm}|}{}\\
\multicolumn{4}{|p{\dimexpr 16cm + 6\tabcolsep\relax}|}{Desenvolvimento avançado de projetos extensionistas com foco na resolução de problemas reais, incorporando aprimoramentos técnicos, análise crítica dos resultados e estratégias para sustentabilidade e escalabilidade das soluções.}\\
\multicolumn{4}{|p{16cm}|}{}\\
\hline
\multicolumn{4}{|p{16cm}|}{Conteúdo programático:}\\*
\multicolumn{4}{|p{\dimexpr 16cm + 6\tabcolsep\relax}|}{\begin{enumerate}[start=1]\item Avaliação e análise crítica dos projetos anteriores.\end{enumerate}}\\
\multicolumn{4}{|p{\dimexpr 16cm + 6\tabcolsep\relax}|}{\begin{enumerate}[start=2]\item Incorporação de melhorias técnicas e funcionais nas soluções.\end{enumerate}}\\
\multicolumn{4}{|p{\dimexpr 16cm + 6\tabcolsep\relax}|}{\begin{enumerate}[start=3]\item Estudo de viabilidade para escalabilidade e sustentabilidade.\end{enumerate}}\\
\multicolumn{4}{|p{\dimexpr 16cm + 6\tabcolsep\relax}|}{\begin{enumerate}[start=4]\item Integração de novas tecnologias e abordagens inovadoras.\end{enumerate}}\\
\multicolumn{4}{|p{\dimexpr 16cm + 6\tabcolsep\relax}|}{\begin{enumerate}[start=5]\item Documentação atualizada e elaboração de relatórios finais.\end{enumerate}}\\
\multicolumn{4}{|p{\dimexpr 16cm + 6\tabcolsep\relax}|}{\begin{enumerate}[start=6]\item Preparação para divulgação acadêmica e comunitária.\end{enumerate}}\\
\multicolumn{4}{|p{\dimexpr 16cm + 6\tabcolsep\relax}|}{\begin{enumerate}[start=7]\item Planejamento estratégico para continuidade dos projetos.\end{enumerate}}\\
\multicolumn{4}{|p{\dimexpr 16cm + 6\tabcolsep\relax}|}{\begin{enumerate}[start=8]\item Considerações éticas e responsabilidade social.\end{enumerate}}\\
\multicolumn{4}{|p{\dimexpr 16cm + 6\tabcolsep\relax}|}{\begin{enumerate}[start=9]\item Trabalho colaborativo multidisciplinar e com a comunidade.\end{enumerate}}\\
\multicolumn{4}{|p{\dimexpr 16cm + 6\tabcolsep\relax}|}{\begin{enumerate}[start=10]\item Reflexões sobre o impacto social e tecnológico dos projetos.\end{enumerate}}\\
\multicolumn{4}{|p{\dimexpr 16cm + 6\tabcolsep\relax}|}{\begin{enumerate}[start=11]\item \textbf{Perfil da Comunidade:} A disciplina atua em parceria com comunidades locais, órgãos públicos, escolas, ONGs, associações comunitárias e pequenos negócios que enfrentam desafios reais relacionados a tecnologia, gestão ou inclusão digital. Os projetos são desenvolvidos de forma colaborativa, com participação da comunidade na identificação das necessidades, definição das soluções e avaliação dos resultados, priorizando contextos com demandas sociais claras e potencial de impacto.\end{enumerate}}\\
\multicolumn{4}{|p{\dimexpr 16cm + 6\tabcolsep\relax}|}{\begin{enumerate}[start=12]\item \textbf{Objetivos Extencionistas:} A disciplina visa (i) identificar e co-construir soluções tecnológicas para problemas reais em conjunto com comunidades e instituições locais, (ii) aplicar técnicas de computação e inteligência artificial de modo a gerar impacto social e apoiar demandas comunitárias, (iii) promover a inclusão digital e a democratização do conhecimento aproximando a universidade da sociedade, (iv) desenvolver competências sociotécnicas nos discentes, como comunicação, ética, cooperação e responsabilidade social, e (v) fortalecer o papel da UFOP como agente de transformação social por meio de projetos colaborativos e tecnologicamente orientados.\end{enumerate}}\\
\multicolumn{4}{|p{\dimexpr 16cm + 6\tabcolsep\relax}|}{}\\
\hline
\multicolumn{4}{|p{16cm}|}{Bibliografia Básica:}\\*
\multicolumn{4}{|p{\dimexpr 16cm + 6\tabcolsep\relax}|}{\begin{itemize}\item BRASIL. Ministério da Educação. Conselho Nacional de Educação. Câmara de Educação Superior. Resolução nº 7, de 18 de dezembro de 2018: estabelece as Diretrizes para a Extensão na Educação Superior Brasileira. Diário Oficial da União, Brasília, DF, 19 dez. 2018, Seção 1, p. 49-50.\end{itemize}}\\
\multicolumn{4}{|p{\dimexpr 16cm + 6\tabcolsep\relax}|}{\begin{itemize}\item MORAIS, Regis de (org.). Filosofia da ciência e da tecnologia: introdução metodológica e crítica. Campinas: Papirus, 2013.\end{itemize}}\\
\multicolumn{4}{|p{\dimexpr 16cm + 6\tabcolsep\relax}|}{\begin{itemize}\item BENDER, William N. Aprendizagem baseada em projetos: educação diferenciada para o século XXI. Tradução de Fernando de Siqueira Rodrigues. Porto Alegre: Penso, 2014.\end{itemize}}\\
\multicolumn{4}{|p{16cm}|}{}\\
\hline
\multicolumn{4}{|p{16cm}|}{Bibliografia Complementar:}\\*
\multicolumn{4}{|p{\dimexpr 16cm + 6\tabcolsep\relax}|}{\begin{itemize}\item DAGNINO, Renato (org.). Tecnologia social: ferramenta para construir outra sociedade. 2. ed. rev. e ampl. Campinas: Komedi, 2010.\end{itemize}}\\
\multicolumn{4}{|p{\dimexpr 16cm + 6\tabcolsep\relax}|}{\begin{itemize}\item BROWN, Tim. Change by Design: How Design Thinking Creates New Alternatives for Business and Society. HarperBusiness, 2009.\end{itemize}}\\
\multicolumn{4}{|p{\dimexpr 16cm + 6\tabcolsep\relax}|}{\begin{itemize}\item AXELOS. Managing successful projects with PRINCE2. 6. ed. Norwich: TSO (The Stationery Office), 2017.\end{itemize}}\\
\multicolumn{4}{|p{\dimexpr 16cm + 6\tabcolsep\relax}|}{\begin{itemize}\item RUSSELL, Stuart; NORVIG, Peter. Inteligência Artificial. 3. ed. Rio de Janeiro: Elsevier, 2013. (Capítulo sobre aplicações práticas).\end{itemize}}\\
\multicolumn{4}{|p{\dimexpr 16cm + 6\tabcolsep\relax}|}{\begin{itemize}\item MUNIZ, Antonio et al. Inteligência artificial: entenda como a IA pode impactar no mercado de trabalho e na sociedade. Brasport, 2024.\end{itemize}}\\
\hline
\end{longtable}
\end{center}

\clearpage
\begin{center}
\begin{tabular}{lccr}
 \multirow{3}{*}{\includegraphics[height=2.7cm]{brasao.png}} &
 \multicolumn{2}{c}{\bfseries UNIVERSIDADE FEDERAL DE OURO PRETO} &
 \ \ \ \ \multirow{3}{*}{\includegraphics[height=2.7cm]{ufop.png}} \\
 & \multicolumn{2}{c}{\bfseries PRÓ-REITORIA DE GRADUAÇÃO} & \\
 & \multicolumn{2}{c}{} & \\
 & \multicolumn{2}{c}{\Large\bfseries PROGRAMA DE DISCIPLINA} & \\
\end{tabular}
\end{center}

\begin{center}
\begin{longtable}{|p{4cm}|p{4cm}|p{4cm}|p{4cm}|}
\hline
\multicolumn{3}{|p{12cm}|}{Nome do Componente Curricular em Português:} &
\multicolumn{1}{p{4cm}|}{Código:} \\ 
\multicolumn{3}{|p{12cm}|}{\textbf{Monografia}} &
\textbf{BIA205}\\ 
\multicolumn{3}{|p{12cm}|}{Nome do Componente Curricular em Inglês:} & \\ 
\multicolumn{3}{|p{12cm}|}{\textbf{Undergraduate Thesis}} & \\ 
\hline
\multicolumn{3}{|p{12cm}|}{Nome e Sigla do Departamento} & Unidade Acadêmica: \\ 
\multicolumn{3}{|p{12cm}|}{Departamento de Computação (DECOM)} & {ICEB} \\ 
\hline
\multicolumn{4}{|p{16cm}|}{Modalidade de Oferta:
[X] presencial \hspace{0.6cm}
[ ] semipresencial \hspace{0.6cm}
[ ] à distância}\\
\hline
\multicolumn{2}{|p{8cm}|}{Carga horária semestral} &
\multicolumn{2}{p{8cm}|}{Carga horária semanal}\\
\hline
\multicolumn{1}{|p{4cm}|}{Total} &
\multicolumn{1}{p{4cm}|}{Extensionista} &
\multicolumn{1}{p{4cm}|}{Teórica} &
\multicolumn{1}{p{4cm}|}{Prática} \\ 
\multicolumn{1}{|p{4cm}|}{120\,horas} &
\multicolumn{1}{p{4cm}|}{0\;horas} &
\multicolumn{1}{p{4cm}|}{8\;horas/aula} &
\multicolumn{1}{p{4cm}|}{0\;horas/aula} \\ 
\hline
\multicolumn{4}{|p{16cm}|}{Ementa:}\\
\multicolumn{4}{|p{16cm}|}{}\\
\multicolumn{4}{|p{\dimexpr 16cm + 6\tabcolsep\relax}|}{Desenvolvimento de projeto memorial dos projetos de resolução de problemas do mundo real. Entrega de relatórios parciais ao longo do semestre. Entrega da monografia ao final do semestre. Apresentação oral dos resultados da monografia.}\\
\multicolumn{4}{|p{16cm}|}{}\\
\hline
\multicolumn{4}{|p{16cm}|}{Conteúdo programático:}\\*
\multicolumn{4}{|p{\dimexpr 16cm + 6\tabcolsep\relax}|}{\begin{enumerate}[start=1]\item Construção do projeto memorial.\end{enumerate}}\\
\multicolumn{4}{|p{\dimexpr 16cm + 6\tabcolsep\relax}|}{\begin{enumerate}[start=2]\item Redação de relatórios das atividades.\end{enumerate}}\\
\multicolumn{4}{|p{\dimexpr 16cm + 6\tabcolsep\relax}|}{\begin{enumerate}[start=3]\item Elaboração da monografia.\end{enumerate}}\\
\multicolumn{4}{|p{\dimexpr 16cm + 6\tabcolsep\relax}|}{\begin{enumerate}[start=4]\item Apresentação oral dos resultados finais obtidos.\end{enumerate}}\\
\multicolumn{4}{|p{\dimexpr 16cm + 6\tabcolsep\relax}|}{}\\
\hline
\multicolumn{4}{|p{16cm}|}{Bibliografia Básica:}\\*
\multicolumn{4}{|p{\dimexpr 16cm + 6\tabcolsep\relax}|}{\begin{itemize}\item WAZLAWICK, Raul Sidnei. Metodologia de pesquisa para Ciência da Computação. 2. ed. Rio de Janeiro: Elsevier, 2014.\end{itemize}}\\
\multicolumn{4}{|p{\dimexpr 16cm + 6\tabcolsep\relax}|}{\begin{itemize}\item MARCONI, Marina de Andrade; LAKATOS, Eva Maria. Fundamentos de metodologia científica. 8. ed. São Paulo: Atlas, 2017.\end{itemize}}\\
\multicolumn{4}{|p{\dimexpr 16cm + 6\tabcolsep\relax}|}{\begin{itemize}\item Associação Brasileira de Normas Técnicas. NBR 14724: Informação e Documentação - Trabalhos acadêmicos - Apresentação. Rio de Janeiro, 2011.\end{itemize}}\\
\multicolumn{4}{|p{16cm}|}{}\\
\hline
\multicolumn{4}{|p{16cm}|}{Bibliografia Complementar:}\\*
\multicolumn{4}{|p{\dimexpr 16cm + 6\tabcolsep\relax}|}{\begin{itemize}\item GIL, Antônio Carlos. Como elaborar projetos de pesquisa. 6. ed. São Paulo: Atlas, 2017.\end{itemize}}\\
\multicolumn{4}{|p{\dimexpr 16cm + 6\tabcolsep\relax}|}{\begin{itemize}\item FRANÇA, Júnia Lessa; VASCONCELLOS, Ana Cristina de. Manual para normalização de publicações técnico-científicas. 8. ed. Belo Horizonte: Editora UFMG, 2009.\end{itemize}}\\
\multicolumn{4}{|p{\dimexpr 16cm + 6\tabcolsep\relax}|}{\begin{itemize}\item MEDEIROS, João Bosco. Redação Científica: práticas de fichamentos, resumos, resenhas. 13. ed. São Paulo: Atlas, 2019.\end{itemize}}\\
\multicolumn{4}{|p{\dimexpr 16cm + 6\tabcolsep\relax}|}{\begin{itemize}\item MARCONI, Marina de Andrade; LAKATOS, Eva Maria. Metodologia científica. 7. ed. São Paulo: Atlas, 2017.\end{itemize}}\\
\multicolumn{4}{|p{\dimexpr 16cm + 6\tabcolsep\relax}|}{\begin{itemize}\item MATIAS-PEREIRA, José. Metodologia da pesquisa científica. 4. ed. São Paulo: Atlas, 2016.\end{itemize}}\\
\multicolumn{4}{|p{\dimexpr 16cm + 6\tabcolsep\relax}|}{\begin{itemize}\item FERNANDEZ, Brena Paula Magno. Métodos e técnicas de pesquisa. 7. ed. Editora Saraiva, 2012.\end{itemize}}\\
\multicolumn{4}{|p{\dimexpr 16cm + 6\tabcolsep\relax}|}{\begin{itemize}\item CERVO, Amado L.; BERVIAN, Pedro A.; SILVA, Roberto da. Metodologia Científica. 6. ed. São Paulo: Pearson, 2006.\end{itemize}}\\
\hline
\end{longtable}
\end{center}

\clearpage
\begin{center}
\begin{tabular}{lccr}
 \multirow{3}{*}{\includegraphics[height=2.7cm]{brasao.png}} &
 \multicolumn{2}{c}{\bfseries UNIVERSIDADE FEDERAL DE OURO PRETO} &
 \ \ \ \ \multirow{3}{*}{\includegraphics[height=2.7cm]{ufop.png}} \\
 & \multicolumn{2}{c}{\bfseries PRÓ-REITORIA DE GRADUAÇÃO} & \\
 & \multicolumn{2}{c}{} & \\
 & \multicolumn{2}{c}{\Large\bfseries PROGRAMA DE DISCIPLINA} & \\
\end{tabular}
\end{center}

\begin{center}
\begin{longtable}{|p{4cm}|p{4cm}|p{4cm}|p{4cm}|}
\hline
\multicolumn{3}{|p{12cm}|}{Nome do Componente Curricular em Português:} &
\multicolumn{1}{p{4cm}|}{Código:} \\ 
\multicolumn{3}{|p{12cm}|}{\textbf{Cálculo Diferencial e Integral A}} &
\textbf{MTM700}\\ 
\multicolumn{3}{|p{12cm}|}{Nome do Componente Curricular em Inglês:} & \\ 
\multicolumn{3}{|p{12cm}|}{\textbf{Differential and Integral Calculus A}} & \\ 
\hline
\multicolumn{3}{|p{12cm}|}{Nome e Sigla do Departamento} & Unidade Acadêmica: \\ 
\multicolumn{3}{|p{12cm}|}{Departamento de Matemática (DEMAT)} & {ICEB} \\ 
\hline
\multicolumn{4}{|p{16cm}|}{Modalidade de Oferta:
[X] presencial \hspace{0.6cm}
[ ] semipresencial \hspace{0.6cm}
[ ] à distância}\\
\hline
\multicolumn{2}{|p{8cm}|}{Carga horária semestral} &
\multicolumn{2}{p{8cm}|}{Carga horária semanal}\\
\hline
\multicolumn{1}{|p{4cm}|}{Total} &
\multicolumn{1}{p{4cm}|}{Extensionista} &
\multicolumn{1}{p{4cm}|}{Teórica} &
\multicolumn{1}{p{4cm}|}{Prática} \\ 
\multicolumn{1}{|p{4cm}|}{60\,horas} &
\multicolumn{1}{p{4cm}|}{0\;horas} &
\multicolumn{1}{p{4cm}|}{4\;horas/aula} &
\multicolumn{1}{p{4cm}|}{0\;horas/aula} \\ 
\hline
\multicolumn{4}{|p{16cm}|}{Ementa:}\\
\multicolumn{4}{|p{16cm}|}{}\\
\multicolumn{4}{|p{\dimexpr 16cm + 6\tabcolsep\relax}|}{Limite. Continuidade. Derivadas e Aplicações. Integral.}\\
\multicolumn{4}{|p{16cm}|}{}\\
\hline
\multicolumn{4}{|p{16cm}|}{Conteúdo programático:}\\*
\multicolumn{4}{|p{\dimexpr 16cm + 6\tabcolsep\relax}|}{\begin{enumerate}[start=1]\item Limite: Definição; Limites Laterais; Limites no Infinito; Limites Infinitos; Propriedades; Limites Fundamentais\end{enumerate}}\\
\multicolumn{4}{|p{\dimexpr 16cm + 6\tabcolsep\relax}|}{\begin{enumerate}[start=2]\item Continuidade: Definição; Propriedades; Teorema do Valor Intermediário\end{enumerate}}\\
\multicolumn{4}{|p{\dimexpr 16cm + 6\tabcolsep\relax}|}{\begin{enumerate}[start=3]\item Derivadas: Retas Tangentes a curvas; Definição; Regras de Derivação; Derivadas de Funções Trigonométricas e Exponenciais; Derivadas de Funções Inversas; Derivadas das Funções Trigonométricas Inversas e Logarítmica; Derivada de Funções Compostas (Regra da Cadeia)\end{enumerate}}\\
\multicolumn{4}{|p{\dimexpr 16cm + 6\tabcolsep\relax}|}{\begin{enumerate}[start=4]\item Aplicações da Derivada: Máximos e Mínimos; Teorema do Valor Médio; Crescimento e Concavidade de Funções; Gráfico de Funções; Problemas de Máximos e Mínimos; Regra de L'Hospital\end{enumerate}}\\
\multicolumn{4}{|p{\dimexpr 16cm + 6\tabcolsep\relax}|}{\begin{enumerate}[start=5]\item Integrais: Integral Indefinida e Suas Propriedades; Integral Definida e Suas Propriedades; Técnicas de Integração (Substituição, Por Partes, Frações Parciais, Substituição Trigonométrica)\end{enumerate}}\\
\multicolumn{4}{|p{\dimexpr 16cm + 6\tabcolsep\relax}|}{}\\
\hline
\multicolumn{4}{|p{16cm}|}{Bibliografia Básica:}\\*
\multicolumn{4}{|p{\dimexpr 16cm + 6\tabcolsep\relax}|}{\begin{itemize}\item Swokowski, Earl W.; tradução Alfredo Alves de Faria; revisão técnica Victor Hugo Teixeira Rodrigues e Antonio Gabriel da Silva St. Aubyn. Cálculo com geometria analítica: volume 1. 2. ed. São Paulo: McGraw-Hill do Brasil, 1995.\end{itemize}}\\
\multicolumn{4}{|p{\dimexpr 16cm + 6\tabcolsep\relax}|}{\begin{itemize}\item Thomas, George Brinton. Cálculo volume 1. 10. ed. São Paulo: Addison Wesley, 2002. v.1.\end{itemize}}\\
\multicolumn{4}{|p{16cm}|}{}\\
\hline
\multicolumn{4}{|p{16cm}|}{Bibliografia Complementar:}\\*
\multicolumn{4}{|p{\dimexpr 16cm + 6\tabcolsep\relax}|}{\begin{itemize}\item Stewart, James; tradução Helena Maria Ávila de Castro. Cálculo: volume 1. 8. ed. São Paulo: Cengage Learning, c2017.\end{itemize}}\\
\multicolumn{4}{|p{\dimexpr 16cm + 6\tabcolsep\relax}|}{\begin{itemize}\item Leithold, Louis. O cálculo com geometria analítica. 2. ed. v.1. São Paulo: HARBRA, c1982.\end{itemize}}\\
\hline
\end{longtable}
\end{center}

\clearpage
\begin{center}
\begin{tabular}{lccr}
 \multirow{3}{*}{\includegraphics[height=2.7cm]{brasao.png}} &
 \multicolumn{2}{c}{\bfseries UNIVERSIDADE FEDERAL DE OURO PRETO} &
 \ \ \ \ \multirow{3}{*}{\includegraphics[height=2.7cm]{ufop.png}} \\
 & \multicolumn{2}{c}{\bfseries PRÓ-REITORIA DE GRADUAÇÃO} & \\
 & \multicolumn{2}{c}{} & \\
 & \multicolumn{2}{c}{\Large\bfseries PROGRAMA DE DISCIPLINA} & \\
\end{tabular}
\end{center}

\begin{center}
\begin{longtable}{|p{4cm}|p{4cm}|p{4cm}|p{4cm}|}
\hline
\multicolumn{3}{|p{12cm}|}{Nome do Componente Curricular em Português:} &
\multicolumn{1}{p{4cm}|}{Código:} \\ 
\multicolumn{3}{|p{12cm}|}{\textbf{Cálculo Diferencial e Integral B}} &
\textbf{MTM702}\\ 
\multicolumn{3}{|p{12cm}|}{Nome do Componente Curricular em Inglês:} & \\ 
\multicolumn{3}{|p{12cm}|}{\textbf{Differential and Integral Calculus B}} & \\ 
\hline
\multicolumn{3}{|p{12cm}|}{Nome e Sigla do Departamento} & Unidade Acadêmica: \\ 
\multicolumn{3}{|p{12cm}|}{Departamento de Matemática (DEMAT)} & {ICEB} \\ 
\hline
\multicolumn{4}{|p{16cm}|}{Modalidade de Oferta:
[X] presencial \hspace{0.6cm}
[ ] semipresencial \hspace{0.6cm}
[ ] à distância}\\
\hline
\multicolumn{2}{|p{8cm}|}{Carga horária semestral} &
\multicolumn{2}{p{8cm}|}{Carga horária semanal}\\
\hline
\multicolumn{1}{|p{4cm}|}{Total} &
\multicolumn{1}{p{4cm}|}{Extensionista} &
\multicolumn{1}{p{4cm}|}{Teórica} &
\multicolumn{1}{p{4cm}|}{Prática} \\ 
\multicolumn{1}{|p{4cm}|}{60\,horas} &
\multicolumn{1}{p{4cm}|}{0\;horas} &
\multicolumn{1}{p{4cm}|}{4\;horas/aula} &
\multicolumn{1}{p{4cm}|}{0\;horas/aula} \\ 
\hline
\multicolumn{4}{|p{16cm}|}{Ementa:}\\
\multicolumn{4}{|p{16cm}|}{}\\
\multicolumn{4}{|p{\dimexpr 16cm + 6\tabcolsep\relax}|}{Aplicações da Integral. Integrais Impróprias. Sequências e Séries Infinitas. Funções de Várias Variáveis.}\\
\multicolumn{4}{|p{16cm}|}{}\\
\hline
\multicolumn{4}{|p{16cm}|}{Conteúdo programático:}\\*
\multicolumn{4}{|p{\dimexpr 16cm + 6\tabcolsep\relax}|}{\begin{enumerate}[start=1]\item Aplicações de Integral: Área de Regiões Planas; Comprimento de Curvas; Volume de Sólidos de Revolução; Área de Superfícies de Revolução\end{enumerate}}\\
\multicolumn{4}{|p{\dimexpr 16cm + 6\tabcolsep\relax}|}{\begin{enumerate}[start=2]\item Integrais Impróprias: Definição e Exemplos\end{enumerate}}\\
\multicolumn{4}{|p{\dimexpr 16cm + 6\tabcolsep\relax}|}{\begin{enumerate}[start=3]\item Sequências e Séries Numéricas: Limites de Sequências; Critérios de Convergência para Séries; Propriedades da Convergência de Séries; Valor Aproximado do Limite de Séries\end{enumerate}}\\
\multicolumn{4}{|p{\dimexpr 16cm + 6\tabcolsep\relax}|}{\begin{enumerate}[start=4]\item Séries de Potências: Intervalos de Convergência; Derivação e Integração de Séries de Potência; Polinômios e Séries de Taylor\end{enumerate}}\\
\multicolumn{4}{|p{\dimexpr 16cm + 6\tabcolsep\relax}|}{\begin{enumerate}[start=5]\item Funções Reais de Várias Variáveis: Domínio e Imagem; Limite e Continuidade; Derivadas Parciais e Diferencial; Regra da Cadeia; Derivada Direcional e Gradiente; Derivadas Parciais de Ordem Superior; Máximos e Mínimos; Aplicações: Otimização\end{enumerate}}\\
\multicolumn{4}{|p{\dimexpr 16cm + 6\tabcolsep\relax}|}{}\\
\hline
\multicolumn{4}{|p{16cm}|}{Bibliografia Básica:}\\*
\multicolumn{4}{|p{\dimexpr 16cm + 6\tabcolsep\relax}|}{\begin{itemize}\item Thomas, George B.; Weir, Maurice D.; Hass, Joel; tradução Carlos Scalici; revisão técnica Claudio Hirofume Asano. Cálculo: volume 2. 12. ed. São Paulo: Pearson Education do Brasil, c2013.\end{itemize}}\\
\multicolumn{4}{|p{\dimexpr 16cm + 6\tabcolsep\relax}|}{\begin{itemize}\item Thomas, George B.; Weir, Maurice D.; Hass, Joel; tradução Kleber Roberto Pedroso, Regina Célia Simille de Macedo; revisão técnica Claudio Hirofume Asano. Cálculo: volume 1. 12. ed. São Paulo: Pearson Education do Brasil, c2013.\end{itemize}}\\
\multicolumn{4}{|p{\dimexpr 16cm + 6\tabcolsep\relax}|}{\begin{itemize}\item Simmons, George F. Cálculo com geometria analítica, v.1. São Paulo: Makron Books, 1987.\end{itemize}}\\
\multicolumn{4}{|p{\dimexpr 16cm + 6\tabcolsep\relax}|}{\begin{itemize}\item Simmons, George F. Cálculo com geometria analítica. São Paulo: McGraw-Hill, 1988.\end{itemize}}\\
\multicolumn{4}{|p{\dimexpr 16cm + 6\tabcolsep\relax}|}{\begin{itemize}\item Leithold, Louis. O cálculo com geometria analítica. 3. ed. v.2. São Paulo: HARBRA, c1994.\end{itemize}}\\
\multicolumn{4}{|p{\dimexpr 16cm + 6\tabcolsep\relax}|}{\begin{itemize}\item Leithold, Louis. O cálculo com geometria analítica. 3. ed. v.1. São Paulo: HARBRA, c1994.\end{itemize}}\\
\multicolumn{4}{|p{\dimexpr 16cm + 6\tabcolsep\relax}|}{\begin{itemize}\item Stewart, James. Cálculo volume I. 4. ed. São Paulo: Pioneira, 2001.\end{itemize}}\\
\multicolumn{4}{|p{\dimexpr 16cm + 6\tabcolsep\relax}|}{\begin{itemize}\item Stewart, James. Cálculo: volume 2. São Paulo: Pioneira, 2005.\end{itemize}}\\
\multicolumn{4}{|p{16cm}|}{}\\
\hline
\multicolumn{4}{|p{16cm}|}{Bibliografia Complementar:}\\*
\multicolumn{4}{|p{\dimexpr 16cm + 6\tabcolsep\relax}|}{\begin{itemize}\item[] (nenhum item informado)\end{itemize}}\\
\hline
\end{longtable}
\end{center}

\clearpage
\begin{center}
\begin{tabular}{lccr}
 \multirow{3}{*}{\includegraphics[height=2.7cm]{brasao.png}} &
 \multicolumn{2}{c}{\bfseries UNIVERSIDADE FEDERAL DE OURO PRETO} &
 \ \ \ \ \multirow{3}{*}{\includegraphics[height=2.7cm]{ufop.png}} \\
 & \multicolumn{2}{c}{\bfseries PRÓ-REITORIA DE GRADUAÇÃO} & \\
 & \multicolumn{2}{c}{} & \\
 & \multicolumn{2}{c}{\Large\bfseries PROGRAMA DE DISCIPLINA} & \\
\end{tabular}
\end{center}

\begin{center}
\begin{longtable}{|p{4cm}|p{4cm}|p{4cm}|p{4cm}|}
\hline
\multicolumn{3}{|p{12cm}|}{Nome do Componente Curricular em Português:} &
\multicolumn{1}{p{4cm}|}{Código:} \\ 
\multicolumn{3}{|p{12cm}|}{\textbf{Geometria Analítica e Álgebra Linear}} &
\textbf{MTM730}\\ 
\multicolumn{3}{|p{12cm}|}{Nome do Componente Curricular em Inglês:} & \\ 
\multicolumn{3}{|p{12cm}|}{\textbf{Analytical Geometry and Linear Algebra}} & \\ 
\hline
\multicolumn{3}{|p{12cm}|}{Nome e Sigla do Departamento} & Unidade Acadêmica: \\ 
\multicolumn{3}{|p{12cm}|}{Departamento de Matemática (DEMAT)} & {ICEB} \\ 
\hline
\multicolumn{4}{|p{16cm}|}{Modalidade de Oferta:
[X] presencial \hspace{0.6cm}
[ ] semipresencial \hspace{0.6cm}
[ ] à distância}\\
\hline
\multicolumn{2}{|p{8cm}|}{Carga horária semestral} &
\multicolumn{2}{p{8cm}|}{Carga horária semanal}\\
\hline
\multicolumn{1}{|p{4cm}|}{Total} &
\multicolumn{1}{p{4cm}|}{Extensionista} &
\multicolumn{1}{p{4cm}|}{Teórica} &
\multicolumn{1}{p{4cm}|}{Prática} \\ 
\multicolumn{1}{|p{4cm}|}{60\,horas} &
\multicolumn{1}{p{4cm}|}{0\;horas} &
\multicolumn{1}{p{4cm}|}{4\;horas/aula} &
\multicolumn{1}{p{4cm}|}{0\;horas/aula} \\ 
\hline
\multicolumn{4}{|p{16cm}|}{Ementa:}\\
\multicolumn{4}{|p{16cm}|}{}\\
\multicolumn{4}{|p{\dimexpr 16cm + 6\tabcolsep\relax}|}{Álgebra Vetorial. Retas e Planos. Matrizes, Sistemas Lineares e Determinantes. Espaços Vetoriais. Autovalores e Autovetores. Diagonalização.}\\
\multicolumn{4}{|p{16cm}|}{}\\
\hline
\multicolumn{4}{|p{16cm}|}{Conteúdo programático:}\\*
\multicolumn{4}{|p{\dimexpr 16cm + 6\tabcolsep\relax}|}{\begin{enumerate}[start=1]\item Álgebra Vetorial: Vetor (definição e notação); Operações fundamentais com vetores (adição e multiplicação por número real); Combinação linear, dependência e independência linear de vetores; Bases Ortogonais e Ortonormais; Multiplicação escalar de dois vetores; Multiplicação vetorial de dois vetores\end{enumerate}}\\
\multicolumn{4}{|p{\dimexpr 16cm + 6\tabcolsep\relax}|}{\begin{enumerate}[start=2]\item A Reta e o Plano no Espaço: Equações da reta; Equação do plano; Interseção de dois planos; Distâncias (de um ponto a um plano, de um ponto a uma reta, entre duas retas); Ângulos\end{enumerate}}\\
\multicolumn{4}{|p{\dimexpr 16cm + 6\tabcolsep\relax}|}{\begin{enumerate}[start=3]\item Matrizes e Sistemas Lineares: Definição; Operações com matrizes e suas propriedades; Resolução de Sistemas de Equações Lineares (Matrizes Escalonadas, Eliminação de Gauss-Jordan); Inversas de Matrizes\end{enumerate}}\\
\multicolumn{4}{|p{\dimexpr 16cm + 6\tabcolsep\relax}|}{\begin{enumerate}[start=4]\item Determinantes: Definição por Cofatores; Propriedades; Regra de Cramer\end{enumerate}}\\
\multicolumn{4}{|p{\dimexpr 16cm + 6\tabcolsep\relax}|}{\begin{enumerate}[start=5]\item Espaços Vetoriais: Definição; Subespaços Vetoriais; Dependência e Independência linear; Bases e dimensão; Espaço-linha, espaço-coluna e posto de uma matriz; Produto Interno em um espaço vetorial (desigualdade de Cauchy-Schwarz)\end{enumerate}}\\
\multicolumn{4}{|p{\dimexpr 16cm + 6\tabcolsep\relax}|}{\begin{enumerate}[start=6]\item Autovalores, Autovetores e Diagonalização: Definição; Polinômio Característico; Diagonalização; Diagonalização de Matrizes Simétricas\end{enumerate}}\\
\multicolumn{4}{|p{\dimexpr 16cm + 6\tabcolsep\relax}|}{}\\
\hline
\multicolumn{4}{|p{16cm}|}{Bibliografia Básica:}\\*
\multicolumn{4}{|p{\dimexpr 16cm + 6\tabcolsep\relax}|}{\begin{itemize}\item Anton, Howard. Álgebra Linear. Rio de Janeiro: Campus.\end{itemize}}\\
\multicolumn{4}{|p{\dimexpr 16cm + 6\tabcolsep\relax}|}{\begin{itemize}\item Lehmann, Charles H. Geometria analítica. 8. ed. São Paulo: Globo, 1995.\end{itemize}}\\
\multicolumn{4}{|p{\dimexpr 16cm + 6\tabcolsep\relax}|}{\begin{itemize}\item Santos, Reginaldo J. Um curso de geometria analítica e álgebra linear. Belo Horizonte (MG): Imprensa Universitária da UFMG, 2003.\end{itemize}}\\
\multicolumn{4}{|p{\dimexpr 16cm + 6\tabcolsep\relax}|}{\begin{itemize}\item Steinbruch, Alfredo. Geometria analítica. 2. ed. São Paulo: Makron Books, c1987.\end{itemize}}\\
\multicolumn{4}{|p{\dimexpr 16cm + 6\tabcolsep\relax}|}{\begin{itemize}\item Steinbruch, Alfredo. Álgebra linear. 2. ed. São Paulo: McGraw-Hill, 1987.\end{itemize}}\\
\multicolumn{4}{|p{\dimexpr 16cm + 6\tabcolsep\relax}|}{\begin{itemize}\item Kolman, Bernard. Álgebra linear. 3. ed. Livros Técnicos e Científicos.\end{itemize}}\\
\multicolumn{4}{|p{16cm}|}{}\\
\hline
\multicolumn{4}{|p{16cm}|}{Bibliografia Complementar:}\\*
\multicolumn{4}{|p{\dimexpr 16cm + 6\tabcolsep\relax}|}{\begin{itemize}\item[] (nenhum item informado)\end{itemize}}\\
\hline
\end{longtable}
\end{center}

\clearpage
\begin{center}
\begin{tabular}{lccr}
 \multirow{3}{*}{\includegraphics[height=2.7cm]{brasao.png}} &
 \multicolumn{2}{c}{\bfseries UNIVERSIDADE FEDERAL DE OURO PRETO} &
 \ \ \ \ \multirow{3}{*}{\includegraphics[height=2.7cm]{ufop.png}} \\
 & \multicolumn{2}{c}{\bfseries PRÓ-REITORIA DE GRADUAÇÃO} & \\
 & \multicolumn{2}{c}{} & \\
 & \multicolumn{2}{c}{\Large\bfseries PROGRAMA DE DISCIPLINA} & \\
\end{tabular}
\end{center}

\begin{center}
\begin{longtable}{|p{4cm}|p{4cm}|p{4cm}|p{4cm}|}
\hline
\multicolumn{3}{|p{12cm}|}{Nome do Componente Curricular em Português:} &
\multicolumn{1}{p{4cm}|}{Código:} \\ 
\multicolumn{3}{|p{12cm}|}{\textbf{Planejamento e Análise de Experimentos}} &
\textbf{EST326}\\ 
\multicolumn{3}{|p{12cm}|}{Nome do Componente Curricular em Inglês:} & \\ 
\multicolumn{3}{|p{12cm}|}{\textbf{Experimental Design and Analysis}} & \\ 
\hline
\multicolumn{3}{|p{12cm}|}{Nome e Sigla do Departamento} & Unidade Acadêmica: \\ 
\multicolumn{3}{|p{12cm}|}{Departamento de Estatística (DEEST)} & {ICEB} \\ 
\hline
\multicolumn{4}{|p{16cm}|}{Modalidade de Oferta:
[ ] presencial \hspace{0.6cm}
[X] semipresencial \hspace{0.6cm}
[ ] à distância}\\
\hline
\multicolumn{2}{|p{8cm}|}{Carga horária semestral} &
\multicolumn{2}{p{8cm}|}{Carga horária semanal}\\
\hline
\multicolumn{1}{|p{4cm}|}{Total} &
\multicolumn{1}{p{4cm}|}{Extensionista} &
\multicolumn{1}{p{4cm}|}{Teórica} &
\multicolumn{1}{p{4cm}|}{Prática} \\ 
\multicolumn{1}{|p{4cm}|}{60\,horas} &
\multicolumn{1}{p{4cm}|}{0\;horas} &
\multicolumn{1}{p{4cm}|}{4\;horas/aula} &
\multicolumn{1}{p{4cm}|}{0\;horas/aula} \\ 
\hline
\multicolumn{4}{|p{16cm}|}{Ementa:}\\
\multicolumn{4}{|p{16cm}|}{}\\
\multicolumn{4}{|p{\dimexpr 16cm + 6\tabcolsep\relax}|}{Princípios básicos de experimentação. Experimentos com um fator: efeitos fixos e aleatórios. Planejamentos em blocos aleatorizados, quadrados latinos. Verificação das pressuposições do modelo. Regressão na análise de variância. Planejamentos fatoriais: dois fatores e três fatores. Experimentos com restrições na aleatorização. Análise de Covariância.}\\
\multicolumn{4}{|p{16cm}|}{}\\
\hline
\multicolumn{4}{|p{16cm}|}{Conteúdo programático:}\\*
\multicolumn{4}{|p{\dimexpr 16cm + 6\tabcolsep\relax}|}{\begin{enumerate}[start=1]\item Unidade 1 - Princípios básicos de experimentação: conceitos importantes em experimentação, repetição/réplica, casualização/randomização e controle local; delineamentos experimentais\end{enumerate}}\\
\multicolumn{4}{|p{\dimexpr 16cm + 6\tabcolsep\relax}|}{\begin{enumerate}[start=2]\item Unidade 2 - Experimentos com um fator: modelos fixos, aleatórios e mistos; análise de variância e teste F; contrastes, comparações múltiplas e método de agrupamento de Scott-Knott\end{enumerate}}\\
\multicolumn{4}{|p{\dimexpr 16cm + 6\tabcolsep\relax}|}{\begin{enumerate}[start=3]\item Unidade 3 - Delineamento em blocos casualizados e em quadrado latino\end{enumerate}}\\
\multicolumn{4}{|p{\dimexpr 16cm + 6\tabcolsep\relax}|}{\begin{enumerate}[start=4]\item Unidade 4 - Análise de resíduos em experimentos e verificação das pressuposições do modelo linear geral de Gauss-Markov\end{enumerate}}\\
\multicolumn{4}{|p{\dimexpr 16cm + 6\tabcolsep\relax}|}{\begin{enumerate}[start=5]\item Unidade 5 - Dimensionamento de amostra em experimentos; experimentos com subamostragem; uso de regressão em experimentos quando os tratamentos são de natureza quantitativa\end{enumerate}}\\
\multicolumn{4}{|p{\dimexpr 16cm + 6\tabcolsep\relax}|}{\begin{enumerate}[start=6]\item Unidade 6 - Ensaios fatoriais com dois e três fatores; introdução à técnica de superfícies de resposta\end{enumerate}}\\
\multicolumn{4}{|p{\dimexpr 16cm + 6\tabcolsep\relax}|}{\begin{enumerate}[start=7]\item Unidade 7 - Experimentos com restrições na aleatorização: ensaios em parcelas subdivididas e em faixas\end{enumerate}}\\
\multicolumn{4}{|p{\dimexpr 16cm + 6\tabcolsep\relax}|}{\begin{enumerate}[start=8]\item Unidade 8 - Análise de Covariância: aplicações e análise\end{enumerate}}\\
\multicolumn{4}{|p{\dimexpr 16cm + 6\tabcolsep\relax}|}{\begin{enumerate}[start=9]\item Conteúdos trabalhados de forma semipresencial: 53 horas presenciais e 7 horas a distância\end{enumerate}}\\
\multicolumn{4}{|p{\dimexpr 16cm + 6\tabcolsep\relax}|}{}\\
\hline
\multicolumn{4}{|p{16cm}|}{Bibliografia Básica:}\\*
\multicolumn{4}{|p{\dimexpr 16cm + 6\tabcolsep\relax}|}{\begin{itemize}\item Montgomery, D. C. Design and Analysis of Experiments. 8. ed. New York: J. Wiley, 2013.\end{itemize}}\\
\multicolumn{4}{|p{\dimexpr 16cm + 6\tabcolsep\relax}|}{\begin{itemize}\item Montgomery, D. C.; Kowalski, S. M. Design and Analysis of Experiments: MINITAB Companion. 7. ed. Wiley, 2010.\end{itemize}}\\
\multicolumn{4}{|p{\dimexpr 16cm + 6\tabcolsep\relax}|}{\begin{itemize}\item Pimentel-Gomes, F. Curso de Estatística Experimental. 15. ed. Piracicaba: FEALQ, 2009.\end{itemize}}\\
\multicolumn{4}{|p{\dimexpr 16cm + 6\tabcolsep\relax}|}{\begin{itemize}\item Hinkelmann, Klaus; Kempthorne, Oscar. Design and Analysis of Experiments: introduction to experimental design. v.1. Hoboken: John Wiley \& Sons, 2008. 631 p.\end{itemize}}\\
\multicolumn{4}{|p{16cm}|}{}\\
\hline
\multicolumn{4}{|p{16cm}|}{Bibliografia Complementar:}\\*
\multicolumn{4}{|p{\dimexpr 16cm + 6\tabcolsep\relax}|}{\begin{itemize}\item Hinkelmann, Klaus; Kempthorne, Oscar. Design and analysis of experiments: advanced experimental design. v. 2. New York: Wiley-Interscience, 2005. 780 p. (Wiley Series in Probability and Statistics).\end{itemize}}\\
\multicolumn{4}{|p{\dimexpr 16cm + 6\tabcolsep\relax}|}{\begin{itemize}\item Calegare, A. J. A. Introdução ao Delineamento de Experimentos. 2. ed. São Paulo: Edgard Blucher, 2009.\end{itemize}}\\
\multicolumn{4}{|p{\dimexpr 16cm + 6\tabcolsep\relax}|}{\begin{itemize}\item Dean, A. M.; Voss, D. Design and Analysis of Experiments. Springer, 2000.\end{itemize}}\\
\multicolumn{4}{|p{\dimexpr 16cm + 6\tabcolsep\relax}|}{\begin{itemize}\item Werkema, M. C. C.; Aguiar, S. Planejamento e Análise de Experimentos: Como Identificar as Principais Variáveis Influentes em um Processo. Belo Horizonte, MG: Fundação Christiano Ottoni, E.E. UFMG, 1996.\end{itemize}}\\
\multicolumn{4}{|p{\dimexpr 16cm + 6\tabcolsep\relax}|}{\begin{itemize}\item Vieira, S. Análise de Variância: Anova. Atlas, 2006.\end{itemize}}\\
\multicolumn{4}{|p{\dimexpr 16cm + 6\tabcolsep\relax}|}{\begin{itemize}\item Cochran, W. G.; Cox, G. M. Experimental Designs. 2. ed. Wiley, 1992.\end{itemize}}\\
\hline
\end{longtable}
\end{center}

\clearpage
\begin{center}
\begin{tabular}{lccr}
 \multirow{3}{*}{\includegraphics[height=2.7cm]{brasao.png}} &
 \multicolumn{2}{c}{\bfseries UNIVERSIDADE FEDERAL DE OURO PRETO} &
 \ \ \ \ \multirow{3}{*}{\includegraphics[height=2.7cm]{ufop.png}} \\
 & \multicolumn{2}{c}{\bfseries PRÓ-REITORIA DE GRADUAÇÃO} & \\
 & \multicolumn{2}{c}{} & \\
 & \multicolumn{2}{c}{\Large\bfseries PROGRAMA DE DISCIPLINA} & \\
\end{tabular}
\end{center}

\begin{center}
\begin{longtable}{|p{4cm}|p{4cm}|p{4cm}|p{4cm}|}
\hline
\multicolumn{3}{|p{12cm}|}{Nome do Componente Curricular em Português:} &
\multicolumn{1}{p{4cm}|}{Código:} \\ 
\multicolumn{3}{|p{12cm}|}{\textbf{Análise de Séries Temporais e Aprendizado Dinâmico}} &
\textbf{EST334}\\ 
\multicolumn{3}{|p{12cm}|}{Nome do Componente Curricular em Inglês:} & \\ 
\multicolumn{3}{|p{12cm}|}{\textbf{Time Series Analysis and Dynamic Learning}} & \\ 
\hline
\multicolumn{3}{|p{12cm}|}{Nome e Sigla do Departamento} & Unidade Acadêmica: \\ 
\multicolumn{3}{|p{12cm}|}{Departamento de Estatística (DEEST)} & {ICEB} \\ 
\hline
\multicolumn{4}{|p{16cm}|}{Modalidade de Oferta:
[ ] presencial \hspace{0.6cm}
[X] semipresencial \hspace{0.6cm}
[ ] à distância}\\
\hline
\multicolumn{2}{|p{8cm}|}{Carga horária semestral} &
\multicolumn{2}{p{8cm}|}{Carga horária semanal}\\
\hline
\multicolumn{1}{|p{4cm}|}{Total} &
\multicolumn{1}{p{4cm}|}{Extensionista} &
\multicolumn{1}{p{4cm}|}{Teórica} &
\multicolumn{1}{p{4cm}|}{Prática} \\ 
\multicolumn{1}{|p{4cm}|}{60\,horas} &
\multicolumn{1}{p{4cm}|}{0\;horas} &
\multicolumn{1}{p{4cm}|}{4\;horas/aula} &
\multicolumn{1}{p{4cm}|}{0\;horas/aula} \\ 
\hline
\multicolumn{4}{|p{16cm}|}{Ementa:}\\
\multicolumn{4}{|p{16cm}|}{}\\
\multicolumn{4}{|p{\dimexpr 16cm + 6\tabcolsep\relax}|}{Introdução. Conceitos importantes. Modelos para Séries Estacionárias. Os modelos ARIMA. Os modelos SARIMA. Identificação de Modelos. Estimação e Verificação. Previsão com modelos ARIMA. Métodos Automáticos de Previsão. Tópicos Especiais.}\\
\multicolumn{4}{|p{16cm}|}{}\\
\hline
\multicolumn{4}{|p{16cm}|}{Conteúdo programático:}\\*
\multicolumn{4}{|p{\dimexpr 16cm + 6\tabcolsep\relax}|}{\begin{enumerate}[start=1]\item Unidade 1 - Introdução: definições, exemplos e notação; classificação de séries temporais; objetivos e enfoques da Análise de Séries Temporais; tipos de variações\end{enumerate}}\\
\multicolumn{4}{|p{\dimexpr 16cm + 6\tabcolsep\relax}|}{\begin{enumerate}[start=2]\item Unidade 2 - Conceitos importantes: séries temporais e processos estocásticos; média, covariância e correlação; processos estacionários; processo ruído branco\end{enumerate}}\\
\multicolumn{4}{|p{\dimexpr 16cm + 6\tabcolsep\relax}|}{\begin{enumerate}[start=3]\item Unidade 3 - Modelos para Séries Estacionárias: processo linear geral; processos auto-regressivos; processos médias móveis; modelos mistos ARMA(p,q); condições de estacionariedade e invertibilidade\end{enumerate}}\\
\multicolumn{4}{|p{\dimexpr 16cm + 6\tabcolsep\relax}|}{\begin{enumerate}[start=4]\item Unidade 4 - Os modelos ARIMA: séries não-estacionárias\end{enumerate}}\\
\multicolumn{4}{|p{\dimexpr 16cm + 6\tabcolsep\relax}|}{\begin{enumerate}[start=5]\item Unidade 5 - Os modelos SARIMA: séries não-estacionárias com sazonalidade\end{enumerate}}\\
\multicolumn{4}{|p{\dimexpr 16cm + 6\tabcolsep\relax}|}{\begin{enumerate}[start=6]\item Unidade 6 - Identificação de Modelos: propriedades da função de autocorrelação amostral; função de autocorrelação parcial; função de autocorrelação inversa; outros procedimentos de identificação\end{enumerate}}\\
\multicolumn{4}{|p{\dimexpr 16cm + 6\tabcolsep\relax}|}{\begin{enumerate}[start=7]\item Unidade 7 - Estimação e Verificação: método de Momentos; método dos Mínimos Quadrados; método da Máxima Verossimilhança; propriedades dos estimadores; análise de resíduos\end{enumerate}}\\
\multicolumn{4}{|p{\dimexpr 16cm + 6\tabcolsep\relax}|}{\begin{enumerate}[start=8]\item Unidade 8 - Previsão com modelos ARIMA: cálculo das previsões; formas básicas de previsão; intervalos de confiança\end{enumerate}}\\
\multicolumn{4}{|p{\dimexpr 16cm + 6\tabcolsep\relax}|}{\begin{enumerate}[start=9]\item Unidade 9 - Métodos Automáticos de Previsão: modelos de alisamento exponencial (método de Holt-Winters)\end{enumerate}}\\
\multicolumn{4}{|p{\dimexpr 16cm + 6\tabcolsep\relax}|}{\begin{enumerate}[start=10]\item Unidade 10 - Tópicos Especiais: introdução a modelos de volatilidade (ARCH e GARCH); modelos lineares dinâmicos; modelos de regressão dinâmica; modelos de função de transferência\end{enumerate}}\\
\multicolumn{4}{|p{\dimexpr 16cm + 6\tabcolsep\relax}|}{\begin{enumerate}[start=11]\item Conteúdos trabalhados de forma semipresencial: 53 horas presenciais e 7 horas a distância\end{enumerate}}\\
\multicolumn{4}{|p{\dimexpr 16cm + 6\tabcolsep\relax}|}{}\\
\hline
\multicolumn{4}{|p{16cm}|}{Bibliografia Básica:}\\*
\multicolumn{4}{|p{\dimexpr 16cm + 6\tabcolsep\relax}|}{\begin{itemize}\item Enders, Walter. Applied econometric time series. 2. ed. New York: John Wiley, c2004.\end{itemize}}\\
\multicolumn{4}{|p{\dimexpr 16cm + 6\tabcolsep\relax}|}{\begin{itemize}\item Morettin, Pedro Alberto. Análise de séries temporais. 2. ed. São Paulo: E. Blücher, 2008.\end{itemize}}\\
\multicolumn{4}{|p{\dimexpr 16cm + 6\tabcolsep\relax}|}{\begin{itemize}\item Brockwell, P. J. Introduction to time series and forecasting. 2. ed. New York: Springer-Verlag, c2002.\end{itemize}}\\
\multicolumn{4}{|p{\dimexpr 16cm + 6\tabcolsep\relax}|}{\begin{itemize}\item Chatfield, C. The Analysis of Time Series: An Introduction. 6. ed. Chapman and Hall/CRC, 2003.\end{itemize}}\\
\multicolumn{4}{|p{\dimexpr 16cm + 6\tabcolsep\relax}|}{\begin{itemize}\item Box, G. E. P.; Jenkins, G. M. Time Series Analysis: Forecasting and Control. 3. ed. Prentice Hall, 1994.\end{itemize}}\\
\multicolumn{4}{|p{16cm}|}{}\\
\hline
\multicolumn{4}{|p{16cm}|}{Bibliografia Complementar:}\\*
\multicolumn{4}{|p{\dimexpr 16cm + 6\tabcolsep\relax}|}{\begin{itemize}\item Wei, William W. S. Time series analysis: univariate and multivariate methods. 2. ed. Boston: Pearson Education Addison Wesley, c2006. 614 p. ISBN 0321322169.\end{itemize}}\\
\multicolumn{4}{|p{\dimexpr 16cm + 6\tabcolsep\relax}|}{\begin{itemize}\item Pole, Andy. Applied bayesian forecasting and time series analysis. Boca Raton, FL: Chapman \& Hall/CRC, 1994.\end{itemize}}\\
\multicolumn{4}{|p{\dimexpr 16cm + 6\tabcolsep\relax}|}{\begin{itemize}\item Hamilton, J. D. Time Series Analysis. Princeton University Press, 1994.\end{itemize}}\\
\multicolumn{4}{|p{\dimexpr 16cm + 6\tabcolsep\relax}|}{\begin{itemize}\item Montgomery, D. C.; Johnson, L. A. Forecasting and Time Series Analysis. New York: McGraw-Hill, 1976.\end{itemize}}\\
\multicolumn{4}{|p{\dimexpr 16cm + 6\tabcolsep\relax}|}{\begin{itemize}\item Pankratz, A. Forecasting with univariate Box and Jenkins Models. New York: John Wiley \& Sons, 1983.\end{itemize}}\\
\multicolumn{4}{|p{\dimexpr 16cm + 6\tabcolsep\relax}|}{\begin{itemize}\item Nelson, C. R. Applied Time Series Analysis. San Francisco: Holden Day, 1973.\end{itemize}}\\
\hline
\end{longtable}
\end{center}

\clearpage
\end{document}